%!TEX root = LL_CV.tex
%%
%% Shared preamble for every document in this repository.
%%
%%   LL_CV.tex            the full CV, in an external and an internal edition
%%   LL_Publications.tex  the publication list on its own
%%
%% Both are just \documentclass, %!TEX root = LL_CV.tex
%%
%% Shared preamble for every document in this repository.
%%
%%   LL_CV.tex            the full CV, in an external and an internal edition
%%   LL_Publications.tex  the publication list on its own
%%
%% Both are just \documentclass, %!TEX root = LL_CV.tex
%%
%% Shared preamble for every document in this repository.
%%
%%   LL_CV.tex            the full CV, in an external and an internal edition
%%   LL_Publications.tex  the publication list on its own
%%
%% Both are just \documentclass, %!TEX root = LL_CV.tex
%%
%% Shared preamble for every document in this repository.
%%
%%   LL_CV.tex            the full CV, in an external and an internal edition
%%   LL_Publications.tex  the publication list on its own
%%
%% Both are just \documentclass, \input{LL_Preamble.tex}, and a body, so page
%% geometry, fonts, list styles and the custom commands below are only ever
%% edited in one place.  A document that wants to differ from a default
%% overrides it after the \input.
%%
%% Template attribution: Jason R. Blevins - Curriculum Vitae
%% Copyright (C) 2004-2013 Jason R. Blevins


%% ===========================================================================
%% Packages
%% ===========================================================================
%% hyperref goes last: it patches internals of most other packages.

\usepackage[T1]{fontenc}
\usepackage{CormorantGaramond}    % body and heading face
\usepackage{microtype}            % subtler justification, fewer overfull lines
\usepackage{wasysym}              % $\RHD$, the "lead role" marker
\usepackage{graphicx}             % for the (currently disabled) portrait

\usepackage{geometry}             % page size and margins
\usepackage{fancyhdr}             % running headers
\usepackage{lastpage}             % "page n of m"
\usepackage{titlesec}             % section spacing
\usepackage{sectsty}              % section fonts
\usepackage{enumitem}             % list configuration

\usepackage[hidelinks]{hyperref}  % clickable links, printed without decoration


%% ===========================================================================
%% Identity
%% ===========================================================================

\def\name{Lawrence Lee, PhD}
\def\email{Lawrence.Lee.Jr@cern.ch}
\def\phone{+1 856 765 4622}
\def\homepage{cern.ch/larry}
\def\orcid{0000-0002-5590-335X}


%% ===========================================================================
%% Internal and external editions
%% ===========================================================================
%% The internal edition carries material that is not for general circulation:
%% the current student and postdoc roster, thesis committee membership,
%% individual refereeing assignments, and student travel awards.  The external
%% edition is the default and needs no flag; the internal one is selected by
%% defining \CVinternal before the document is read:
%%
%%     pdflatex -jobname=LL_CV_internal "\def\CVinternal{}\input{LL_CV.tex}"
%%
%% which is what build.sh and .github/workflows/build-cv.yml do.

\newif\ifCVinternal
\ifdefined\CVinternal \CVinternaltrue \else \CVinternalfalse \fi

% \internalonly{<content>} -- included only in the internal edition.
\newcommand{\internalonly}[1]{\ifCVinternal #1\fi}

% \externalonly{<content>} -- included only in the external edition.
\newcommand{\externalonly}[1]{\ifCVinternal\else #1\fi}

% The internal edition says so in its PDF title, so the two files stay
% distinguishable once they have been detached from their filenames.
\ifCVinternal
  \def\LLedition{\name: Curriculum Vitae (internal edition)}
\else
  \def\LLedition{\name: Curriculum Vitae}
\fi

\hypersetup{
  colorlinks  = false,
  pdfauthor   = {\name},
  pdftitle    = {\LLedition},
  pdfsubject  = {Curriculum Vitae},
  pdfkeywords = {},
  pdfpagemode = UseNone,
}


%% ===========================================================================
%% Page layout
%% ===========================================================================

\geometry{
  body       = {6.5in, 9.0in},
  left       = 1.0in,
  top        = 1.0in,
  headheight = 14pt,   % enough room for the running header; silences fancyhdr
}

\pagestyle{fancy}
\fancyhf{}
\renewcommand{\headrulewidth}{0pt}
\fancyhead[L]{\name}
\fancyhead[R]{\thepage}

\setlength{\parindent}{0em}

\titlespacing{\section}{0pt}{0pt}{0pt}
\sectionfont{\rmfamily\mdseries\Large}
\subsectionfont{\rmfamily\mdseries\itshape\large}


%% ===========================================================================
%% Lists
%% ===========================================================================
%% The CV is built almost entirely from bullet-free lists that differ only in
%% how tightly they are set.  itemize itself is redefined, rather than
%% configured through enumitem, so that a plain \begin{itemize} anywhere picks
%% up the CV style.
%%
%%   itemize               default spacing between entries
%%   itemizetight          no extra space between entries
%%   itemizetighter        entries set as close as the leading allows
%%   itemizetightrightpad  tight, and inset from the right margin, for the
%%                         commentary set underneath an entry
%%
%% \LLlist{<left>}{<right>}{<itemsep>}{<parskip>}{<parsep>}

\newcommand{\LLlist}[5]{%
  \list{}{%
    \setlength{\leftmargin}{#1}%
    \setlength{\rightmargin}{#2}%
    \setlength{\itemsep}{#3}%
    \setlength{\parskip}{#4}%
    \setlength{\parsep}{#5}%
  }%
}

\renewenvironment{itemize}%
  {\LLlist{1.5em}{0em}{0.25em}{0pt}{0.25em}}{\endlist}
\newenvironment{itemizetight}%
  {\LLlist{1.5em}{0em}{0em}{-1pt}{0em}}{\endlist}
\newenvironment{itemizetighter}%
  {\LLlist{1.5em}{0em}{0em}{-0.4em}{0em}}{\endlist}
\newenvironment{itemizetightrightpad}%
  {\LLlist{1.5em}{3em}{0em}{-1pt}{0em}}{\endlist}

% Hanging indent for a run-on entry: \parshape 2 3.2em \linewidth \listindent
% \dimexpr\linewidth-\listindent\relax sets the first line full width and
% indents every line after it to \listindent.
\newlength{\listindent}
\setlength{\listindent}{4.8em}

% \changemargin{<left>}{<right>} ... \endchangemargin -- indents a block of the
% document; used to nest the collaboration sub-lists under their headings.
\def\changemargin#1#2{\list{}{\rightmargin#2\leftmargin#1}\item[]}
\let\endchangemargin=\endlist


%% ===========================================================================
%% The right-hand date column
%% ===========================================================================
%% Nearly every entry ends with a date set against the right margin.  \when
%% supplies both the \hfill and the fixed-width box that lines those dates up
%% into one column: the box is left-aligned, so "2024" and "2022--24" begin at
%% the same horizontal position.
%%
%% A \makebox is only as wide as it is declared, so an over-long date used to
%% overflow it silently and print out in the right margin -- six of the dates in
%% this CV ("2021--2026", "2022--23, 24--25", ...) are wider than the column.
%% \when measures the date first and sets an over-wide one flush right instead,
%% which costs the shared left edge on those few entries but keeps every date
%% inside the text block.

\newlength{\datecolumn}
\setlength{\datecolumn}{3.5em}

%% The glue ahead of a date is \datesep plus infinite stretch.  The stretch is
%% what pushes the date to the right margin; the fixed \datesep is a floor, so
%% an entry that grows long enough to reach its date reports an overfull box
%% instead of quietly printing its text hard against the year.  Every entry
%% currently clears it, so the floor costs nothing today -- it is there to catch
%% the next entry that does not.  \whenflush is the fix when one turns up: it
%% gives back the empty right-hand part of the date box, about 1.6em for a
%% four-digit year.
\newlength{\datesep}
\setlength{\datesep}{0.4em}

%% \LLdateglue is the glue that carries the date to the right margin.  It reads
%% as one \hfill most of the time, but is written out so that an entry too long
%% to hold its date can break: TeX takes the \penalty50, leaves the text on the
%% line it filled, and the \hbox{} -- which survives the break where glue would
%% not -- carries the second \hfill onto the next line, so the date still lands
%% flush right instead of stranded against the left margin.  \unskip drops the
%% source space ahead of the date, so the gap is \datesep however the entry was
%% typed.
\newcommand{\LLdateglue}{%
  \unskip\nobreak\hspace{\datesep}\hfill\penalty50 \hbox{}\nobreak\hfill
}

\newlength{\LLdatewidth}
\newcommand{\when}[1]{%
  \LLdateglue
  \settowidth{\LLdatewidth}{#1}%
  \ifdim\LLdatewidth>\datecolumn #1\else\makebox[\datecolumn][l]{#1}\fi
}

% \whenflush{<date>} -- for an entry whose text runs so close to the right
% margin that the shared left edge of the date column will not fit.  Sets the
% date flush right, giving up the alignment to keep the gap.
\newcommand{\whenflush}[1]{\LLdateglue #1}

% \rightnote{<text>} -- a right-aligned italic annotation (an award, a current
% position) hung underneath an entry.  The trailing pad makes it end where a
% four-digit year in the date column ends, rather than at the margin itself.
\newcommand{\rightnote}[1]{%
  \settowidth{\LLdatewidth}{2024}%
  \hfill\emph{#1}\hspace{\dimexpr\datecolumn-\LLdatewidth\relax}%
}


%% ===========================================================================
%% Live publication metrics from INSPIRE-HEP
%% ===========================================================================
%% Citation counts go stale the moment they are typed, so they are not written
%% into the CV by hand.  tools/fetch_inspire.py queries the INSPIRE-HEP API and
%% writes LL_InspireData.tex, a flat list of declarations:
%%
%%     \inspiresetcites{2690093}{3}         % citations of one paper
%%     \inspiresetstat{citations}{207,000}  % an author-level summary figure
%%
%% which the CV body reads back through \inspirepub and \inspirestat.  The
%% values set just below are fallbacks for a build that has never run the
%% fetcher -- an offline checkout, or Overleaf -- so the document always
%% compiles with plausible numbers.  build.sh and CI refresh the data first.

% Declarations written by tools/fetch_inspire.py.  The keys live in \csname, so
% the LL@ namespace prefix needs no \makeatletter here.
\newcommand{\inspiresetcites}[2]{\expandafter\gdef\csname LL@cites@#1\endcsname{#2}}
\newcommand{\inspiresetstat}[2]{\expandafter\gdef\csname LL@stat@#1\endcsname{#2}}

% Fallbacks: the last hand-checked figures, overridden by LL_InspireData.tex.
\inspiresetstat{papers}{1400}
\inspiresetstat{citations}{202,000}
\inspiresetstat{hindex}{209}

% \inspirestat{<key>} -- an author-level figure: papers, citations or hindex.
% A missing key is reported loudly rather than left silently blank.
\newcommand{\inspirestat}[1]{%
  \ifcsname LL@stat@#1\endcsname\csname LL@stat@#1\endcsname\else\textbf{[?~#1]}\fi
}

% \inspirecites{<record>} -- "[n citations]" for a record whose count is known.
% Nothing is printed for an unknown count, or for one below \citesmin, so a
% brand-new paper is left unannotated rather than advertising that nobody has
% cited it yet.
%
% The annotation costs about a dozen characters on every entry, which is enough
% to push an entry that already filled its line onto a second one.  Raising
% \citesmin, or shortening the label, is the dial if that happens too often.
\newcommand{\citesmin}{1}
\newcommand{\LLprintcites}[1]{%
  \ifnum#1<\citesmin\relax\else
    \nobreakspace{\small[#1\nobreakspace\ifnum#1=1\relax citation\else citations\fi]}%
  \fi
}
\newcommand{\inspirecites}[1]{%
  \ifcsname LL@cites@#1\endcsname\LLprintcites{\csname LL@cites@#1\endcsname}\fi
}

% \inspirepub{<record>}{<title>} -- a publication title in bold, linked to its
% INSPIRE-HEP record and followed by that record's live citation count.
\newcommand{\inspirepub}[2]{%
  \href{https://inspirehep.net/literature/#1}{\textbf{#2}}\inspirecites{#1}%
}

% \pub{<url>}{<title>} -- the same, for an entry with no INSPIRE-HEP record:
% arXiv-only preprints, CDS and CERN notes, bare DOI links.
\newcommand{\pub}[2]{\href{#1}{\textbf{#2}}}

% The generated data, if it has been fetched.  \IfFileExists keeps a fresh
% checkout compiling on the fallbacks above.
\IfFileExists{LL_InspireData.tex}{\input{LL_InspireData.tex}}{%
  \typeout{^^JLL_CV: LL_InspireData.tex not found; using fallback INSPIRE
figures. Run tools/fetch_inspire.py to refresh them.^^J}%
}


%% ===========================================================================
%% Title block
%% ===========================================================================

% \cvheader -- the name, contact details and ORCID that open every document.
\newcommand{\cvheader}{%
  \begin{minipage}[t]{0.4\textwidth}
    {\huge \name}
  \end{minipage}%
  \begin{minipage}[t]{0.6\textwidth}
    \begin{flushright}
      \href{mailto:\email}{\email}\\
      \phone\\
      \href{https://\homepage}{\homepage}\\
      \href{https://orcid.org/\orcid}{\orcid}
      % A postal address, if one is ever wanted here:
      %   CERN CH-1211\\
      %   B\^atiment 40-1-D11\\
      %   23 Gen\`{e}ve, Switzerland
    \end{flushright}
  \end{minipage}%
}

% \cvfooter -- the "last updated" line that closes every document.
\newcommand{\cvfooter}{%
  \medskip
  \begin{center}
    \begin{small}
      Last updated: \today
    \end{small}
  \end{center}%
}
, and a body, so page
%% geometry, fonts, list styles and the custom commands below are only ever
%% edited in one place.  A document that wants to differ from a default
%% overrides it after the \input.
%%
%% Template attribution: Jason R. Blevins - Curriculum Vitae
%% Copyright (C) 2004-2013 Jason R. Blevins


%% ===========================================================================
%% Packages
%% ===========================================================================
%% hyperref goes last: it patches internals of most other packages.

\usepackage[T1]{fontenc}
\usepackage{CormorantGaramond}    % body and heading face
\usepackage{microtype}            % subtler justification, fewer overfull lines
\usepackage{wasysym}              % $\RHD$, the "lead role" marker
\usepackage{graphicx}             % for the (currently disabled) portrait

\usepackage{geometry}             % page size and margins
\usepackage{fancyhdr}             % running headers
\usepackage{lastpage}             % "page n of m"
\usepackage{titlesec}             % section spacing
\usepackage{sectsty}              % section fonts
\usepackage{enumitem}             % list configuration

\usepackage[hidelinks]{hyperref}  % clickable links, printed without decoration


%% ===========================================================================
%% Identity
%% ===========================================================================

\def\name{Lawrence Lee, PhD}
\def\email{Lawrence.Lee.Jr@cern.ch}
\def\phone{+1 856 765 4622}
\def\homepage{cern.ch/larry}
\def\orcid{0000-0002-5590-335X}


%% ===========================================================================
%% Internal and external editions
%% ===========================================================================
%% The internal edition carries material that is not for general circulation:
%% the current student and postdoc roster, thesis committee membership,
%% individual refereeing assignments, and student travel awards.  The external
%% edition is the default and needs no flag; the internal one is selected by
%% defining \CVinternal before the document is read:
%%
%%     pdflatex -jobname=LL_CV_internal "\def\CVinternal{}%% Lawrence Lee -- Curriculum Vitae
%%
%% Two editions are produced from this one source:
%%
%%   LL_CV.pdf           external, for general circulation
%%   LL_CV_internal.pdf  internal; additionally carries the current student and
%%                       postdoc roster, thesis committee membership, individual
%%                       refereeing assignments, and student travel awards
%%
%% Build both with ./build.sh.  Everything shared with LL_Publications.tex --
%% packages, page layout, list styles, the \internalonly / \externalonly
%% switches, and the live INSPIRE-HEP citation commands -- lives in
%% LL_Preamble.tex.  See README.md for the rest.

\documentclass[11pt,letterpaper]{article}

\input{LL_Preamble.tex}


\begin{document}
\thispagestyle{empty} % no running header on the first page

\cvheader

% Alternative title blocks, kept for reference:
%   \centerline{\huge \bf \name}          % name centred and bold
%   \begin{picture}(0,0)                  % portrait in the top right corner
%     \put(380,-30){\hbox{\includegraphics[height=12em]{Portrait.jpg}}}
%   \end{picture}

\bigskip

% Spacer that holds the vertical gap ahead of the first section.  It pairs with
% the summary paragraph below, which is currently disabled.
\begin{minipage}[t]{1.3em}
  \hspace{1em}
\end{minipage}
% \begin{minipage}[t]{0.75\textwidth}
% \emph{Broadly-trained particle physicist with expertise in searching for physics beyond the standard model at particle colliders. Particular expertise in supersymmetry, long-lived particles, and unconventional searches. Technical skills include expertise in detector data acquisition system, analysis software, and detector simulation.}
% \end{minipage}


\section*{Professional Experience}

  \begin{itemizetight}


    \item Associate Professor, University of Tennessee, Knoxville \when{2026--}
    \item Assistant Professor, University of Tennessee, Knoxville \when{2021--26}

    % \begin{itemizetighter}
    %   \item {University of Tennessee}, Knoxville, TN
    % \end{itemizetighter}

    \item Research Associate, Harvard University - \emph{Supervisor: J. Huth} \when{2016--21}

    % \begin{itemizetighter}
    %   \item {Harvard University}, Cambridge, MA -- \emph{Supervisor: J. Huth}
    % \end{itemizetighter}

    \item Research Associate, University of Adelaide - \emph{Supervisor: P. Jackson} \when{2014--16}

    % \begin{itemizetighter}
    %   \item {The University of Adelaide}, Adelaide, Australia -- \emph{Supervisor: P. Jackson}
    % \end{itemizetighter}

  \end{itemizetight}


\section*{Education}

  \begin{itemizetight}
    \item Ph.D., M.Phil., M.S., Physics, {Yale University} - \emph{Supervisor: T. Golling} \when{2014}
    % \begin{itemizetighter}
    %       \item \emph{Supervisor: T. Golling}
    % \end{itemizetighter}
    % \item M.S., M.Phil. Physics, {Yale University} \when{2012}
    \item B.S. Physics, {Rutgers University} - \emph{Supervisors: R. Ransome, R. Gilman, R. Tumulka} \when{2009}

  %   \begin{itemizetighter}
  %         \item \emph{Supervisors: R. Ransome, R. Gilman, R. Tumulka}
  %   \end{itemizetighter}
  \end{itemizetight}


% \section*{Fields of Research Interest}

%     \begin{itemizetightrightpad}
%     \item High Energy Collider Physics, Experimental searches for physics Beyond the Standard Model, Supersymmetry, Long-Lived Particles, Foundations of Quantum Mechanics
%     % \item Quantum Information, Philosophical Foundations of Quantum Mechanics
%     \end{itemizetightrightpad}


\section*{Awards \textit{\&} Honors}

\emph{External}

  \begin{itemizetighter}
    \item Cottrell Lectureship \when{2026}
    \item Fermilab LHC Physics Center Senior Distinguished Researcher \when{2026}
    \item Breakthrough Prize in Fundamental Physics Laureate  \when{2025}
    \item \hspace{2em}\emph{as part of the ATLAS and CMS Collaborations}
    \item Cottrell Scholar - Research Corporation for Science Advancement \when{2025}
    \item NSF CAREER Award \when{2023}
    \item European Physical Society's High Energy and Particle Physics Prize \when{2013}
    \item \hspace{2em}\emph{as part of the ATLAS Collaboration}
  \end{itemizetighter}


\emph{University of Tennessee, Knoxville}

\begin{itemizetighter}
  \item University of Tennessee Alumni Association -  Alumni Outstanding Teacher Award \when{2025}
  \item UTK Division of Student Success - Faculty Research Mentor Award \when{2025}
  \item UTK Provost - Junior Faculty Fellow \when{2023--25}
  \item UTK Physics - Research Advisor of the Year \when{2023--24}
  \item UTK College of Arts and Sciences - Faculty Academic Outreach Teaching Award \when{2022}
  \item UTK Physics - Teacher of the Year \when{2021--22}
\end{itemizetighter}


\section*{Funding}

\emph{Multi-Year Grants}

  \begin{itemizetighter}
    \item BNL Subcontract on Accelerator Physics - \$190k \when{2025--29}
    \item Cottrell Scholar - Research Corporation for Science Advancement - \$120k \when{2025--28}
    \item DOE Award for UTK HEP (4 PIs) - \$1.425M \when{2024--27}
    \item NSF CAREER Award - \$1M \when{2023--28}
    \item USCMS Upgrade Project Funding - \$73k \when{2022--24}
    \item Single-PI DOE Award - \$250k \when{2022--24}
  \end{itemizetighter}


  \emph{Other Funding}

  \begin{itemizetighter}
    \item Fermilab LHC Physics Center - \$56k\when{2026}
    \item Universities Research Association - \$7k\when{2026}
    \item Fermilab LHC Physics Center (Postdoc Emery Nibigira Led) - \$43k\when{2024}
    \item Universities Research Association (Postdoc Emery Nibigira Led) - \$20k\when{2024}
    \item UT-Battelle (ORNL Management Contractor) - \$74k \when{2024}
    \internalonly{
      \item UTK Graduate Student Senate Travel Award - John Lawless - \$1k\when{2024}
      \item SESAPS Travel Grant - Caroline Riggall - \$325 \when{2024}
      \item SESAPS Travel Grant - Charles Bell - \$500 \when{2023}
      \item SESAPS Travel Grant - Lacey Dishman - \$500 \when{2023}
      \item UTK Faculty Research Assistants Funding - Tanner Wolfe - \$2.3k \when{2023}
      \item FNAL LPC Guest \textit{\&} Visitor Funding - Emery Nibigira - \$5.2k \when{2023}
      \item FNAL LPC Guest \textit{\&} Visitor Funding - John Lawless - \$3.9k \when{2023}
      \item UTK Faculty Research Assistants Funding - Taylor Sussmane - \$2.3k \when{2022}
      \item SESAPS Travel Grant - Taylor Sussmane - \$500 \when{2022}
      \item SESAPS Travel Grant - Alexander Sizemore - \$500 \when{2022}
      \item FNAL LPC Guest \textit{\&} Visitor Funding - Ramon Ogaz - \$2.8k \when{2022}
      \item FNAL LPC Guest \textit{\&} Visitor Funding - Philipp Wagenknecht - \$5.2k \when{2022}
      \item UTK Physics Summer Research Fellowship - Matthew Durkee - \$6k \when{2022}
    }


    \item George Southgate Fellowship, The University of Adelaide - 5k AUD \when{2015}
  \end{itemizetighter}


\section*{Research Experience}

  \begin{itemize}


  \item \textsc{Member of the CMS Collaboration} \when{2021--}


  \begin{changemargin}{0in}{0.0in}
    \begin{itemize}

    \item \textsc{Physics Analysis}

    \vspace{0.1em}
    \begin{itemizetighter}

        \item \textit{Co-convener} LHCC Long-Lived Particle Working Group \when{2022--25}

        \item Heavy Stable Charged Particle Search \when{2022--}

        % \begin{itemizetightrightpad}
        % \item \it Searching for detector-stable particles carrying electric charge and leaving anomalous ionization and timing signatures in the CMS detector.
        % \end{itemizetightrightpad}


        \item Multijet Search \when{2023--}

        % \begin{itemizetightrightpad}
        % \item \it Searching for new physics in multijet final states using conventional techniques and machine learning
        % \end{itemizetightrightpad}

        \item Recursive Jigsaw Reconstruction Search \when{2022}

        % \begin{itemizetightrightpad}
        % \item \it Contributed to calculation of exclusion limits for a search for compressed supersymmetry spectra using the Recursive Jigsaw Reconstruction method.
        % \end{itemizetightrightpad}

    \end{itemizetighter}


    \item \textsc{Phase-II Outer Tracker Upgrade}

    \vspace{0.1em}
    \begin{itemizetighter}

      \item \textit{US CMS Level-3 Manager} for OT Electronics \when{2024--}
      \item \textit{US CMS Control Account Manager} for OT DAQ \when{2024--}

      \item Phase-II Tracker DAQ Software \when{2021--}

      \item Outer Tracker Module Assembly and Grading \when{2022--}

    \end{itemizetighter}


    \end{itemize}
  \end{changemargin}


  \item \textsc{Member of the ATLAS Collaboration} \when{2009--21}

  \begin{changemargin}{0in}{0.0in}
    \begin{itemize}
      % \parshape 2 3.2em \linewidth \listindent \dimexpr\linewidth-\listindent\relax

    % \item { \small [{\it Nota Bene:} Labels in \textsc{small caps} denote official ATLAS designations.] }


    % SUSY Software Review, Contours
    % Maintainer of big framework

    % NSW DAQ

    % \vspace{.5em}

    \item \textsc{Physics Analysis}

    \vspace{0.1em}
    \begin{itemizetighter}

        \item Common BSM Result Interpretation Framework \when{2018--21}

        % \begin{itemizetightrightpad}
        % \item \it Responsible for standards of exclusion interpretation and limit interpolation within the ATLAS SUSY working group, implementing new methods for presentation of results
        % \end{itemizetightrightpad}

        \item \textit{Co-convener} of SUSY RPV/LL sub-group \when{2017--19}

        % \begin{itemizetightrightpad}
        % \item \it Defined standards and direction of searches for long-lived particles and $R$-parity-violating supersymmetry. Served two terms.
        % \end{itemizetightrightpad}

        \item Searches for Long-Lived Particles \when{2016--21}

        % \begin{itemizetightrightpad}
        % \item \it Led many analysis teams in searches for long-lived particles in a variety of signatures. See Publications.
        % \end{itemizetightrightpad}
        % \pagebreak

        \item Common Analysis Software Development \when{2015--21}

        % \begin{itemizetightrightpad}
        % \item \it Primary developer of common analysis software used by roughly ten analysis teams
        % \end{itemizetightrightpad}


        \item Inclusive searches for Supersymmetry \when{2014--16}

        % \begin{itemizetightrightpad}
        % \item \it Led flagship search for supersymmetry in inclusive zero-lepton final states in 2016 including leading the first searches for new physics using ``Recursive Jigsaw'' variables. Created and commissioned novel triggers for use during Run-2 of the LHC geared toward finding new physics in low-mass-splitting scenarios.
        % \end{itemizetightrightpad}


        \item All-Hadronic BSM Signatures \when{2011--14}

        % \begin{itemizetightrightpad}
        % \item \it Led two searches for new phenomena in the context of $R$-Parity violating supersymmetry in fully hadronic final states
        % \end{itemizetightrightpad}


        \item Quark-vs-Gluon Jet Tagging \when{2010--12}

        % \begin{itemizetightrightpad}
        % \item \it Developed the first quark- vs. gluon-jet discriminator at a hadron collider
        % \end{itemizetightrightpad}

    \end{itemizetighter}

    \item \textsc{New Small Wheel Muon Spectrometer Upgrade}

    \vspace{0.1em}
    \begin{itemizetighter}

      \item \textit{Coordinator} Micromegas Trigger \when{2020}

      % \begin{itemizetightrightpad}
      % \item \it Integration of trigger electronics for the micromegas detector of the New Small Wheel (NSW) upgrade for the ATLAS muon spectrometer
      % \end{itemizetightrightpad}

      \item Online Software Coordination \when{2019--20} % \hspace{0.9em}}

      % \begin{itemizetightrightpad}
      % \item \it Designed, implemented, and/or maintain the readout, configuration, and calibration software systems for the NSW
      % \end{itemizetightrightpad}

      % \item Micromegas Digitization \when{2017--18}

      % \begin{itemizetightrightpad}
      % \item \it Responsible for digitization and simulation of the readout and trigger paths for the micromegas detector for the NSW
      % \end{itemizetightrightpad}

    \end{itemizetighter}


    \end{itemize}
  \end{changemargin}

  \item \textsc{Muon Collider Projects} \when{2020--}

  \begin{changemargin}{0in}{0.0in}
    \begin{itemizetight}

      \item \emph{USMCC Deputy Communication Coordinator} \when{2025--}
      \item 10 TeV MAIA Detector Concept Characterization \when{2023--}
      \item \emph{Taskforce Member} IMCC Software Taskforce \when{2024}

      \item Automating Professional 3D Rendering for Collider Event Displays \when{2023--}

      \item Simulation Software Workflow Development \when{2020--22}

    \end{itemizetight}
  \end{changemargin}


  \item \textsc{Accelerator Physics Projects} \when{2024--}

  \begin{changemargin}{0in}{0.0in}
    \begin{itemizetight}

    \item Helical FOFO Design for Muon Cooling for a Future Muon Collider \when{2024--}

    \item Using Self-Consistent Beam Distributions to Reduce Intra-Beam Stripping \when{2024--}

    \end{itemizetight}
  \end{changemargin}


  \item \textsc{External Collider Physics}


    \begin{itemizetight}

    \item Experimental Impact of Jet Fragmentation Reference Frames At Particle Colliders \when{2022--}


    \item Future collider projections for photon-induced production of charged BSM particles \when{2018--21}

    \item Bell Inequalities at Future $ee$ Colliders \when{2013--}

    \end{itemizetight}


  \item \textsc{Nucleon Physics Research} \when{2007--09}

  % \item Developed the prototype trigger and data acquisition systems for Fermilab E-906/SeaQuest \when{2008--2009}

  \begin{changemargin}{0in}{0.0in}
    \begin{itemizetightrightpad}
            \parshape 2 3.2em \linewidth \listindent \dimexpr\linewidth-\listindent\relax

    \item Developed the prototype trigger and data acquisition systems for Fermilab E-906/SeaQuest\vspace{0.2em}
    \item Assembled the photomultiplier tube units used in MINERvA\vspace{0.2em}
    \item Upgraded a Fabry-Perot cavity for use in a JLAB Compton polarimeter\vspace{0.2em}
    \end{itemizetightrightpad}
  \end{changemargin}

\end{itemize}


% \pagebreak

\section*{Publications}

  \begin{itemizetight}

  \input{LL_PubInclude.tex}

  \subsection*{One CMS Analysis Review Committee (ARC)}

  \subsection*{Nine ATLAS Editorial Boards (EBs)}

  % \subsection*{Journal Referee for Physical Review Letters, Physical Review D, Letters in High Energy Physics}

  % \begin{itemize}

  % \item Served on 8 ATLAS Editorial Boards

  % % \item {\it\href{https://atlas.web.cern.ch/Atlas/GROUPS/PHYSICS/CONFNOTES/ATLAS-CONF-2019-027/}{Soft $b$-hadron tagging for compressed SUSY scenarios}},~ATLAS-CONF-2019-027 (2019)

  % % \item {\it Search for new phenomena in a lepton plus high jet multiplicity final state with the ATLAS experiment using $\sqrt{s}$ = 13~TeV proton-proton collision data},~ATLAS-CONF-2016-094, ATLAS-CONF-2017-013, JHEP 09 (2017) 88

  % % \item {\it A search for pair-produced resonances in four jets final states at $\sqrt{s}=13$~TeV with the ATLAS detector},~ATLAS-CONF-2016-022 (2016)

  % % \item {\it\href{https://cds.cern.ch/record/2152392}{A search for R-parity violating decays of the top squark in four jet final states with the ATLAS detector at $\sqrt{s}=13$~TeV}},~ATLAS-CONF-2016-022 (2016)

  % % \item {\it\href{http://link.springer.com/article/10.1140/epjc/s10052-016-4065-1}{A new method to distinguish hadronically decaying boosted Z bosons from W bosons using the ATLAS detector}},~Eur. Phys. J. C \textbf{76} 238 (2016)

  % % \item {\it\href{https://cds.cern.ch/record/2017303}{Constraints on promptly decaying supersymmetric particles with lepton-number- and R-parity-violating interactions using Run-1 ATLAS data}},~ATLAS-CONF-2015-018 (2015)

  % \end{itemize}


  % \subsection*{Other Unpublished Work}

  % \begin{itemize}

  % % FQM Stuff
  % \item $\RHD$ GRW Theory and the Thermodynamic Arrow of Time (2010)
  % \item $\RHD$ Relativistic Considerations of the GRW Theory \emph{w/ R. Tumulka} (2008)
  % \item $\RHD$ Trouble with Locality: Adding Nonlocal Correlations \emph{w/ R. Tumulka} (2008)
  % \item $\RHD$ GRW Theory \emph{w/ R. Tumulka} (2008)

  % \end{itemize}


\end{itemizetight}

% \newpage

% \section*{Grants, Fellowships, \& Awards}

%   \begin{itemize}
%   \item George Southgate Fellowship, The University of Adelaide, 2015
%   \item 2nd International Spring School on Particle Physics and Philosophy Travel Award, 2014
%   \item AAAS/Science Program for Excellence in Science, 2013
%   % \item New York Academy of Science Sponsored Membership, 2012
%   \item Yale University Graduate Student Assembly Conference Travel Fund, 2012
%   % \item NSF Graduate Research Fellowship, Honorable Mention, 2011, 2010
%   \item Division of Nuclear Physics, American Physical Society Fall Meeting, Conference Travel Fund, 2008
%   % \item Summer Research Fellowship, Jefferson Laboratory, Rutgers University, 2008
%   \end{itemize}


\section*{Professional Activities}


  \begin{itemizetighter}
    \internalonly{
      \item \emph{Journal Referee:} PRD \when{2025}
    }
    \item RCSA Reviewer \when{2026} % For Cottrell
    \item Lead organizer, $\mu$CATALYST summer accelerator program \when{2026}
    \item Chair of Fermilab \emph{Colliders of Tomorrow} Committee \when{2025--}
    \item US LHC Users Association (USLUA) Executive Committee Member \when{2022--25}
    \item NSF Reviewer \when{2025} % For Kathy
    \item NSF Review Panel \when{2025}
    \item DOE Reviewer \when{2025}
    \item USCMS Phase-II Upgrade Level-3 Manager for Outer Tracker DAQ \when{2024--}
    \internalonly{
      \item \emph{Journal Referee:} PRD \when{2024}
      \item \emph{Journal Referee:} EPJC \when{2024}
    }

    \item Swiss National Science Foundation Reviewer \when{2024}
    \item International Muon Collider Collaboration Software Task Force \when{2024}
    \item US Muon Collider Collaboration Membership Coordinator \when{2024}
    \item US CMS (DOE+NSF) Project Review Panel \when{2024}
    \item NSF Reviewer ($\times$3 occasions) \when{2024}
    \item NSF Review Panel \when{2024}
    \internalonly{
      \item \emph{Journal Referee:} PRD \when{2024}
    }

    \item NSF Reviewer \when{2023}
    \item DOE Reviewer \when{2023}
    \item LHC LLP Working Group Convener \when{2022--25}
    \item SESAPS Nominating Committee \when{2022--23}
    \item Physics Today Consultant \when{2022}
    \internalonly{
      \item \emph{Journal Referee:} Universe \when{2022}
      \item \emph{Journal Referee:} Universe \when{2022}
      \item \emph{Journal Referee:} Instruments \when{2021}
      \item \emph{Journal Referee:} PRD \when{2021}
      \item \emph{Journal Referee:} LHEP \when{2021}
      \item \emph{Journal Referee:} PRL \when{2020}
      \item \emph{Journal Referee:} PRL \when{2018}
    }

    \externalonly{
      \item Journal Referee for Physical Review Letters, Physical Review D, \when{2018--}
      \item \hspace{2em} EPJC, Instruments, Letters in High Energy Physics, Universe
    }
  \end{itemizetighter}

\section*{Selected Scientific Presentations}

  \emph{Invited Seminars and Colloquiua}


  % \begin{itemize}
  \begin{itemizetighter}


    % \item University of Florida -- \textit{Physics Colloquium, Forthcoming} \when{2026}

    \item UC San Diego -- \textit{Physics Colloquium, Cottrell Lecture} \when{2026} % May 7, 2026

    \item San Diego State University -- \textit{Physics Colloquium, Cottrell Lecture} \when{2026} % May 8, 2026

    \item Cornell University -- \textit{HEP Seminar} \when{2026} % April 24, 2026

    \item Fermilab -- \textit{Wine \& Cheese Seminar} \when{2026} % 16 Jan 2026

    \item Yale University -- \textit{NPA Seminar} \when{2024} % 9 May 2024

    \item New York University -- \textit{Physics Colloquium} \when{2024} % March 2024

    \item University of California, Irvine, Virtual -- \textit{AI Seminar} \when{2022}
    % July 12, 2022

    \item University of Kansas, Lawrence, KS -- \emph{Particle Physics on the Plains} \when{2022}
    % March 8, 2022

    \item University of Tennessee, Knoxville, TN -- \emph{HEP Seminar} \when{2021}%
    % April 21, 2021
    % \item \textbf{R-Parity Violating SUSY is Hard to Find}
    %   \begin{itemizetighter}
    %   \item University of Tennessee, \emph{Seminar}, Knoxville, TN \when{2021}%
    %   \end{itemizetighter}

    \item University of Tennessee, Knoxville, TN -- \emph{Physics Colloquium} \when{2020}%
    %  Nov 23, 2020
    % \item \textbf{What To Do When Nature Is Unnatural}
    %   \begin{itemizetighter}
    %   \item University of Tennessee, \emph{Physics Colloquium}, Knoxville, TN \when{2020} % 11 December 2018
    %   \end{itemizetighter}

    \item Fermilab, Batavia, IL -- \emph{LPC Physics Forum} \when{2020}
    % 19 Nov 2020
    % \item \textbf{Displaced Vertex Searches at ATLAS}
    %   \begin{itemizetighter}
    %   \item Fermilab, \emph{LPC Physics Forum}, Batavia, IL \when{2020} % 11 December 2018
    %   \end{itemizetighter}


    \item Lund University, Sweden -- \emph{Seminar} \when{2020} % Mar 5, 2020
    \item Michigan State University, East Lansing, MI -- \emph{Physics Colloquium} \when{2019} % 5 February 2019
    \item University of Illinois, Urbana-Champaign, IL -- \emph{Physics Colloquium} \when{2018} % 25 October 2018
    \item Fermilab, Batavia, IL --  \emph{Topic of the Week} \when{2018} % 22 October 2018
    \item SLAC, Palo Alto, CA -- \emph{Seminar} \when{2018} % 18 October 2018
    \item University of Oregon, Eugene, OR -- \emph{Seminar} \when{2018} % 15 October 2018
    \item University of Pennsylvania, Philadelphia, PA -- \emph{Seminar} \when{2018} % 7 August 2018
    \item Rutgers University, Piscataway, NJ -- \emph{Seminar} \when{2018} % 6 August 2018

    % \item \textbf{How SUSY Can Still Save The Day / The SUSY Swindle}
    %   \begin{itemizetighter}
    %   \item Lund University, Sweden \when{2020}
    %   \item Michigan State University, East Lansing, MI \when{2019} % 5 February 2019
    %   \item University of Illinois, Urbana-Champaign, IL \when{2018} % 25 October 2018
    %   \item Fermilab \emph{Topic of the Week}, Batavia, IL \when{2018} % 22 October 2018
    %   \item SLAC, Palo Alto, CA \when{2018} % 18 October 2018
    %   \item University of Oregon, Eugene, OR \when{2018} % 15 October 2018
    %   \item University of Pennsylvania, Philadelphia, PA \when{2018} % 7 August 2018
    %   \item Rutgers University, Piscataway, NJ \when{2018} % 6 August 2018
    %   \end{itemizetighter}


    \item Harvard University, Cambridge, MA -- \emph{LPPC Seminar} \when{2018} % 11 October 2018
    % \item \textbf{On the Higgs Branching Ratios}
    %   \begin{itemizetighter}
    %   \item Harvard University - LPPC Seminar, Cambridge, MA\when{2018} % 11 October 2018
    %   \end{itemizetighter}


    \item Oskar Klein Centre, Stockholm, Sweden -- \emph{HEP Seminar}\when{2018} % 22 May 2018
    \item Lawrence Berkeley National Laboratory, Berkeley, CA  -- \emph{RPM Seminar}\when{2017} % 17 August 2017
    % \item \textbf{Searches for Sneaky SUSY at the ATLAS Experiment \& A Toy Entropic Explanation for the Hierarchy Problem}
    %   \begin{itemizetighter}
    %   \item Oskar Klein Centre, Stockholm, Sweden \when{2018} % 22 May 2018
    %   \item Lawrence Berkeley National Laboratory, Berkeley, CA \when{2017} % 17 August 2017
    %   \end{itemizetighter}

    \item Harvard University, Cambridge, MA -- \emph{LPPC Seminar}\when{2015} % 30 September 2015
  % \item \textbf{Searches for R-Parity Violating SUSY at the ATLAS Experiment}
  %   \begin{itemizetighter}
  %   \item Harvard University - LPPC Seminar, Cambridge, MA\when{2015}% 30 September 2015
  %   \end{itemizetighter}

    \item University of Adelaide, Australia -- \emph{CSSM Seminar}\when{2014} % 13 August 2014
    \item University of Melbourne, Australia -- \emph{Seminar}\when{2014} % 10 July 2014
    \item Institut de F\'isica d'Altes Energies (IFAE), Barcelona, Spain -- \emph{Seminar}\when{2014} % 7 January 2014

  % \item \textbf{A Search for B-Violating Supersymmetry in Multijet Signatures and Light-quark vs. Gluon Jet Tagging}
  %   \begin{itemizetighter}
  %   \item University of Adelaide, Australia\when{2014}% 13 August 2014
  %   \item University of Melbourne, Australia\when{2014}% 10 July 2014
  %   \item Institut de F\'isica d'Altes Energies (IFAE), Barcelona, Spain\when{2014} % 7 January 2014
  %   \item Santa Cruz Institute of Particle Physics Seminar, Geneva, Switzerland\when{2013}% 12
  %   \end{itemizetighter}

    \item 2nd International Spring School on Particle Physics and Philosophy -- \emph{Seminar}\when{2014}
  % \item \textbf{...But What If I Like Naturalness?}
  %   \begin{itemizetighter}
  %   \item 2nd International Spring School on Particle Physics and Philosophy
  %   \item Wuppertal, Germany\when{2014}% 31 March 2014
  %   \end{itemizetighter}

    \item Santa Cruz Institute of Particle Physics Seminar, Geneva, Switzerland -- \emph{Seminar}\when{2013} % 12

  % \item \textbf{RPV Stops at the LHC}
  %   \begin{itemizetighter}
  %   \item LPC Workshop on Exotic Top Partners
  %   \item Fermi National Accelerator Laboratory, LHC Physics Center, Batavia, IL\when{2013}% 26 September 2013
  %   \end{itemizetighter}


% \end{itemize}
\end{itemizetighter}

  \emph{Conference Presentations}

  % \begin{itemize}
  \begin{itemizetighter}

  \item DPF, \textit{Energy Frontier Session}, Fermilab \when{2026} % July 23 2026
  \item DPF, \textit{Outreach \& Engagement Session}, Fermilab \when{2026} % July 21 2026

  \item AI4MuC -- \textit{Overview Talk} \when{2026} % June 29, 2026


  \item IDEC Conference, Columbia College of Chicago  \when{2026} % Mar 2026

  \item Good $\nu$s, University of Pittsburgh -- \emph{Invited}  \when{2025} % 29 Oct 2025

  \item Physics At The Highest Energies With Colliders, Galileo Galilei Institute, Florence, Italy -- \emph{Invited}  \when{2025} % 29 July 2025


  \item LHCP2025, NTU, Taipei, Taiwan -- \emph{Invited}  \when{2025} % 4 March 2025


  \item Aspen Center for Physics -- \emph{Invited Muon Collider Overview Plenary}  \when{2025} % 4 March 2025

  \item Inaugural US Muon Collider Meeting, Fermilab -- \emph{Invited Experimental Overview Plenary}  \when{2024} % 7 August 2024

  \item Workshop on Long-Lived Particles, University of Tokyo, Tokyo, Japan -- \emph{o.b.o. CMS}  \when{2024} % 4 July 2024

  \item DPF/Pheno, Pittsburgh, PA \when{2024} % May 2024

  \item SESAPS2023, EKU, Richmond, KY -- \emph{Invited `Highlight' Talk}\when{2023} % XX Nov 2023

  \item Fermilab Users Meeting -- \emph{Invited Plenary} \when{2023}
  % June 28, 2023

  \item UCSB, Kavli Institute of Theoretical Physics, Muon Collider Workshop -- \emph{Invited Plenary} \when{2023}
  % March 8, 2023

  \item CMS Exotica Workshop -- \emph{Invited Plenary} \when{2022} % 11 Dec 2022

  \item ML4Jets, Rutgers University, Piscataway, NJ \when{2022} % 1 Nov 2022

  \item Snowmass Community Summer Study Workshop, Seattle, WA \when{2022} % 21 July 2022

  \item APS April Meeting, New York City, NY \when{2022} % 10 April 2022

  \item CMS Exotica Workshop -- \emph{Invited Plenary} \when{2021} % 24 Nov 2021

  \item SESAPS2021, FSU, Tallahassee, Fl -- \emph{Invited `Highlight' Talk}\when{2021} % 19 Nov 2021


  % \item \textbf{ColliderScope: Reaching New Audiences with Electronic(s) Music}
  %   \begin{itemizetighter}
  %   \item 18th International Particle Physics Outreach Group meeting, CERN \when{2019} % 28 November 2019
  %   \item International Conference on New Frontiers in Physics, Kolymbari, Crete \when{2019} % 21 August 2019
  %   \item APS Division of Particles and Fields Meeting, Boston, MA \when{2019} % 30 July 2019
  %   \end{itemizetighter}

  \item LHCP2021, Virtual -- \emph{Invited Plenary, o.b.o. ATLAS and CMS}\when{2021}
  % \item 6/9/2021 Neutral FIPs at the LHC and Beyond

  \item APS April Meeting, Virtual -- \emph{Invited Plenary, o.b.o. ATLAS}\when{2021}
  % \item \textbf{New Long-Lived Particle Results from the LHC} -- \emph{Plenary Talk} \when{2021} % 4/17/2021
  %   \begin{itemizetighter}
  %   \item APS April Meeting, Virtual
  %   \end{itemizetighter}

  \item A Rainbow Of Dark Sectors, Aspen Center for Physics, Virtual -- \emph{Invited Plenary}\when{2021}
  % \item \textbf{The Lifetime of the Long-Lived Particle Era at the LHC} -- \emph{Plenary Talk} \when{2021} % 24 Mar 2021
  %   \begin{itemizetighter}
  %   \item A Rainbow Of Dark Sectors, Aspen Center for Physics, Virtual
  %   \end{itemizetighter}


  \item LHC Experiments Committee (LHCC) Meeting, CERN, Geneva -- \emph{Open Session} \when{2020}
    % Nov 18, 2020
    % \item \textbf{Updates from the ATLAS Experiment}
    %   \begin{itemizetighter}
    %   \item LHC Experiments Committee (LHCC) Meeting Open Session, CERN, Geneva \when{2020}
    %   \end{itemizetighter}


  \item ICNFP, Kolymbari, Crete, Greece -- \emph{o.b.o. ATLAS}\when{2019} % 21 August 2019
  % \item \textbf{Searches for Long-Lived Particles with the ATLAS Detector}
  %   \begin{itemizetighter}
  %   \item \emph{On behalf of the ATLAS Collaboration}
  %   \item International Conference on New Frontiers in Physics, Kolymbari, Crete \when{2019} % 21 August 2019
  %   \end{itemizetighter}


  \item Higgs Cross-Section Working Group Meeting, CERN, Geneva -- \emph{Invited Plenary} \when{2018} % 11 December 2018
  % \item \textbf{Long-Lived Particles and the Higgs}
  %   \begin{itemizetighter}
  %   \item Higgs Cross-Section Working Group Meeting, CERN, Geneva \when{2018} % 11 December 2018
  %   \end{itemizetighter}

    \item USATLAS Workshop, Pittsburgh, PA -- \emph{Invited Plenary} \when{2018} % 1 August 2018
    % \item \textbf{SUSY and Exotics Overview} -- ``\emph{Headliner}'' Talk
    %   \begin{itemizetighter}
    %   \item USATLAS Workshop, University of Pittsburgh, Pittsburgh, PA \when{2018} % 1 August 2018
    %   \end{itemizetighter}

    % \item \textbf{The SUSY Paradox \& A Toy Entropic Explanation for the Hierarchy Problem}
    %   \begin{itemizetighter}
    %   \item Oskar Klein Centre, Stockholm, Sweden \when{2018} % 22 May 2018
    %   \end{itemizetighter}

  \item SUSY2017, Mumbai, India -- \emph{o.b.o. ATLAS} \when{2017} % 11 December 2017
  % \item \textbf{Searches for squarks and gluinos in scenarios with R-parity violating sparticle decays, or long- lived sparticles with ATLAS}
  %   \begin{itemizetighter}
  %   \item \emph{On behalf of the ATLAS Collaboration}
  %    \item SUSY2017, Mumbai, India \when{2017} % 11 December 2017
  %   \end{itemizetighter}

  \item IHEP-T2E, Kuala Lumpur, Malaysia -- \emph{Plenary, o.b.o. ATLAS} \when{2017} % 27 February 2017
  % \item \textbf{SUSY Searches in the ATLAS Experiment} -- \emph{Plenary Talk}
  %   \begin{itemizetighter}
  %   \item \emph{On behalf of the ATLAS Collaboration}
  %   \item IHEP-T2E, Kuala Lumpur, Malaysia \when{2017} % 27 February 2017
  %   \end{itemizetighter}

  \item SUSY2016, Melbourne, Victoria, Australia\when{2016} % 3 July 2016
  % \item \textbf{The Recursive Jigsaw Reconstruction Technique}
  %   \begin{itemizetighter}
  %   \item SUSY2016, Melbourne, Victoria, Australia \when{2016} % 3 July 2016
  %   \end{itemizetighter}

  % \item \textbf{New Searches for Strongly-Produced SUSY at the ATLAS Experiment}
  %   \begin{itemizetighter}
  %   \item CoEPP Annual Workshop, Torquay, Victoria, Australia\when{2016} % 18 February 2016
  %   \end{itemizetighter}

  \item SUSY2015, Lake Tahoe, CA\when{2015} % 24 August 2015
  % \item \textbf{Applications of the Recursive Jigsaw Technique}
  %   \begin{itemizetighter}
  %   \item SUSY2015, Lake Tahoe, CA\when{2015}% 24 August 2015
  %   \end{itemizetighter}

  \item Kruger2014, Kruger Gate, South Africa -- \emph{o.b.o. ATLAS} \when{2014} % 6 December 2014
  % \item \textbf{SUSY Searches in the ATLAS Experiment} -- \emph{Plenary Talk}
  %   \begin{itemizetighter}
  %   \item \emph{On behalf of the ATLAS Collaboration}
  %   \item Kruger2014, Kruger Gate, South Africa\when{2014}% 6 December 2014
  %   \end{itemizetighter}


  \item Fermilab, LPC Workshop on Exotic Top Partners, Batavia, IL -- \emph{Invited Plenary}\when{2013} % 26 September 2013

  \item SUSY2012, Peking University, Beijing, China -- \emph{o.b.o. ATLAS} \when{2012} % 17 August 2012
  % \item \textbf{Search for supersymmetry in resonant production and R-parity violating signatures with the ATLAS detector}
  %   \begin{itemizetighter}
  %   \item \emph{On behalf of the ATLAS Collaboration}
  %   \item SUSY2012, Peking University, Beijing, China\when{2012}% 17 August 2012
  %   \end{itemizetighter}

% \end{itemize}
\end{itemizetighter}


  \emph{Conference and Workshop Organization}

\begin{itemizetighter}

  \item \textit{USMCC 2026}, Stanford University, \emph{Co-chair (w/ M. Peskin), Program Committee} \when{2026} % Dec 2026
  \item \textit{USMCC 2025}, University of Chicago, IL, \emph{Detector Design and Performance Session Organizer [of 3]} \when{2025} % Aug 2025
  \item \textit{LHCP 2025}, Taipei, Taiwan, \emph{Poster Committee Member} \when{2025} % May 2025
  \item \textit{CPAD 2024}, Knoxville, TN, \emph{Chair of Local Organizing Committee [of 5]} \when{2024} % Nov 2024
  \item \textit{LHCP 2024}, Northeastern University, Boston, MA, \emph{Poster Committee Member} \when{2024} % May 2024

  \item \textit{The Future of High Energy Physics}, Aspen Center for Physics, \emph{Conference Organizer [of 4]} \when{2024} % Mar 2024
  \item \textit{PITT PACC Muon Collider Physics Benchmark Workshop}, Pittsburgh, PA, \emph{Invited Closing Remarks} \when{2023} % Nov 2023
  \item \textit{ML4Jets}, Rutgers University, \emph{Session Chair} \when{2022} % Oct 2022
  \item \textit{APS April Meeting}, New York, NY, \emph{Session Chair} \when{2022} % April 10 2022
  \item \textit{ATLAS SUSY and Exotics Workshop}, Virtual \when{2020} % September 2020
  \item \textit{ATLAS Reaching New Heights Workshop}, CERN \when{2019} % December 2019
  \item \textit{ATLAS SUSY Workshop}, Lecce, Italy \when{2019} % September 2019
  \item \textit{APS Division of Particles and Fields Meeting}, Northeastern University \when{2019} % 29 July - 2 August 2019
  \item \textit{ATLAS SUSY Workshop}, Stockholm, Sweden \when{2018} % 23-25 May 2018
  \item \textit{ATLAS SUSY and Exotics Workshop}, Bucharest, Romania \when{2017} % 8-12 May 2017
  \item \textit{Recursive Jigsaw Workshop}, Harvard University \when{2015} % 21 September 2015
  \item \textit{Workshop on LHC Searches}, LBNL \when{2014} % 30 January 2014

  \end{itemizetighter}

% \clearpage


\section*{Teaching}

  \begin{itemizetighter}
    \item Physics of Cosmic Rays \emph{Seminar \& Design/Build Course w/ UT Architecture \& Design} (PHYS 494) \when{2025}
    \item Classical Mechanics I \textit{\&} II (PHYS 311/312) \when{2021--24}
    \item UTK HEP/Astro Seminar    \when{2022}
    \item Graduate Research Participation Seminar  \when{2021}
  \end{itemizetighter}


  \internalonly{

    \section*{Current Student \textit{\&} Postdoc Mentorship}

    \begin{itemizetighter}
      \item Olivia Clark [Undergrad] \when{2025--}
      \item Micah Hillman [Undergrad] \when{2025--}
      \item \rightnote{UTK EUReCA 2nd Place Presenter Award}

      \item Allie Dabney [Undergrad] \when{2025--}
      \item Shivam [Ph.D.] \when{2024--}
      \item Caroline Riggall [Ph.D.] \when{2024--}
      \item Cordney Nash [Undergrad] \when{2024--}
      \item \rightnote{UTK `Volunteer of Distinction'}

      \item Colby Thompson [Ph.D.] \when{2023--}
      \item Adam Vendrasco [Ph.D.] \when{2022--}
      \item John Lawless [Ph.D.] \when{2022--}
      \item \rightnote{Stelson Fellowship for Outstanding Beginning Research}

      \item Emery Nibigira [Postdoc] \when{2022--}
      \item \rightnote{Fermilab LPC Distinguished Researcher}

      \item \rightnote{Lindau Nobel Laureate Meeting Selected Participant}

    \end{itemizetighter}


    \section*{Student Thesis Committee Membership}

    \begin{itemizetighter}
      \item Colby Thompson [Ph.D.], \emph{chair} \when{2025--}
      \item Raghav Chari [Undergrad, College Honors] \when{2025}
      \item John Lawless [Ph.D.], \emph{chair} \when{2024--}
      \item Brandi Skipworth [Ph.D.] \when{2024--}
      \item Christal Martin [Ph.D.] \when{2024--}
      \item Ewa Glimos [Ph.D.] \when{2024--}
      \item Taylor Sussmane [Undergrad, Honors], \emph{chair} \when{2024}
      \item John Melhorn [Undergrad, Honors] \when{2024}
      \item Luke Carpenter [Undergrad, Honors] \when{2024}
      \item Rebecca Godri [Ph.D.] \when{2023--}
      \item Daniel Ally [Ph.D.] \when{2022--}
      \item Jesse Harris [Ph.D.] \when{2022--}
      \item Satyajit Roy [Ph.D.] \when{2021--24}
      \item Austin Schmier [Ph.D.] \when{2021--24}
      \item Himal Acharya [Ph.D.] \when{2021}
      \item External Ph.D. Dissertation Review for Yale University \when{2021}
      \item External M.S. Thesis Review for University of Adelaide \when{2021}
    \end{itemizetighter}

  }

  \section*{Past Student \textit{\&} Postdoc Mentorship}


  \textit{n.b. Square brackets denote initial position of mentorship relationship. Current position in italics.}

    \begin{itemizetighter}
      \item Dr. Alexander Tuna [Postdoc, UTK] \hfill \emph{Research Staff, UCSD}
      \item Dr. Ankush Kanuganti [Postdoc, UTK] \hfill \emph{Postdoc, BNL}
      \item Micah Hillman [Undergrad, UTK] \hfill \emph{Physics Postbacc, Princeton}
      \item Charles Bell [Undergrad, UTK] \hfill \emph{Physics PhD Student, UMichigan, NSF GRFP Awardee, Cover of Science,}
        \item \hfill \emph{UTK Talley Award for Outstanding Undergrad Research,}
        \item \hfill \emph{UTK Douglas V. Roseberry Distinguished Upper Class Major Award}
      \item Lacey Dishman [Undergrad, UTK] \hfill \emph{Physics PhD Student, Boston University}
      \item Taylor Sussmane [Undergrad, UTK] \hfill \emph{Physics PhD Student, University of Wisconsin, Madison,}
        \item \hfill \emph{Douglas V. Roseberry Distinguished Upper Class Major Award,}
        \item \hfill \emph{UTK Talley Award for Outstanding Undergrad Leadership}
      \item Everett Hagen [Undergrad, UTK]
      \item Jack Peltier [Undergrad, UTK]\hfill \emph{Goldwater Scholar}
      \item Puck Cravens [Undergrad, UTK]
      \item Tanner Wolfe [Undergrad, UTK]
      \item Alexander Sizemore [Undergrad, UTK]
      \item Matthew Durkee [Undergrad, UTK] \hfill \emph{Defense Contractor}
      \item Philipp Wagenknecht [Ph.D., UTK] \hfill \emph{HEP Lab Technician, Brown University}
      \item Dr. Tamas Vami [Ph.D., Johns Hopkins] \hfill \emph{Postdoc, UCSB}
      \item Prof. Karri Folan DiPetrillo [Ph.D., Harvard] \hfill \emph{Asst. Professor, University of Chicago}
      \item Prof. Jennifer Roloff [Ph.D., Harvard] \hfill \emph{Asst. Professor, Brown University}
      \item Dr. Anthony Badea [Ph.D., Harvard] \hfill \emph{Schmidt AI Fellow, University of Chicago}
      \item Dr. Anne Fortman [Ph.D., Harvard] \hfill \emph{Postdoc, LBNL}
      \item Dr. Aaron White [Postdoc, Harvard] \hfill \emph{}
      \item Xiaohe Jia [Ph.D., Harvard] \hfill \emph{Industry}
      \item Madeline Bernstein [Undergrad, Harvard] \hfill \emph{Physics PhD Student, UC Berkeley}
      \item Dr. Stefanie Morgenstern [Ph.D., Dresden/CERN] \hfill \emph{CERN Fellow}
      \item Dr. Russell Smith [Ph.D., Columbia] \hfill \emph{Squarepoint, Financial Sector}
      \item Prof. Emma Alexander [Undergrad, Yale] \hfill \emph{Asst. Professor, Northwestern University}
      \item Dr. Ariel Ekblaw [Undergrad, Yale] \hfill \emph{Founder of the MIT Space Exploration Initiative, CEO of Aurelia Institute}
    \end{itemizetighter}


\section*{Selected Outreach Activities}

  \begin{itemize}

    \item \textbf{\href{https://cosmicrayn.utk.edu}{Cosmic Rayn: Cosmic Ray Educational Experience}} w/ UTK Architecture \textit{\&} Design \when{2024--26}

    \item \textbf{Quantum Canvases: Physics, the Arts, and the Humanities} \when{2023}

      \begin{itemizetighter}
      \item Lead organizer of a series of events in downtown Knoxville on the intersection of physics, the arts, and the humanities in collaboration with UTK Physics \textit{\&} Astronomy, CERN, and the UT Humanities Center.
      \item Electronic music event (``Harmonic Motion: Physics x Electronic Music'')
      \item Events around a collection of science fiction short stories
      \item A panel discussion on the intersection of physics and science fiction
      \item A humanities-facing physics colloquium event
      \end{itemizetighter}


    \item \textbf{Higgs Tenth Anniversary Celebration} \when{2022}

      \begin{itemizetighter}
      \item Primary organizer for a documentary screening of \emph{Particle Fever} and open panel session for the tenth anniversary of the discovery of the Higgs Boson. Held at Central Cinema in Knoxville, TN.
      \end{itemizetighter}


    \item \textbf{STEAM Trailer Design Project} \when{2021--22}

      \begin{itemizetighter}
      \item Worked with the UTK College of Architecture \textit{\&} Design to design and construct an educational experience for local middle schoolers with a focus on the physics of sound and visualization.
      \end{itemizetighter}

    \item \textbf{Keynote presentation at event for Eureka Science Museum}, Grand Junction, CO \when{2021}


    \item \href{http://www.cern.ch/ColliderScope}{\textbf{ColliderScope -- Particle-Physics-Inspired Electronic Music Project}} \when{2019--}

      \begin{itemizetighter}
      \item Music with sound waves that show particle physics images in Lissajous figures
      \item Live performances:
      \item \hspace{2em} Luck Reunion (at Willie Nelson's ranch), Austin, TX (\emph{w/ Micah Nelson, Daniel Lanois}) \when{2026}
      \item \hspace{2em} Fly By Night, Knoxville, TN (\emph{w/ signal:noise; VJ Set}) \when{2026}
      \item \hspace{2em} Fly By Night, Knoxville, TN (\emph{w/ signal:noise; VJ Set}) \when{2025}
      \item \hspace{2em} Fly By Night, Knoxville, TN (\emph{w/ signal:noise; VJ Set}) \when{2025}
      \item \hspace{2em} Scruffy City Hall, Knoxville, TN (\emph{w/ Harmonic Motion}) \when{2025}
      \item \hspace{2em} Scruffy City Hall, Knoxville, TN (\emph{w/ Harmonic Motion}) \when{2024}
      \item \hspace{2em} Old City Performing Arts Center, Knoxville, TN (\emph{w/ Harmonic Motion}) \when{2023}
      \item \hspace{2em} Blue Moon Tavern, Seattle, WA (\emph{w/ Snowmass summer meeting}) \when{2022}
      \item \hspace{2em} Roskilde Festival, Denmark \when{2022}
      \item \hspace{2em} Cross Club, Prague, YouTube, Live from the CERN Control Centre (\emph{w/ ICHEP2020}) \when{2020}
      \item \hspace{2em} Lund, Sweden \when{2020}
      \item \hspace{2em} Aeronaut Brewing, Somerville, MA \when{2019}
      \item \hspace{2em} ATLAS Experiment, Geneva, Switzerland \when{2019}
      \item \hspace{2em} CERN Open Days, Geneva, Switzerland \when{2019}
      \item \hspace{2em} Pohoda Music Festival, Tren\v{c}\'{i}n, Slovakia \when{2019}
      \item Commissioned Compositions:

      \item \hspace{2em} Intro Videos for the CERN Science Gateway Educational Program \when{2023}

      \item \hspace{2em} Intro Videos for CERN Music Festival Program \when{2019}

      \item Presented at multiple international and national conferences
    %   \item 18th International Particle Physics Outreach Group meeting, CERN \when{2019 % 28 November 2019}
    %   \item International Conference on New Frontiers in Physics, Kolymbari, Crete \when{2019 % 21 August 2019}
    %   \item APS Division of Particles and Fields Meeting, Boston, MA \when{2019 % 30 July 2019}
      \end{itemizetighter}


    \item \textbf{Presentations to Grade School Students}
    \begin{itemizetighter}
    \item L\&N STEM Academy, Knoxville, TN \when{2023}
    \item Tel Aviv University -- Future Scientists CERN Tour \when{2017--19}
    \item Five presentations to NJ and NY schools \when{2011--16}
    \end{itemizetighter}

    \item \textbf{Harvard NSW Activities at CERN} - \emph{Video} (\href{https://youtu.be/VHrJAkNsMhM}{YouTube})
    \begin{itemizetighter}
    \item Fully produced video for recruitment \when{2018}
    \end{itemizetighter}

  \item \href{http://www.cern.ch/inparticular}{\textbf{In Particular}} - \emph{Podcast}\when{2015--16}

    \begin{itemizetighter}
    \item Video, audio, and original music production
    % \item Production of ``Video Shorts''
    % \item Web development, graphic design, and episode production
    \item Tens of thousands of downloads ($\sim11$~TB served)
    \item Featured by CERN, iTunes, The Guardian Science Blogs, Ars Technica, Vice, and others
    \end{itemizetighter}


  \item \textbf{Visiting Scholar, Dharma Drum Mountain, Taiwan}
    \begin{itemizetighter}
    \item Several weeks of conversations with monastery's monks and academics \when{2014}
    \end{itemizetighter}


  \end{itemize}


\section*{Press \textit{\&} Interviews}

      \begin{itemizetighter}

      \item \href{https://physics.utk.edu/tova-holmes-and-larry-lee-selected-as-fermilab-distinguished-researchers/}{UTK Physics -- Tova Holmes And Larry Lee Selected As Fermilab Distinguished Researchers}\when{2026}

      \item \href{https://news.utk.edu/2025/02/12/university-of-tennessee-physicist-named-cottrell-scholar/}{UTK News -- University of Tennessee Physicist Named Cottrell Scholar}\when{2025}

      \item \emph{Science} -- L.~Lee, C.~Bell, \textbf{\href{https://www.science.org/toc/science/383/6690}{Cover of Issue 383 6690}}  \when{2024}


      \begin{itemizetightrightpad}
        \item \emph{Muon collider simulation work and renderings made in Unreal Engine were featured on the cover and in the cover article in the issue from March 29, 2024.}
      \end{itemizetightrightpad}

      \item \href{https://engage.aps.org/dpb/resources/newsletters/24-25-newsletter/f10-article}{APS DPB Annual Newsletter -- The Muon Collider: A Modern Machine for a Modern World}\when{2024}
      \item \href{https://cerncourier.com/a/a-new-generation-a-new-vision/}{CERN Courier -- A New Generation, A New Vision} \when{2024}
      \item \href{https://physics.utk.edu/the-particle-detectorists/}{UTK Physics -- The (Particle) Detectorists} \when{2024}
      \item \href{https://physics.utk.edu/four-ut-physics-students-win-doe-scgsr-support/}{UTK Physics -- NSF GRFP Awardee} \when{2024}
      \item \href{https://physics.utk.edu/the-art-of-muon-collisions/}{UTK Physics -- The Art of Muon Collisions} \when{2024}
      \item Interview for article for MIT Tech Review \when{2023}
      \item Interview for article for Symmetry Magazine on CERN music festival program \when{2023}
      \item Interview with Sparks Magazine on Asian Americans in STEM\when{2023}
      \item \href{https://physics.utk.edu/in-my-own-words-undergraduate-taylor-sussmane/}{UTK Physics -- In My Own Words: Undergraduate Taylor Sussmane} \when{2023}
      \item \href{https://physics.utk.edu/challenging-the-standard-model-lawrence-lee-wins-nsf-career-award/}{UTK Physics -- Challenging The Standard Model: Lawrence Lee Wins NSF CAREER Award}\when{2023}
      \item \href{https://physics.utk.edu/the-art-of-science/}{UTK Physics -- The Art of Science} \when{2023}
      \item \href{https://physics.utk.edu/the-yeti-returns/}{UTK Physics -- The YETI Returns} \when{2023}
      \item Interview with Danish podcast TechTopia \when{2022}
      \item \textit{ColliderScope} featured in \href{https://www.symmetrymagazine.org/article/lhc-music-through-the-colliderscope}{Symmetry Magazine} \when{2019}
      \item Interview for article for Symmetry Magazine on naturalness \when{2019}
      \item CERN Courier -- ATLAS Pushes the Limits on Supersymmetry \when{2019}

      \item Interview for Swedish-language podcast \textit{Professor Magenta} about science and art \when{2016}
      \item Public screening of film Particle Fever w/ panel Q\textit{\&}A\when{2014}


      \end{itemizetighter}


\section*{University, College, and Departmental Committees and Other Service Roles}

  \begin{itemizetighter}

    \item UTK Academy for Global Engagement - Mentor \when{2026--27}

    \item CAS Climate, Culture, and Community Committee \when{2025--}

    \item Faculty Search Committee - Particle Theory \when{2025--26}
    \item Faculty Search Committee - Lecturer \when{2025--26}

    \item ORIED Launch Pad Panel on Prestigious Fellowships \& Awards \when{2025}
    \item University NSF CAREER Workshop Panel \when{2025}
    \item Faculty Search Committee - Nuclear Theory \when{2024--25}

    \item Coordinator of Department Web Services \when{2024--}

    \item 4 UTK Senior Thesis Committees \when{2024--}

    \item University NSF CAREER Workshop Panel \when{2024}

    \item Provost Junior Faculty Fellows (\emph{UTK}) \when{2023--25}

    \item Graduate Recruitment Committee  \when{2023--}

    \item SPS Faculty Advisor  \when{2023--}

    \item Assessment Committee  \when{2023--24}

    \item Graduate Advising Committee \when{2022--}

    \item Planning Committee \when{2022--24}

    \item Faculty Search Committee - Medium Energy Nuclear Experiment \when{2022--23}

    \item Qualifying Exam Committee \when{2022--23, 24--25}

    \externalonly{
      \item 9 UTK PhD Thesis Committees \when{2021--}

    }

    \item Community Connections Committee (\emph{Chair}) \when{2021--23}

    \item Department Head Search Committee  (\emph{UTK College of Arts \& Sciences}) \when{2021--22}
  \end{itemizetighter}


% \section*{Professional Activities}

%   \begin{itemize}
%   \item Member, American Association for the Advancement of Science, 2013--Present
%   \item Member, New York Academy of Sciences, 2012--Present
%   \item Member, American Physical Society, 2008--2009
%   \end{itemize}

% \section*{Advising}


%   \begin{itemizetight}

%   \item Harvard University:
%   \begin{itemizetighter}
%     \item Brendon Bullard (PhD Candidate) \hfill[\emph{w/ M. Morii}]
%     \item Jennifer Roloff (2019 PhD) \hfill[\emph{w/ J. Huth}]
%     \item Karri Folan DiPetrillo (2019 PhD) \hfill[\emph{w/ M. Franklin}]
%     \item Adriana Rotaru (2019 Summer Undergraduate)  \hfill[\emph{w/ J. Huth}]
%     \item Madeline Bernstein (2017 Summer Undergraduate)  \hfill[\emph{w/ M. Franklin}]
%   \end{itemizetighter}

%   % \item Uppsala University, Students:
%   % \begin{itemizetighter}
%   %   \item Andreas Gillgren (2018-2019 Visiting Engineering Physics Student)  \hfill[\emph{w/ J. Huth}]
%   % \end{itemizetighter}

%   \item Columbia University:
%   \begin{itemizetighter}
%     \item Russell Smith (2016 PhD)  \hfill[\emph{w/ E. Hughes}]
%   \end{itemizetighter}

%   \item University of Adelaide, Three PhD and Two Undergraduate Students
%   % \begin{itemizetighter}
%   %     \item Jason Oliver (2019 PhD) \hfill[\emph{w/ P. Jackson}]
%   %     \item Anum Qureshi (2019 PhD) \hfill[\emph{'' ''}]
%   %     \item Damir Duvnjak (2014 Honours, 2019 PhD) \hfill[\emph{'' ''}]
%   %     \item Jens Kr\"oger (2015 Visiting) \hfill[\emph{'' ''}]
%   %     \item Finley Borgas (2014 Honours) \hfill[\emph{'' ''}]
%   % \end{itemizetighter}

%   \item Yale University, Five Undergraduate Honors Students
%   % \begin{itemizetight}
%   %     \item Madeleine Barrow (2012)
%   %     \item Ari Brill (2012)
%   %     \item Emma Alexander (2011)
%   %     \item James Giammona (2011)
%   %     \item Alexandra Turrini (2011)
%   % \end{itemizetight}
%   \end{itemizetight}

% \section*{Teaching}

% % \begin{changemargin}{1.5em}{0.0in}

%   \begin{itemizetight}
%     \item Yale University
%     \begin{itemizetight}
%       \item \emph{Fundamentals of Physics}, lecturing assistant for Prof. P. Tipton \when{2011}
%       \item \emph{Physics for the Life Sciences}, lecturing assistant for Prof. S. Mochrie \when{2010}
%       \item \emph{Modern Physical Measurement}, assistant for Prof. J. Harris \when{2010}
%       \item \emph{General Physics Laboratory}, assistant for Prof. K. Baker \when{2009}
%     \end{itemizetight}
%   \end{itemizetight}
% \end{changemargin}

% \clearpage
% \section*{References}

%   \begin{itemize}
%   \item \textbf{Prof. Melissa Franklin} - {\it Mallinckrodt Professor of Physics} - Harvard University
%   \item \hspace{1.6em} \href{mailto:Franklin@Physics.Harvard.edu}{Franklin@Physics.Harvard.edu}
%   \item \textbf{Prof. Tobias Golling} - {\it Associate Professor} - University of Geneva
%   \item \hspace{1.6em} \href{mailto:Tobias.Golling@unige.ch}{Tobias.Golling@unige.ch}
%   \item \textbf{Dr. Andreas Hoecker} - {\it Senior Research Physicist, ATLAS Spokesperson-Elect} - CERN
%   \item \hspace{1.6em} \href{mailto:Andreas.Hoecker@cern.ch}{Andreas.Hoecker@cern.ch}
%   \item \textbf{Prof. John Huth} - {\it Donner Professor of Science} - Harvard University
%   \item \hspace{1.6em} \href{mailto:Huth@g.Harvard.edu}{Huth@g.Harvard.edu}

%     % \pagebreak

%     % \item {\textit{\hspace{-1.6em} [References continued...] }}

%   \item \textbf{Prof. Paul Jackson} - {\it Associate Professor} - University of Adelaide
%   \item \hspace{1.6em} \href{mailto:P.Jackson@adelaide.edu.au}{P.Jackson@adelaide.edu.au}
%   \item \textbf{Dr. Zachary Marshall} - {\it Senior Scientist} - Lawrence Berkeley National Lab
%   \item \hspace{1.6em} \href{mailto:Zach.Marshall@cern.ch}{Zach.Marshall@cern.ch}

%   % \vspace{1em}
%   % \item \textit{Additional References}
%   % \begin{itemize}
%   % \item \textbf{Prof. Paul Tipton} - {\it Chair of Physics Department} - Yale University
%   % \item \hspace{1.6em} \href{mailto:paul.tipton@yale.edu}{Paul.Tipton@Yale.edu}
%   % \end{itemize}

%   \end{itemize}

\cvfooter

\end{document}
"
%%
%% which is what build.sh and .github/workflows/build-cv.yml do.

\newif\ifCVinternal
\ifdefined\CVinternal \CVinternaltrue \else \CVinternalfalse \fi

% \internalonly{<content>} -- included only in the internal edition.
\newcommand{\internalonly}[1]{\ifCVinternal #1\fi}

% \externalonly{<content>} -- included only in the external edition.
\newcommand{\externalonly}[1]{\ifCVinternal\else #1\fi}

% The internal edition says so in its PDF title, so the two files stay
% distinguishable once they have been detached from their filenames.
\ifCVinternal
  \def\LLedition{\name: Curriculum Vitae (internal edition)}
\else
  \def\LLedition{\name: Curriculum Vitae}
\fi

\hypersetup{
  colorlinks  = false,
  pdfauthor   = {\name},
  pdftitle    = {\LLedition},
  pdfsubject  = {Curriculum Vitae},
  pdfkeywords = {},
  pdfpagemode = UseNone,
}


%% ===========================================================================
%% Page layout
%% ===========================================================================

\geometry{
  body       = {6.5in, 9.0in},
  left       = 1.0in,
  top        = 1.0in,
  headheight = 14pt,   % enough room for the running header; silences fancyhdr
}

\pagestyle{fancy}
\fancyhf{}
\renewcommand{\headrulewidth}{0pt}
\fancyhead[L]{\name}
\fancyhead[R]{\thepage}

\setlength{\parindent}{0em}

\titlespacing{\section}{0pt}{0pt}{0pt}
\sectionfont{\rmfamily\mdseries\Large}
\subsectionfont{\rmfamily\mdseries\itshape\large}


%% ===========================================================================
%% Lists
%% ===========================================================================
%% The CV is built almost entirely from bullet-free lists that differ only in
%% how tightly they are set.  itemize itself is redefined, rather than
%% configured through enumitem, so that a plain \begin{itemize} anywhere picks
%% up the CV style.
%%
%%   itemize               default spacing between entries
%%   itemizetight          no extra space between entries
%%   itemizetighter        entries set as close as the leading allows
%%   itemizetightrightpad  tight, and inset from the right margin, for the
%%                         commentary set underneath an entry
%%
%% \LLlist{<left>}{<right>}{<itemsep>}{<parskip>}{<parsep>}

\newcommand{\LLlist}[5]{%
  \list{}{%
    \setlength{\leftmargin}{#1}%
    \setlength{\rightmargin}{#2}%
    \setlength{\itemsep}{#3}%
    \setlength{\parskip}{#4}%
    \setlength{\parsep}{#5}%
  }%
}

\renewenvironment{itemize}%
  {\LLlist{1.5em}{0em}{0.25em}{0pt}{0.25em}}{\endlist}
\newenvironment{itemizetight}%
  {\LLlist{1.5em}{0em}{0em}{-1pt}{0em}}{\endlist}
\newenvironment{itemizetighter}%
  {\LLlist{1.5em}{0em}{0em}{-0.4em}{0em}}{\endlist}
\newenvironment{itemizetightrightpad}%
  {\LLlist{1.5em}{3em}{0em}{-1pt}{0em}}{\endlist}

% Hanging indent for a run-on entry: \parshape 2 3.2em \linewidth \listindent
% \dimexpr\linewidth-\listindent\relax sets the first line full width and
% indents every line after it to \listindent.
\newlength{\listindent}
\setlength{\listindent}{4.8em}

% \changemargin{<left>}{<right>} ... \endchangemargin -- indents a block of the
% document; used to nest the collaboration sub-lists under their headings.
\def\changemargin#1#2{\list{}{\rightmargin#2\leftmargin#1}\item[]}
\let\endchangemargin=\endlist


%% ===========================================================================
%% The right-hand date column
%% ===========================================================================
%% Nearly every entry ends with a date set against the right margin.  \when
%% supplies both the \hfill and the fixed-width box that lines those dates up
%% into one column: the box is left-aligned, so "2024" and "2022--24" begin at
%% the same horizontal position.
%%
%% A \makebox is only as wide as it is declared, so an over-long date used to
%% overflow it silently and print out in the right margin -- six of the dates in
%% this CV ("2021--2026", "2022--23, 24--25", ...) are wider than the column.
%% \when measures the date first and sets an over-wide one flush right instead,
%% which costs the shared left edge on those few entries but keeps every date
%% inside the text block.

\newlength{\datecolumn}
\setlength{\datecolumn}{3.5em}

%% The glue ahead of a date is \datesep plus infinite stretch.  The stretch is
%% what pushes the date to the right margin; the fixed \datesep is a floor, so
%% an entry that grows long enough to reach its date reports an overfull box
%% instead of quietly printing its text hard against the year.  Every entry
%% currently clears it, so the floor costs nothing today -- it is there to catch
%% the next entry that does not.  \whenflush is the fix when one turns up: it
%% gives back the empty right-hand part of the date box, about 1.6em for a
%% four-digit year.
\newlength{\datesep}
\setlength{\datesep}{0.4em}

%% \LLdateglue is the glue that carries the date to the right margin.  It reads
%% as one \hfill most of the time, but is written out so that an entry too long
%% to hold its date can break: TeX takes the \penalty50, leaves the text on the
%% line it filled, and the \hbox{} -- which survives the break where glue would
%% not -- carries the second \hfill onto the next line, so the date still lands
%% flush right instead of stranded against the left margin.  \unskip drops the
%% source space ahead of the date, so the gap is \datesep however the entry was
%% typed.
\newcommand{\LLdateglue}{%
  \unskip\nobreak\hspace{\datesep}\hfill\penalty50 \hbox{}\nobreak\hfill
}

\newlength{\LLdatewidth}
\newcommand{\when}[1]{%
  \LLdateglue
  \settowidth{\LLdatewidth}{#1}%
  \ifdim\LLdatewidth>\datecolumn #1\else\makebox[\datecolumn][l]{#1}\fi
}

% \whenflush{<date>} -- for an entry whose text runs so close to the right
% margin that the shared left edge of the date column will not fit.  Sets the
% date flush right, giving up the alignment to keep the gap.
\newcommand{\whenflush}[1]{\LLdateglue #1}

% \rightnote{<text>} -- a right-aligned italic annotation (an award, a current
% position) hung underneath an entry.  The trailing pad makes it end where a
% four-digit year in the date column ends, rather than at the margin itself.
\newcommand{\rightnote}[1]{%
  \settowidth{\LLdatewidth}{2024}%
  \hfill\emph{#1}\hspace{\dimexpr\datecolumn-\LLdatewidth\relax}%
}


%% ===========================================================================
%% Live publication metrics from INSPIRE-HEP
%% ===========================================================================
%% Citation counts go stale the moment they are typed, so they are not written
%% into the CV by hand.  tools/fetch_inspire.py queries the INSPIRE-HEP API and
%% writes LL_InspireData.tex, a flat list of declarations:
%%
%%     \inspiresetcites{2690093}{3}         % citations of one paper
%%     \inspiresetstat{citations}{207,000}  % an author-level summary figure
%%
%% which the CV body reads back through \inspirepub and \inspirestat.  The
%% values set just below are fallbacks for a build that has never run the
%% fetcher -- an offline checkout, or Overleaf -- so the document always
%% compiles with plausible numbers.  build.sh and CI refresh the data first.

% Declarations written by tools/fetch_inspire.py.  The keys live in \csname, so
% the LL@ namespace prefix needs no \makeatletter here.
\newcommand{\inspiresetcites}[2]{\expandafter\gdef\csname LL@cites@#1\endcsname{#2}}
\newcommand{\inspiresetstat}[2]{\expandafter\gdef\csname LL@stat@#1\endcsname{#2}}

% Fallbacks: the last hand-checked figures, overridden by LL_InspireData.tex.
\inspiresetstat{papers}{1400}
\inspiresetstat{citations}{202,000}
\inspiresetstat{hindex}{209}

% \inspirestat{<key>} -- an author-level figure: papers, citations or hindex.
% A missing key is reported loudly rather than left silently blank.
\newcommand{\inspirestat}[1]{%
  \ifcsname LL@stat@#1\endcsname\csname LL@stat@#1\endcsname\else\textbf{[?~#1]}\fi
}

% \inspirecites{<record>} -- "[n citations]" for a record whose count is known.
% Nothing is printed for an unknown count, or for one below \citesmin, so a
% brand-new paper is left unannotated rather than advertising that nobody has
% cited it yet.
%
% The annotation costs about a dozen characters on every entry, which is enough
% to push an entry that already filled its line onto a second one.  Raising
% \citesmin, or shortening the label, is the dial if that happens too often.
\newcommand{\citesmin}{1}
\newcommand{\LLprintcites}[1]{%
  \ifnum#1<\citesmin\relax\else
    \nobreakspace{\small[#1\nobreakspace\ifnum#1=1\relax citation\else citations\fi]}%
  \fi
}
\newcommand{\inspirecites}[1]{%
  \ifcsname LL@cites@#1\endcsname\LLprintcites{\csname LL@cites@#1\endcsname}\fi
}

% \inspirepub{<record>}{<title>} -- a publication title in bold, linked to its
% INSPIRE-HEP record and followed by that record's live citation count.
\newcommand{\inspirepub}[2]{%
  \href{https://inspirehep.net/literature/#1}{\textbf{#2}}\inspirecites{#1}%
}

% \pub{<url>}{<title>} -- the same, for an entry with no INSPIRE-HEP record:
% arXiv-only preprints, CDS and CERN notes, bare DOI links.
\newcommand{\pub}[2]{\href{#1}{\textbf{#2}}}

% The generated data, if it has been fetched.  \IfFileExists keeps a fresh
% checkout compiling on the fallbacks above.
\IfFileExists{LL_InspireData.tex}{%% Publication metrics fetched from the INSPIRE-HEP API.
%% GENERATED FILE -- do not edit; run tools/fetch_inspire.py instead.
%%
%% \inspiresetcites{<record>}{<count>}  per-paper citation count
%% \inspiresetstat{<key>}{<value>}      author-level summary figure
%% Both are defined in LL_Preamble.tex.

%% Author-level summary, as shown at the top of the Publications section.
%% Papers and citations are rounded down because the CV says "over".
\inspiresetstat{papers}{1400}    % exact: 1,467
\inspiresetstat{citations}{207,000}  % exact: 207,783
\inspiresetstat{hindex}{211}

%% Per-paper citation counts, keyed by INSPIRE record id.
\inspiresetcites{3159704}{0}
\inspiresetcites{3137624}{2}
\inspiresetcites{3119390}{0}
\inspiresetcites{3103093}{4}
\inspiresetcites{2962374}{9}
\inspiresetcites{2958460}{1}
\inspiresetcites{2874678}{26}
\inspiresetcites{2840007}{27}
\inspiresetcites{2752508}{22}
\inspiresetcites{2690093}{3}
\inspiresetcites{2644916}{13}
\inspiresetcites{2642414}{406}
\inspiresetcites{2628398}{83}
\inspiresetcites{2005570}{23}
\inspiresetcites{1856547}{33}
\inspiresetcites{1831504}{102}
\inspiresetcites{1788448}{119}
\inspiresetcites{1756729}{145}
\inspiresetcites{1710078}{319}
\inspiresetcites{1701002}{245}
\inspiresetcites{1630632}{260}
\inspiresetcites{1609773}{260}
\inspiresetcites{1498566}{50}
\inspiresetcites{1458270}{135}
\inspiresetcites{1345355}{100}
\inspiresetcites{1298030}{152}
\inspiresetcites{1191022}{108}
}{%
  \typeout{^^JLL_CV: LL_InspireData.tex not found; using fallback INSPIRE
figures. Run tools/fetch_inspire.py to refresh them.^^J}%
}


%% ===========================================================================
%% Title block
%% ===========================================================================

% \cvheader -- the name, contact details and ORCID that open every document.
\newcommand{\cvheader}{%
  \begin{minipage}[t]{0.4\textwidth}
    {\huge \name}
  \end{minipage}%
  \begin{minipage}[t]{0.6\textwidth}
    \begin{flushright}
      \href{mailto:\email}{\email}\\
      \phone\\
      \href{https://\homepage}{\homepage}\\
      \href{https://orcid.org/\orcid}{\orcid}
      % A postal address, if one is ever wanted here:
      %   CERN CH-1211\\
      %   B\^atiment 40-1-D11\\
      %   23 Gen\`{e}ve, Switzerland
    \end{flushright}
  \end{minipage}%
}

% \cvfooter -- the "last updated" line that closes every document.
\newcommand{\cvfooter}{%
  \medskip
  \begin{center}
    \begin{small}
      Last updated: \today
    \end{small}
  \end{center}%
}
, and a body, so page
%% geometry, fonts, list styles and the custom commands below are only ever
%% edited in one place.  A document that wants to differ from a default
%% overrides it after the \input.
%%
%% Template attribution: Jason R. Blevins - Curriculum Vitae
%% Copyright (C) 2004-2013 Jason R. Blevins


%% ===========================================================================
%% Packages
%% ===========================================================================
%% hyperref goes last: it patches internals of most other packages.

\usepackage[T1]{fontenc}
\usepackage{CormorantGaramond}    % body and heading face
\usepackage{microtype}            % subtler justification, fewer overfull lines
\usepackage{wasysym}              % $\RHD$, the "lead role" marker
\usepackage{graphicx}             % for the (currently disabled) portrait

\usepackage{geometry}             % page size and margins
\usepackage{fancyhdr}             % running headers
\usepackage{lastpage}             % "page n of m"
\usepackage{titlesec}             % section spacing
\usepackage{sectsty}              % section fonts
\usepackage{enumitem}             % list configuration

\usepackage[hidelinks]{hyperref}  % clickable links, printed without decoration


%% ===========================================================================
%% Identity
%% ===========================================================================

\def\name{Lawrence Lee, PhD}
\def\email{Lawrence.Lee.Jr@cern.ch}
\def\phone{+1 856 765 4622}
\def\homepage{cern.ch/larry}
\def\orcid{0000-0002-5590-335X}


%% ===========================================================================
%% Internal and external editions
%% ===========================================================================
%% The internal edition carries material that is not for general circulation:
%% the current student and postdoc roster, thesis committee membership,
%% individual refereeing assignments, and student travel awards.  The external
%% edition is the default and needs no flag; the internal one is selected by
%% defining \CVinternal before the document is read:
%%
%%     pdflatex -jobname=LL_CV_internal "\def\CVinternal{}%% Lawrence Lee -- Curriculum Vitae
%%
%% Two editions are produced from this one source:
%%
%%   LL_CV.pdf           external, for general circulation
%%   LL_CV_internal.pdf  internal; additionally carries the current student and
%%                       postdoc roster, thesis committee membership, individual
%%                       refereeing assignments, and student travel awards
%%
%% Build both with ./build.sh.  Everything shared with LL_Publications.tex --
%% packages, page layout, list styles, the \internalonly / \externalonly
%% switches, and the live INSPIRE-HEP citation commands -- lives in
%% LL_Preamble.tex.  See README.md for the rest.

\documentclass[11pt,letterpaper]{article}

%!TEX root = LL_CV.tex
%%
%% Shared preamble for every document in this repository.
%%
%%   LL_CV.tex            the full CV, in an external and an internal edition
%%   LL_Publications.tex  the publication list on its own
%%
%% Both are just \documentclass, \input{LL_Preamble.tex}, and a body, so page
%% geometry, fonts, list styles and the custom commands below are only ever
%% edited in one place.  A document that wants to differ from a default
%% overrides it after the \input.
%%
%% Template attribution: Jason R. Blevins - Curriculum Vitae
%% Copyright (C) 2004-2013 Jason R. Blevins


%% ===========================================================================
%% Packages
%% ===========================================================================
%% hyperref goes last: it patches internals of most other packages.

\usepackage[T1]{fontenc}
\usepackage{CormorantGaramond}    % body and heading face
\usepackage{microtype}            % subtler justification, fewer overfull lines
\usepackage{wasysym}              % $\RHD$, the "lead role" marker
\usepackage{graphicx}             % for the (currently disabled) portrait

\usepackage{geometry}             % page size and margins
\usepackage{fancyhdr}             % running headers
\usepackage{lastpage}             % "page n of m"
\usepackage{titlesec}             % section spacing
\usepackage{sectsty}              % section fonts
\usepackage{enumitem}             % list configuration

\usepackage[hidelinks]{hyperref}  % clickable links, printed without decoration


%% ===========================================================================
%% Identity
%% ===========================================================================

\def\name{Lawrence Lee, PhD}
\def\email{Lawrence.Lee.Jr@cern.ch}
\def\phone{+1 856 765 4622}
\def\homepage{cern.ch/larry}
\def\orcid{0000-0002-5590-335X}


%% ===========================================================================
%% Internal and external editions
%% ===========================================================================
%% The internal edition carries material that is not for general circulation:
%% the current student and postdoc roster, thesis committee membership,
%% individual refereeing assignments, and student travel awards.  The external
%% edition is the default and needs no flag; the internal one is selected by
%% defining \CVinternal before the document is read:
%%
%%     pdflatex -jobname=LL_CV_internal "\def\CVinternal{}\input{LL_CV.tex}"
%%
%% which is what build.sh and .github/workflows/build-cv.yml do.

\newif\ifCVinternal
\ifdefined\CVinternal \CVinternaltrue \else \CVinternalfalse \fi

% \internalonly{<content>} -- included only in the internal edition.
\newcommand{\internalonly}[1]{\ifCVinternal #1\fi}

% \externalonly{<content>} -- included only in the external edition.
\newcommand{\externalonly}[1]{\ifCVinternal\else #1\fi}

% The internal edition says so in its PDF title, so the two files stay
% distinguishable once they have been detached from their filenames.
\ifCVinternal
  \def\LLedition{\name: Curriculum Vitae (internal edition)}
\else
  \def\LLedition{\name: Curriculum Vitae}
\fi

\hypersetup{
  colorlinks  = false,
  pdfauthor   = {\name},
  pdftitle    = {\LLedition},
  pdfsubject  = {Curriculum Vitae},
  pdfkeywords = {},
  pdfpagemode = UseNone,
}


%% ===========================================================================
%% Page layout
%% ===========================================================================

\geometry{
  body       = {6.5in, 9.0in},
  left       = 1.0in,
  top        = 1.0in,
  headheight = 14pt,   % enough room for the running header; silences fancyhdr
}

\pagestyle{fancy}
\fancyhf{}
\renewcommand{\headrulewidth}{0pt}
\fancyhead[L]{\name}
\fancyhead[R]{\thepage}

\setlength{\parindent}{0em}

\titlespacing{\section}{0pt}{0pt}{0pt}
\sectionfont{\rmfamily\mdseries\Large}
\subsectionfont{\rmfamily\mdseries\itshape\large}


%% ===========================================================================
%% Lists
%% ===========================================================================
%% The CV is built almost entirely from bullet-free lists that differ only in
%% how tightly they are set.  itemize itself is redefined, rather than
%% configured through enumitem, so that a plain \begin{itemize} anywhere picks
%% up the CV style.
%%
%%   itemize               default spacing between entries
%%   itemizetight          no extra space between entries
%%   itemizetighter        entries set as close as the leading allows
%%   itemizetightrightpad  tight, and inset from the right margin, for the
%%                         commentary set underneath an entry
%%
%% \LLlist{<left>}{<right>}{<itemsep>}{<parskip>}{<parsep>}

\newcommand{\LLlist}[5]{%
  \list{}{%
    \setlength{\leftmargin}{#1}%
    \setlength{\rightmargin}{#2}%
    \setlength{\itemsep}{#3}%
    \setlength{\parskip}{#4}%
    \setlength{\parsep}{#5}%
  }%
}

\renewenvironment{itemize}%
  {\LLlist{1.5em}{0em}{0.25em}{0pt}{0.25em}}{\endlist}
\newenvironment{itemizetight}%
  {\LLlist{1.5em}{0em}{0em}{-1pt}{0em}}{\endlist}
\newenvironment{itemizetighter}%
  {\LLlist{1.5em}{0em}{0em}{-0.4em}{0em}}{\endlist}
\newenvironment{itemizetightrightpad}%
  {\LLlist{1.5em}{3em}{0em}{-1pt}{0em}}{\endlist}

% Hanging indent for a run-on entry: \parshape 2 3.2em \linewidth \listindent
% \dimexpr\linewidth-\listindent\relax sets the first line full width and
% indents every line after it to \listindent.
\newlength{\listindent}
\setlength{\listindent}{4.8em}

% \changemargin{<left>}{<right>} ... \endchangemargin -- indents a block of the
% document; used to nest the collaboration sub-lists under their headings.
\def\changemargin#1#2{\list{}{\rightmargin#2\leftmargin#1}\item[]}
\let\endchangemargin=\endlist


%% ===========================================================================
%% The right-hand date column
%% ===========================================================================
%% Nearly every entry ends with a date set against the right margin.  \when
%% supplies both the \hfill and the fixed-width box that lines those dates up
%% into one column: the box is left-aligned, so "2024" and "2022--24" begin at
%% the same horizontal position.
%%
%% A \makebox is only as wide as it is declared, so an over-long date used to
%% overflow it silently and print out in the right margin -- six of the dates in
%% this CV ("2021--2026", "2022--23, 24--25", ...) are wider than the column.
%% \when measures the date first and sets an over-wide one flush right instead,
%% which costs the shared left edge on those few entries but keeps every date
%% inside the text block.

\newlength{\datecolumn}
\setlength{\datecolumn}{3.5em}

%% The glue ahead of a date is \datesep plus infinite stretch.  The stretch is
%% what pushes the date to the right margin; the fixed \datesep is a floor, so
%% an entry that grows long enough to reach its date reports an overfull box
%% instead of quietly printing its text hard against the year.  Every entry
%% currently clears it, so the floor costs nothing today -- it is there to catch
%% the next entry that does not.  \whenflush is the fix when one turns up: it
%% gives back the empty right-hand part of the date box, about 1.6em for a
%% four-digit year.
\newlength{\datesep}
\setlength{\datesep}{0.4em}

%% \LLdateglue is the glue that carries the date to the right margin.  It reads
%% as one \hfill most of the time, but is written out so that an entry too long
%% to hold its date can break: TeX takes the \penalty50, leaves the text on the
%% line it filled, and the \hbox{} -- which survives the break where glue would
%% not -- carries the second \hfill onto the next line, so the date still lands
%% flush right instead of stranded against the left margin.  \unskip drops the
%% source space ahead of the date, so the gap is \datesep however the entry was
%% typed.
\newcommand{\LLdateglue}{%
  \unskip\nobreak\hspace{\datesep}\hfill\penalty50 \hbox{}\nobreak\hfill
}

\newlength{\LLdatewidth}
\newcommand{\when}[1]{%
  \LLdateglue
  \settowidth{\LLdatewidth}{#1}%
  \ifdim\LLdatewidth>\datecolumn #1\else\makebox[\datecolumn][l]{#1}\fi
}

% \whenflush{<date>} -- for an entry whose text runs so close to the right
% margin that the shared left edge of the date column will not fit.  Sets the
% date flush right, giving up the alignment to keep the gap.
\newcommand{\whenflush}[1]{\LLdateglue #1}

% \rightnote{<text>} -- a right-aligned italic annotation (an award, a current
% position) hung underneath an entry.  The trailing pad makes it end where a
% four-digit year in the date column ends, rather than at the margin itself.
\newcommand{\rightnote}[1]{%
  \settowidth{\LLdatewidth}{2024}%
  \hfill\emph{#1}\hspace{\dimexpr\datecolumn-\LLdatewidth\relax}%
}


%% ===========================================================================
%% Live publication metrics from INSPIRE-HEP
%% ===========================================================================
%% Citation counts go stale the moment they are typed, so they are not written
%% into the CV by hand.  tools/fetch_inspire.py queries the INSPIRE-HEP API and
%% writes LL_InspireData.tex, a flat list of declarations:
%%
%%     \inspiresetcites{2690093}{3}         % citations of one paper
%%     \inspiresetstat{citations}{207,000}  % an author-level summary figure
%%
%% which the CV body reads back through \inspirepub and \inspirestat.  The
%% values set just below are fallbacks for a build that has never run the
%% fetcher -- an offline checkout, or Overleaf -- so the document always
%% compiles with plausible numbers.  build.sh and CI refresh the data first.

% Declarations written by tools/fetch_inspire.py.  The keys live in \csname, so
% the LL@ namespace prefix needs no \makeatletter here.
\newcommand{\inspiresetcites}[2]{\expandafter\gdef\csname LL@cites@#1\endcsname{#2}}
\newcommand{\inspiresetstat}[2]{\expandafter\gdef\csname LL@stat@#1\endcsname{#2}}

% Fallbacks: the last hand-checked figures, overridden by LL_InspireData.tex.
\inspiresetstat{papers}{1400}
\inspiresetstat{citations}{202,000}
\inspiresetstat{hindex}{209}

% \inspirestat{<key>} -- an author-level figure: papers, citations or hindex.
% A missing key is reported loudly rather than left silently blank.
\newcommand{\inspirestat}[1]{%
  \ifcsname LL@stat@#1\endcsname\csname LL@stat@#1\endcsname\else\textbf{[?~#1]}\fi
}

% \inspirecites{<record>} -- "[n citations]" for a record whose count is known.
% Nothing is printed for an unknown count, or for one below \citesmin, so a
% brand-new paper is left unannotated rather than advertising that nobody has
% cited it yet.
%
% The annotation costs about a dozen characters on every entry, which is enough
% to push an entry that already filled its line onto a second one.  Raising
% \citesmin, or shortening the label, is the dial if that happens too often.
\newcommand{\citesmin}{1}
\newcommand{\LLprintcites}[1]{%
  \ifnum#1<\citesmin\relax\else
    \nobreakspace{\small[#1\nobreakspace\ifnum#1=1\relax citation\else citations\fi]}%
  \fi
}
\newcommand{\inspirecites}[1]{%
  \ifcsname LL@cites@#1\endcsname\LLprintcites{\csname LL@cites@#1\endcsname}\fi
}

% \inspirepub{<record>}{<title>} -- a publication title in bold, linked to its
% INSPIRE-HEP record and followed by that record's live citation count.
\newcommand{\inspirepub}[2]{%
  \href{https://inspirehep.net/literature/#1}{\textbf{#2}}\inspirecites{#1}%
}

% \pub{<url>}{<title>} -- the same, for an entry with no INSPIRE-HEP record:
% arXiv-only preprints, CDS and CERN notes, bare DOI links.
\newcommand{\pub}[2]{\href{#1}{\textbf{#2}}}

% The generated data, if it has been fetched.  \IfFileExists keeps a fresh
% checkout compiling on the fallbacks above.
\IfFileExists{LL_InspireData.tex}{\input{LL_InspireData.tex}}{%
  \typeout{^^JLL_CV: LL_InspireData.tex not found; using fallback INSPIRE
figures. Run tools/fetch_inspire.py to refresh them.^^J}%
}


%% ===========================================================================
%% Title block
%% ===========================================================================

% \cvheader -- the name, contact details and ORCID that open every document.
\newcommand{\cvheader}{%
  \begin{minipage}[t]{0.4\textwidth}
    {\huge \name}
  \end{minipage}%
  \begin{minipage}[t]{0.6\textwidth}
    \begin{flushright}
      \href{mailto:\email}{\email}\\
      \phone\\
      \href{https://\homepage}{\homepage}\\
      \href{https://orcid.org/\orcid}{\orcid}
      % A postal address, if one is ever wanted here:
      %   CERN CH-1211\\
      %   B\^atiment 40-1-D11\\
      %   23 Gen\`{e}ve, Switzerland
    \end{flushright}
  \end{minipage}%
}

% \cvfooter -- the "last updated" line that closes every document.
\newcommand{\cvfooter}{%
  \medskip
  \begin{center}
    \begin{small}
      Last updated: \today
    \end{small}
  \end{center}%
}



\begin{document}
\thispagestyle{empty} % no running header on the first page

\cvheader

% Alternative title blocks, kept for reference:
%   \centerline{\huge \bf \name}          % name centred and bold
%   \begin{picture}(0,0)                  % portrait in the top right corner
%     \put(380,-30){\hbox{\includegraphics[height=12em]{Portrait.jpg}}}
%   \end{picture}

\bigskip

% Spacer that holds the vertical gap ahead of the first section.  It pairs with
% the summary paragraph below, which is currently disabled.
\begin{minipage}[t]{1.3em}
  \hspace{1em}
\end{minipage}
% \begin{minipage}[t]{0.75\textwidth}
% \emph{Broadly-trained particle physicist with expertise in searching for physics beyond the standard model at particle colliders. Particular expertise in supersymmetry, long-lived particles, and unconventional searches. Technical skills include expertise in detector data acquisition system, analysis software, and detector simulation.}
% \end{minipage}


\section*{Professional Experience}

  \begin{itemizetight}


    \item Associate Professor, University of Tennessee, Knoxville \when{2026--}
    \item Assistant Professor, University of Tennessee, Knoxville \when{2021--26}

    % \begin{itemizetighter}
    %   \item {University of Tennessee}, Knoxville, TN
    % \end{itemizetighter}

    \item Research Associate, Harvard University - \emph{Supervisor: J. Huth} \when{2016--21}

    % \begin{itemizetighter}
    %   \item {Harvard University}, Cambridge, MA -- \emph{Supervisor: J. Huth}
    % \end{itemizetighter}

    \item Research Associate, University of Adelaide - \emph{Supervisor: P. Jackson} \when{2014--16}

    % \begin{itemizetighter}
    %   \item {The University of Adelaide}, Adelaide, Australia -- \emph{Supervisor: P. Jackson}
    % \end{itemizetighter}

  \end{itemizetight}


\section*{Education}

  \begin{itemizetight}
    \item Ph.D., M.Phil., M.S., Physics, {Yale University} - \emph{Supervisor: T. Golling} \when{2014}
    % \begin{itemizetighter}
    %       \item \emph{Supervisor: T. Golling}
    % \end{itemizetighter}
    % \item M.S., M.Phil. Physics, {Yale University} \when{2012}
    \item B.S. Physics, {Rutgers University} - \emph{Supervisors: R. Ransome, R. Gilman, R. Tumulka} \when{2009}

  %   \begin{itemizetighter}
  %         \item \emph{Supervisors: R. Ransome, R. Gilman, R. Tumulka}
  %   \end{itemizetighter}
  \end{itemizetight}


% \section*{Fields of Research Interest}

%     \begin{itemizetightrightpad}
%     \item High Energy Collider Physics, Experimental searches for physics Beyond the Standard Model, Supersymmetry, Long-Lived Particles, Foundations of Quantum Mechanics
%     % \item Quantum Information, Philosophical Foundations of Quantum Mechanics
%     \end{itemizetightrightpad}


\section*{Awards \textit{\&} Honors}

\emph{External}

  \begin{itemizetighter}
    \item Cottrell Lectureship \when{2026}
    \item Fermilab LHC Physics Center Senior Distinguished Researcher \when{2026}
    \item Breakthrough Prize in Fundamental Physics Laureate  \when{2025}
    \item \hspace{2em}\emph{as part of the ATLAS and CMS Collaborations}
    \item Cottrell Scholar - Research Corporation for Science Advancement \when{2025}
    \item NSF CAREER Award \when{2023}
    \item European Physical Society's High Energy and Particle Physics Prize \when{2013}
    \item \hspace{2em}\emph{as part of the ATLAS Collaboration}
  \end{itemizetighter}


\emph{University of Tennessee, Knoxville}

\begin{itemizetighter}
  \item University of Tennessee Alumni Association -  Alumni Outstanding Teacher Award \when{2025}
  \item UTK Division of Student Success - Faculty Research Mentor Award \when{2025}
  \item UTK Provost - Junior Faculty Fellow \when{2023--25}
  \item UTK Physics - Research Advisor of the Year \when{2023--24}
  \item UTK College of Arts and Sciences - Faculty Academic Outreach Teaching Award \when{2022}
  \item UTK Physics - Teacher of the Year \when{2021--22}
\end{itemizetighter}


\section*{Funding}

\emph{Multi-Year Grants}

  \begin{itemizetighter}
    \item BNL Subcontract on Accelerator Physics - \$190k \when{2025--29}
    \item Cottrell Scholar - Research Corporation for Science Advancement - \$120k \when{2025--28}
    \item DOE Award for UTK HEP (4 PIs) - \$1.425M \when{2024--27}
    \item NSF CAREER Award - \$1M \when{2023--28}
    \item USCMS Upgrade Project Funding - \$73k \when{2022--24}
    \item Single-PI DOE Award - \$250k \when{2022--24}
  \end{itemizetighter}


  \emph{Other Funding}

  \begin{itemizetighter}
    \item Fermilab LHC Physics Center - \$56k\when{2026}
    \item Universities Research Association - \$7k\when{2026}
    \item Fermilab LHC Physics Center (Postdoc Emery Nibigira Led) - \$43k\when{2024}
    \item Universities Research Association (Postdoc Emery Nibigira Led) - \$20k\when{2024}
    \item UT-Battelle (ORNL Management Contractor) - \$74k \when{2024}
    \internalonly{
      \item UTK Graduate Student Senate Travel Award - John Lawless - \$1k\when{2024}
      \item SESAPS Travel Grant - Caroline Riggall - \$325 \when{2024}
      \item SESAPS Travel Grant - Charles Bell - \$500 \when{2023}
      \item SESAPS Travel Grant - Lacey Dishman - \$500 \when{2023}
      \item UTK Faculty Research Assistants Funding - Tanner Wolfe - \$2.3k \when{2023}
      \item FNAL LPC Guest \textit{\&} Visitor Funding - Emery Nibigira - \$5.2k \when{2023}
      \item FNAL LPC Guest \textit{\&} Visitor Funding - John Lawless - \$3.9k \when{2023}
      \item UTK Faculty Research Assistants Funding - Taylor Sussmane - \$2.3k \when{2022}
      \item SESAPS Travel Grant - Taylor Sussmane - \$500 \when{2022}
      \item SESAPS Travel Grant - Alexander Sizemore - \$500 \when{2022}
      \item FNAL LPC Guest \textit{\&} Visitor Funding - Ramon Ogaz - \$2.8k \when{2022}
      \item FNAL LPC Guest \textit{\&} Visitor Funding - Philipp Wagenknecht - \$5.2k \when{2022}
      \item UTK Physics Summer Research Fellowship - Matthew Durkee - \$6k \when{2022}
    }


    \item George Southgate Fellowship, The University of Adelaide - 5k AUD \when{2015}
  \end{itemizetighter}


\section*{Research Experience}

  \begin{itemize}


  \item \textsc{Member of the CMS Collaboration} \when{2021--}


  \begin{changemargin}{0in}{0.0in}
    \begin{itemize}

    \item \textsc{Physics Analysis}

    \vspace{0.1em}
    \begin{itemizetighter}

        \item \textit{Co-convener} LHCC Long-Lived Particle Working Group \when{2022--25}

        \item Heavy Stable Charged Particle Search \when{2022--}

        % \begin{itemizetightrightpad}
        % \item \it Searching for detector-stable particles carrying electric charge and leaving anomalous ionization and timing signatures in the CMS detector.
        % \end{itemizetightrightpad}


        \item Multijet Search \when{2023--}

        % \begin{itemizetightrightpad}
        % \item \it Searching for new physics in multijet final states using conventional techniques and machine learning
        % \end{itemizetightrightpad}

        \item Recursive Jigsaw Reconstruction Search \when{2022}

        % \begin{itemizetightrightpad}
        % \item \it Contributed to calculation of exclusion limits for a search for compressed supersymmetry spectra using the Recursive Jigsaw Reconstruction method.
        % \end{itemizetightrightpad}

    \end{itemizetighter}


    \item \textsc{Phase-II Outer Tracker Upgrade}

    \vspace{0.1em}
    \begin{itemizetighter}

      \item \textit{US CMS Level-3 Manager} for OT Electronics \when{2024--}
      \item \textit{US CMS Control Account Manager} for OT DAQ \when{2024--}

      \item Phase-II Tracker DAQ Software \when{2021--}

      \item Outer Tracker Module Assembly and Grading \when{2022--}

    \end{itemizetighter}


    \end{itemize}
  \end{changemargin}


  \item \textsc{Member of the ATLAS Collaboration} \when{2009--21}

  \begin{changemargin}{0in}{0.0in}
    \begin{itemize}
      % \parshape 2 3.2em \linewidth \listindent \dimexpr\linewidth-\listindent\relax

    % \item { \small [{\it Nota Bene:} Labels in \textsc{small caps} denote official ATLAS designations.] }


    % SUSY Software Review, Contours
    % Maintainer of big framework

    % NSW DAQ

    % \vspace{.5em}

    \item \textsc{Physics Analysis}

    \vspace{0.1em}
    \begin{itemizetighter}

        \item Common BSM Result Interpretation Framework \when{2018--21}

        % \begin{itemizetightrightpad}
        % \item \it Responsible for standards of exclusion interpretation and limit interpolation within the ATLAS SUSY working group, implementing new methods for presentation of results
        % \end{itemizetightrightpad}

        \item \textit{Co-convener} of SUSY RPV/LL sub-group \when{2017--19}

        % \begin{itemizetightrightpad}
        % \item \it Defined standards and direction of searches for long-lived particles and $R$-parity-violating supersymmetry. Served two terms.
        % \end{itemizetightrightpad}

        \item Searches for Long-Lived Particles \when{2016--21}

        % \begin{itemizetightrightpad}
        % \item \it Led many analysis teams in searches for long-lived particles in a variety of signatures. See Publications.
        % \end{itemizetightrightpad}
        % \pagebreak

        \item Common Analysis Software Development \when{2015--21}

        % \begin{itemizetightrightpad}
        % \item \it Primary developer of common analysis software used by roughly ten analysis teams
        % \end{itemizetightrightpad}


        \item Inclusive searches for Supersymmetry \when{2014--16}

        % \begin{itemizetightrightpad}
        % \item \it Led flagship search for supersymmetry in inclusive zero-lepton final states in 2016 including leading the first searches for new physics using ``Recursive Jigsaw'' variables. Created and commissioned novel triggers for use during Run-2 of the LHC geared toward finding new physics in low-mass-splitting scenarios.
        % \end{itemizetightrightpad}


        \item All-Hadronic BSM Signatures \when{2011--14}

        % \begin{itemizetightrightpad}
        % \item \it Led two searches for new phenomena in the context of $R$-Parity violating supersymmetry in fully hadronic final states
        % \end{itemizetightrightpad}


        \item Quark-vs-Gluon Jet Tagging \when{2010--12}

        % \begin{itemizetightrightpad}
        % \item \it Developed the first quark- vs. gluon-jet discriminator at a hadron collider
        % \end{itemizetightrightpad}

    \end{itemizetighter}

    \item \textsc{New Small Wheel Muon Spectrometer Upgrade}

    \vspace{0.1em}
    \begin{itemizetighter}

      \item \textit{Coordinator} Micromegas Trigger \when{2020}

      % \begin{itemizetightrightpad}
      % \item \it Integration of trigger electronics for the micromegas detector of the New Small Wheel (NSW) upgrade for the ATLAS muon spectrometer
      % \end{itemizetightrightpad}

      \item Online Software Coordination \when{2019--20} % \hspace{0.9em}}

      % \begin{itemizetightrightpad}
      % \item \it Designed, implemented, and/or maintain the readout, configuration, and calibration software systems for the NSW
      % \end{itemizetightrightpad}

      % \item Micromegas Digitization \when{2017--18}

      % \begin{itemizetightrightpad}
      % \item \it Responsible for digitization and simulation of the readout and trigger paths for the micromegas detector for the NSW
      % \end{itemizetightrightpad}

    \end{itemizetighter}


    \end{itemize}
  \end{changemargin}

  \item \textsc{Muon Collider Projects} \when{2020--}

  \begin{changemargin}{0in}{0.0in}
    \begin{itemizetight}

      \item \emph{USMCC Deputy Communication Coordinator} \when{2025--}
      \item 10 TeV MAIA Detector Concept Characterization \when{2023--}
      \item \emph{Taskforce Member} IMCC Software Taskforce \when{2024}

      \item Automating Professional 3D Rendering for Collider Event Displays \when{2023--}

      \item Simulation Software Workflow Development \when{2020--22}

    \end{itemizetight}
  \end{changemargin}


  \item \textsc{Accelerator Physics Projects} \when{2024--}

  \begin{changemargin}{0in}{0.0in}
    \begin{itemizetight}

    \item Helical FOFO Design for Muon Cooling for a Future Muon Collider \when{2024--}

    \item Using Self-Consistent Beam Distributions to Reduce Intra-Beam Stripping \when{2024--}

    \end{itemizetight}
  \end{changemargin}


  \item \textsc{External Collider Physics}


    \begin{itemizetight}

    \item Experimental Impact of Jet Fragmentation Reference Frames At Particle Colliders \when{2022--}


    \item Future collider projections for photon-induced production of charged BSM particles \when{2018--21}

    \item Bell Inequalities at Future $ee$ Colliders \when{2013--}

    \end{itemizetight}


  \item \textsc{Nucleon Physics Research} \when{2007--09}

  % \item Developed the prototype trigger and data acquisition systems for Fermilab E-906/SeaQuest \when{2008--2009}

  \begin{changemargin}{0in}{0.0in}
    \begin{itemizetightrightpad}
            \parshape 2 3.2em \linewidth \listindent \dimexpr\linewidth-\listindent\relax

    \item Developed the prototype trigger and data acquisition systems for Fermilab E-906/SeaQuest\vspace{0.2em}
    \item Assembled the photomultiplier tube units used in MINERvA\vspace{0.2em}
    \item Upgraded a Fabry-Perot cavity for use in a JLAB Compton polarimeter\vspace{0.2em}
    \end{itemizetightrightpad}
  \end{changemargin}

\end{itemize}


% \pagebreak

\section*{Publications}

  \begin{itemizetight}

  %!TEX root = LL_CV.tex
%%
%% The publication list, shared by LL_CV.tex and LL_Publications.tex.
%%
%% Titles are written with the commands defined in LL_Preamble.tex:
%%
%%   \inspirepub{<record>}{<title>}  a title linked to its INSPIRE-HEP record,
%%                                   annotated with that record's live citation
%%                                   count (see tools/fetch_inspire.py)
%%   \pub{<url>}{<title>}            a title linked anywhere else: an arXiv-only
%%                                   preprint, a CDS or CERN note, a bare DOI
%%   \inspirestat{<key>}             a live author-level figure used in the
%%                                   summary paragraph below
%%   \whenflush{<date>}              the date, flush right
%%
%% $\RHD$ marks a primary editor, primary analyzer, or intellectual lead role.

  \item  \href{https://inspirehep.net/authors/1071846}{Over \inspirestat{papers} published papers, mostly as part of the ATLAS and CMS Collaborations, with over \inspirestat{citations} citations, and an $h$-index of \inspirestat{hindex} [\textit{INSPIRE}].} Valid authorship in the ATLAS and CMS Collaborations requires a record of and continued contributions to the operation, maintenance, and improvements of the experiments in service of all of the collaboration's activities. Selected publications with significant direct contributions are listed below. \textit{{$\RHD$} indicates a primary editor, primary analyzer, or intellectual lead role.}


  \subsection*{Small Author List Publications}
  \begin{enumerate}

    \item S. Kakkar, L. Lee, N. Evans, A. Hoover, V. Morozov, \inspirepub{3119390}{Using Normalizing Flows in Normal-Form Space to Model Intra-Beam Stripping}, \emph{JACoW IPAC26} WEP4348 \newline \emph{(Proceedings, norm for publication in Accelerator Physics)} \whenflush{2026}

    \item $\RHD$ L.~Lee, J.~Lawless, C.~Riggall, \inspirepub{3137624}{Higgs Boson Spookiness: Probing Quantum Nonlocality with Spacetime-Resolved $H\!\to\!\tau^+\tau^-$ Decays}, \emph{arXiv}, 2603.28868 \whenflush{2026}

    \item $\RHD$ L.~Jeanty, L.~Lee, \inspirepub{3103093}{Rare and Experimentally Challenging Supersymmetry Signatures}, \emph{arXiv}, 2601.06358, \emph{Invited, Accepted by Annual Reviews} \whenflush{2026}

    \item C.~Riggall, R.~Ganesan, L.~Lee, T.~Holmes, K.~Yonehara, D.~Stratakis, D.~Neuffer, M.~Palmer, \inspirepub{3159704}{A Charge-Agnostic Design For 6D Muon Ionization Cooling}, \emph{JACoW COOL2025} \newline \emph{(Proceedings, norm for publication in Accelerator Physics)} \whenflush{2025}

    \item S. Kakkar, V. Morozov, L. Lee, N. Evans, A. Zhukov, A. Hoover, \inspirepub{3119390}{Beam Loss Modeling And Mitigation Due To
Intra-Beam Stripping}, \emph{JACoW NAPAC2025} THP014 \newline \emph{(Proceedings, norm for publication in Accelerator Physics)} \whenflush{2025}

    \item $\RHD$ T. Holmes, L.~Lee, \inspirepub{2958460}{Event Rates at a 10 TeV Muon Collider and Implications for Detector Design: Trigger, Data Acquisition, and Luminosity}, \emph{arXiv}, 2508.06239  \whenflush{2025}

    \begin{itemizetightrightpad}
      \item \emph{Discusses the experimental implications for electroweak process cross sections, particularly with respect to the collision rate and background process rates.}
    \end{itemizetightrightpad}

    \item $\RHD$ L.~Lee, C.~Bell, J.~Lawless, C.~Nash, E.~Nibigira, \inspirepub{2690093}{Experimental Impact of Jet Fragmentation Reference Frames At Particle Colliders}, \emph{PLB}, Volume 866, 139561  \whenflush{2025}

    \begin{itemizetightrightpad}
      \item \emph{Intellectual lead in argument and analysis demonstrating how modern jet physics relies on a frame-specific assumption that can fail in interesting corners of phase space.}
    \end{itemizetightrightpad}

  \item  C.~Bell, et al (\emph{21 others}) \inspirepub{2874678}{MAIA: A new detector concept for a 10 TeV muon collider},
  \\ \emph{arXiv:2502.00181}  \whenflush{2025}

  \begin{itemizetightrightpad}
    \item \emph{Main editor of Section VI. Responsible for all 3D renderings of the detector concept. General input on the overall physics output.}
  \end{itemizetightrightpad}


  \item C.~Accettura, et al (\emph{296 others}), \inspirepub{2642414}{Towards a Muon Collider}, \emph{Eur. Phys. J. C} 83, 864  \whenflush{2023}

  \begin{itemizetightrightpad}
    \item \emph{Contributed studies of Beam-Induced Background mitigation at a future muon collider.}
  \end{itemizetightrightpad}

  \item G.~Iakovidis, et al (\emph{100 others}), \inspirepub{2644916}{The New Small Wheel electronics}, \emph{JINST} 18 P05012 \whenflush{2023}

  \begin{itemizetightrightpad}
    \item \emph{Contributed to readout electronics integration with readout software and FELIX, integration of trigger hardware readout and dataflow, development of control and calibration software, and noise studies.}
  \end{itemizetightrightpad}

  \item $\RHD$ A.~Badea, W.~Fawcett, J.~Huth, T.J.~Khoo, R.~Poggi, L.~Lee, \inspirepub{2005570}{Solving Combinatorial Problems at Particle Colliders Using Machine Learning}, \emph{PRD} 106 1, 016001 \whenflush{2022}

  \begin{itemizetightrightpad}
    \item \emph{Designed physics-informed neural network to solve combinatorial ambiguities in collider events with many objects, especially for decays from pair-produced heavy resonances.}
  \end{itemizetightrightpad}

  \item C.~Bakalis, I.~Grayzman, L.~Lee, L.~Levinson, V.~Polychronakos, \pub{https://doi.org/10.1088/1748-0221/16/06/P06041}{Accessing register spaces in FPGAs within the ATLAS DAQ scheme via the SCA eXtension}, \emph{JINST} 16 P06041\whenflush{2021}

  \begin{itemizetightrightpad}
    \item \emph{Designed the software infrastructure around an FPGA implementation of a CERN GBT-SCA-like register interface for use in detector configuration.}
  \end{itemizetightrightpad}

  \item M.~Bauer, O.~Brandt, L.~Lee, C.~Ohm, \inspirepub{1756729}{ANUBIS: AN Underground Belayed In-Shaft search experiment}, \href{https://arxiv.org/abs/1909.13022}{arXiv: 1909.13022}\whenflush{2020}

  \begin{itemizetightrightpad}
    \item \emph{Responsible for original idea to instrument service shafts above the ATLAS Experiment for a dedicated long-lived particle experiment. Helped shape technical proposal outlined in paper.}
  \end{itemizetightrightpad}


  \item $\RHD$ L.~Lee, C.~Ohm, A.~Soffer, and T.~Yu, \inspirepub{1701002}{Collider Searches for Long-Lived Particles Beyond the Standard Model}, \emph{Progress in Particle and Nuclear Physics} 3695\whenflush{2019}


  \begin{itemizetightrightpad}
    \item \emph{Co-responsible for all experimental aspects of this broad review of collider searches for long-lived particles.}
  \end{itemizetightrightpad}


  \item X. Cid Vidal, et al., \inspirepub{1710078}{Beyond the Standard Model Physics at the HL-LHC and HE-LHC}, CERN Yellow Rep. Monogr. 7 (2019) 585-865 \whenflush{2019}

  \begin{itemizetightrightpad}
    \item \emph{Contributed studies of HL-LHC projections for long-lived gluino searches. Responsible for sensitivity projections.}
  \end{itemizetightrightpad}


  \end{enumerate}


  \subsection*{Selected Peer-Reviewed Publications from the ATLAS and CMS Collaborations}
  \begin{enumerate}[resume]
  \item $\RHD$ CMS Collaboration, \inspirepub{2962374}{A general search for supersymmetric particles in scenarios with compressed mass spectra using proton-proton collisions at $\sqrt{s}$ = 13 TeV}, \emph{Phys. Rev. D} 112 (2025) 11, 112023 \whenflush{2025}

  \begin{itemizetightrightpad}
    \item \emph{Limit setting and model interpretation.}
  \end{itemizetightrightpad}

  \item $\RHD$ CMS Collaboration, \inspirepub{2840007}{Search for heavy long-lived charged particles with large ionization energy loss in proton-proton collisions at $\sqrt{s}$ = 13 TeV}, \emph{JHEP} 04 (2025) 109  \whenflush{2025}

  \begin{itemizetightrightpad}
    \item \emph{Led multiple aspects of this analysis, particularly in ionization analysis design, simulation, and signal MC generation.}
  \end{itemizetightrightpad}

  \item $\RHD$ ATLAS Collaboration, \inspirepub{2752508}{A search for R-parity-violating supersymmetry in final states containing many jets in $pp$ collisions at $\sqrt{s} = 13$~TeV with the ATLAS detector}, \emph{JHEP} 05 (2024) 003  \whenflush{2024}

  \begin{itemizetightrightpad}
    \item \emph{Early designer of multijet analysis, software developer, and started the machine learning effort portion of this search.}
  \end{itemizetightrightpad}


 \item $\RHD$ ATLAS Collaboration, \inspirepub{2628398}{Search for long-lived, massive particles in events with displaced vertices and multiple jets in pp collisions at $\sqrt{s}=13$~TeV with the ATLAS detector}, \emph{JHEP} 06, 200 \whenflush{2023}

  \begin{itemizetightrightpad}
    \item \emph{Led initial design and implementation of a search for decays of BSM particles with displaced vertices and jets, including reconstruction algorithms, analysis software, and background estimation design.}
  \end{itemizetightrightpad}


 \item $\RHD$ ATLAS Collaboration, \inspirepub{1856547}{A search for the decays of stopped long-lived particles at $\sqrt{s}=13$~TeV with the ATLAS detector}, \emph{JHEP} 173 \whenflush{2021}

  \begin{itemizetightrightpad}
    \item \emph{Led search for very long lived particle decays in empty LHC bunch crossings.}
  \end{itemizetightrightpad}

  \item ATLAS Collaboration, \inspirepub{1831504}{Search for displaced leptons in $\sqrt{s} = 13$ TeV $pp$ collisions with the ATLAS detector}, \emph{Phys. Rev. Lett.} 127 051802 \whenflush{2021}

  \begin{itemizetightrightpad}
    \item \emph{Developed analysis software and provided design guidance for a search that extends the limits obtained from LEP for long-lived sleptons.}
  \end{itemizetightrightpad}


  \item $\RHD$ ATLAS Collaboration, \inspirepub{1788448}{Search for long-lived, massive particles in events with a displaced vertex and a muon with large impact parameter in $pp$ collisions at $\sqrt{s}=13$~TeV with the ATLAS detector}, \emph{Phys. Rev. D} 102 032006\whenflush{2020}

  \begin{itemizetightrightpad}
    \item \emph{Co-coordinator of analysis. Responsible for all aspects of the analysis, including background estimation scheme, physics targets, and analysis software.}
  \end{itemizetightrightpad}

  \item $\RHD$ ATLAS Collaboration, \inspirepub{1630632}{Search for long-lived, massive particles in events with displaced vertices and missing transverse momentum in 13~TeV pp collisions with the ATLAS detector}, \emph{Phys. Rev. D} 97 052012\whenflush{2018}

  \begin{itemizetightrightpad}
    \item \emph{Co-coordinator of analysis. Particularly responsible for background estimation methods and their validation, limit-setting, and data-flow model.}
  \end{itemizetightrightpad}

  \item ATLAS Collaboration, \inspirepub{1609773}{Search for new phenomena in high-mass diphoton final states using $37$~fb$^{-1}$ of proton-proton collisions collected at $\sqrt{s} = 13$~TeV with the ATLAS detector}, \emph{Phys. Lett. B} 775 105\whenflush{2017}

  \begin{itemizetightrightpad}
    \item \emph{Responsible for creation and validation of BSM signals and Higgs EFT models.}
  \end{itemizetightrightpad}

  \item ATLAS Collaboration, \inspirepub{1498566}{Search for new phenomena in events containing a same-flavour opposite-sign dilepton pair, jets, and large missing transverse momentum in $\sqrt{s} = 13$~TeV pp collisions with the ATLAS detector}, \emph{Eur. Phys. J C} \textbf{77} 144\whenflush{2017}

  \begin{itemizetightrightpad}
    \item \emph{Responsible for studies of the jet system dynamics.}
  \end{itemizetightrightpad}


    % \pagebreak

    % \item {\textit{\hspace{-1.6em} [Publications continued...] }}\vspace{-0.5em}

  \item $\RHD$ ATLAS Collaboration, \inspirepub{1458270}{Search for squarks and gluinos in final states with jets and missing transverse momentum at $\sqrt{s} = 13$~TeV with the ATLAS detector}, \emph{Eur. Phys. J. C} \textbf{76}: 392\whenflush{2016}

  \begin{itemizetightrightpad}
    \item \emph{Co-coordinator of search. Responsible for all aspects of limit-setting. Responsible for background estimation design and implementation for RJR analysis.}
  \end{itemizetightrightpad}


  \item $\RHD$ ATLAS Collaboration,
    \inspirepub{1345355}{Search for massive supersymmetric particles decaying to many jets using the ATLAS detector in pp collisions at $\sqrt{s} = 8$~TeV},
    \emph{Phys. Rev. D} \textbf{91}, 112016\whenflush{2015}

  \begin{itemizetightrightpad}
    \item \emph{Co-coordinator of search. Responsible for all aspects of the search from background estimation, signal simulation, and limit-setting.}
  \end{itemizetightrightpad}


  \item $\RHD$ ATLAS Collaboration, \inspirepub{1298030}{Light-quark and Gluon Jet Discrimination in pp Collisions at $\sqrt{s} = 7$~TeV with the ATLAS Detector},
    \emph{Eur. Phys. J. C} \textbf{74}: 3023\whenflush{2014}

  \begin{itemizetightrightpad}
    \item \emph{Responsible for analysis design, data-driven template determination and validation.}
  \end{itemizetightrightpad}

  \item $\RHD$ ATLAS Collaboration, \inspirepub{1191022}{Search for pair production of massive particles decaying into three quarks with the ATLAS detector in $\sqrt{s} = 7$~TeV pp collisions at the LHC},~Journal of High Energy Physics 12, 1--42\whenflush{2012}


  \begin{itemizetightrightpad}
    \item \emph{Responsible for all signal simulation, limit-setting, and interpretation.}
  \end{itemizetightrightpad}

  \end{enumerate}


  \subsection*{Selected Public Documents}

  \begin{itemize}


  \item A. Affolder, et al., \pub{https://arxiv.org/abs/2209.03607}{Solid State Detectors and Tracking for Snowmass},
    \emph{Snowmass 2021}, \whenflush{2022}

  \item $\RHD$ D. Ally, et al., \pub{https://arxiv.org/abs/2203.06773}{Strategies for Beam-Induced Background Reduction at Muon Colliders},
    \emph{Snowmass 2021}, \whenflush{2022}

  \item $\RHD$ J. Cochran, et al., \pub{https://arxiv.org/abs/2203.09585}{Particle Physics Outreach at Non-traditional Venues},
    \emph{Snowmass 2021}, \whenflush{2022}

  \item K. Black, et al., \pub{https://arxiv.org/abs/2203.08874}{Prospects for the Measurement of the Standard Model Higgs Pair Production at the Muon Colliders}, \emph{Snowmass 2021}, \whenflush{2022}

  \item N. Bartosik, et al., \pub{https://arxiv.org/abs/2203.07964}{Simulated Detector Performance at the Muon Collider}, \emph{Snowmass 2021}, \whenflush{2022}

  \item J. De Blas, et al., \pub{https://arxiv.org/abs/2203.07261}{The physics case of a 3 TeV muon collider stage}, \emph{Snowmass 2021}, \whenflush{2022}

  \item C. Aime, et al., \pub{https://arxiv.org/abs/2203.07256}{Muon Collider Physics Summary},
    \emph{Snowmass 2021}, \whenflush{2022}

  \item S. Jindariani, et al., \pub{https://arxiv.org/abs/2203.07224}{Promising Technologies and R\&D Directions for the Future Muon Collider Detectors}, \emph{Snowmass 2021}, \whenflush{2022}


  \item ATLAS Collaboration, \pub{https://atlas.web.cern.ch/Atlas/GROUPS/PHYSICS/PUBNOTES/ATL-PHYS-PUB-2019-019/}{Generation and Simulation of $R$-Hadrons},~ATL-PHYS-PUB-2019-019\whenflush{2019}
  % \item $\RHD$ {\it\href{https://atlas.web.cern.ch/Atlas/GROUPS/PHYSICS/CONFNOTES/ATLAS-CONF-2019-006/}{Search for long-lived, massive particles in events with a displaced vertex and a displaced muon in $pp$ collisions at $\sqrt{s}=13$~TeV with the ATLAS detector}},~ATLAS-CONF-2019-006 (2019)
  \item $\RHD$ ATLAS Collaboration, \pub{https://atlas.web.cern.ch/Atlas/GROUPS/PHYSICS/PUBNOTES/ATL-PHYS-PUB-2019-012/}{Supersymmetry Summary Plots}, Main overview and RPV plots \whenflush{2018--20}
  \item ATLAS Collaboration, \pub{https://atlas.web.cern.ch/Atlas/GROUPS/PHYSICS/PUBNOTES/ATL-PHYS-PUB-2019-013/}{Performance of vertex reconstruction algorithms for detection of new long-lived particle decays within the ATLAS inner detector},~ATL-PHYS-PUB-2019-013\whenflush{2019}
  % \item {\it\href{https://atlas.web.cern.ch/Atlas/GROUPS/PHYSICS/PUBNOTES/ATL-PHYS-PUB-2018-033/}{Sensitivity of the ATLAS experiment to long-lived particles with a displaced vertex and $E_T^{\mbox{miss}}$ signature at the HL-LHC}},~ATL-PHYS-PUB-2018-033 (2018)
  \item ATLAS Collaboration, \pub{https://atlas.web.cern.ch/Atlas/GROUPS/PHYSICS/CONFNOTES/ATLAS-CONF-2018-003/}{Reinterpretation of searches for supersymmetry in models with variable R-parity-violating coupling strength and long-lived R-hadrons},~ATLAS-CONF-2018-003\whenflush{2018}
  % \item $\RHD$ {\it\href{https://atlas.web.cern.ch/Atlas/GROUPS/PHYSICS/CONFNOTES/ATLAS-CONF-2017-026/}{Search for long-lived, massive particles in events with displaced vertices and missing transverse momentum in 13~TeV pp collisions with the ATLAS detector}},~ATLAS-CONF-2017-026 (2017)
  % \item {\it\href{https://atlas.web.cern.ch/Atlas/GROUPS/PHYSICS/CONFNOTES/ATLAS-CONF-2017-022/}{Search for squarks and gluinos in final states with jets and missing transverse momentum using $36$~fb$^{-1}$ of $\sqrt{s} = 13$~TeV pp collision data with the ATLAS detector.}}~ATLAS-CONF-2017-022 (2017)
  % \item $\RHD$ {\it\href{https://atlas.web.cern.ch/Atlas/GROUPS/PHYSICS/CONFNOTES/ATLAS-CONF-2016-078/}{Further searches for squarks and gluinos in final states with jets and missing transverse momentum at $\sqrt{s} = 13$~TeV with the ATLAS detector}},~ATLAS-CONF-2016-078 (2016)
  % \item {\it\href{https://atlas.web.cern.ch/Atlas/GROUPS/PHYSICS/CONFNOTES/ATLAS-CONF-2016-059/}{Search for scalar diphoton resonances with $15.4$~fb$^{-1}$ of data collected at $\sqrt{s} = 13$~TeV in 2015 and 2016 with the ATLAS detector}},~ATLAS-CONF-2016-059 (2016)
  % \item {\it\href{https://cds.cern.ch/record/2114828}{Search for squarks and gluinos in final states with jets and missing transverse momentum at $\sqrt{s} = 13$~TeV with the ATLAS detector}},~ATLAS-CONF-2015-062 (2015)
  \item ATLAS Collaboration, \pub{https://cds.cern.ch/record/2037905}{First look at pp collisions data at $\sqrt{s} = 13$~TeV in preparation for a search for squarks and gluinos in final states with jets and missing transverse momentum with the ATLAS detector},~ATL-PHYS-PUB-2015-028\whenflush{2015} \vspace{-1em}
  % \item $\RHD$ {\it\href{https://cds.cern.ch/record/1558689}{Search for massive particles decaying into multiple quarks with the ATLAS detector in $\sqrt{s} = 8$~TeV pp collisions}},~ATLAS-CONF-2013-091 (2013)
  \vspace{1em}
  \item \textit{Additional public documents supporting the publications described above.}

  \end{itemize}


  \subsection*{One CMS Analysis Review Committee (ARC)}

  \subsection*{Nine ATLAS Editorial Boards (EBs)}

  % \subsection*{Journal Referee for Physical Review Letters, Physical Review D, Letters in High Energy Physics}

  % \begin{itemize}

  % \item Served on 8 ATLAS Editorial Boards

  % % \item {\it\href{https://atlas.web.cern.ch/Atlas/GROUPS/PHYSICS/CONFNOTES/ATLAS-CONF-2019-027/}{Soft $b$-hadron tagging for compressed SUSY scenarios}},~ATLAS-CONF-2019-027 (2019)

  % % \item {\it Search for new phenomena in a lepton plus high jet multiplicity final state with the ATLAS experiment using $\sqrt{s}$ = 13~TeV proton-proton collision data},~ATLAS-CONF-2016-094, ATLAS-CONF-2017-013, JHEP 09 (2017) 88

  % % \item {\it A search for pair-produced resonances in four jets final states at $\sqrt{s}=13$~TeV with the ATLAS detector},~ATLAS-CONF-2016-022 (2016)

  % % \item {\it\href{https://cds.cern.ch/record/2152392}{A search for R-parity violating decays of the top squark in four jet final states with the ATLAS detector at $\sqrt{s}=13$~TeV}},~ATLAS-CONF-2016-022 (2016)

  % % \item {\it\href{http://link.springer.com/article/10.1140/epjc/s10052-016-4065-1}{A new method to distinguish hadronically decaying boosted Z bosons from W bosons using the ATLAS detector}},~Eur. Phys. J. C \textbf{76} 238 (2016)

  % % \item {\it\href{https://cds.cern.ch/record/2017303}{Constraints on promptly decaying supersymmetric particles with lepton-number- and R-parity-violating interactions using Run-1 ATLAS data}},~ATLAS-CONF-2015-018 (2015)

  % \end{itemize}


  % \subsection*{Other Unpublished Work}

  % \begin{itemize}

  % % FQM Stuff
  % \item $\RHD$ GRW Theory and the Thermodynamic Arrow of Time (2010)
  % \item $\RHD$ Relativistic Considerations of the GRW Theory \emph{w/ R. Tumulka} (2008)
  % \item $\RHD$ Trouble with Locality: Adding Nonlocal Correlations \emph{w/ R. Tumulka} (2008)
  % \item $\RHD$ GRW Theory \emph{w/ R. Tumulka} (2008)

  % \end{itemize}


\end{itemizetight}

% \newpage

% \section*{Grants, Fellowships, \& Awards}

%   \begin{itemize}
%   \item George Southgate Fellowship, The University of Adelaide, 2015
%   \item 2nd International Spring School on Particle Physics and Philosophy Travel Award, 2014
%   \item AAAS/Science Program for Excellence in Science, 2013
%   % \item New York Academy of Science Sponsored Membership, 2012
%   \item Yale University Graduate Student Assembly Conference Travel Fund, 2012
%   % \item NSF Graduate Research Fellowship, Honorable Mention, 2011, 2010
%   \item Division of Nuclear Physics, American Physical Society Fall Meeting, Conference Travel Fund, 2008
%   % \item Summer Research Fellowship, Jefferson Laboratory, Rutgers University, 2008
%   \end{itemize}


\section*{Professional Activities}


  \begin{itemizetighter}
    \internalonly{
      \item \emph{Journal Referee:} PRD \when{2025}
    }
    \item RCSA Reviewer \when{2026} % For Cottrell
    \item Lead organizer, $\mu$CATALYST summer accelerator program \when{2026}
    \item Chair of Fermilab \emph{Colliders of Tomorrow} Committee \when{2025--}
    \item US LHC Users Association (USLUA) Executive Committee Member \when{2022--25}
    \item NSF Reviewer \when{2025} % For Kathy
    \item NSF Review Panel \when{2025}
    \item DOE Reviewer \when{2025}
    \item USCMS Phase-II Upgrade Level-3 Manager for Outer Tracker DAQ \when{2024--}
    \internalonly{
      \item \emph{Journal Referee:} PRD \when{2024}
      \item \emph{Journal Referee:} EPJC \when{2024}
    }

    \item Swiss National Science Foundation Reviewer \when{2024}
    \item International Muon Collider Collaboration Software Task Force \when{2024}
    \item US Muon Collider Collaboration Membership Coordinator \when{2024}
    \item US CMS (DOE+NSF) Project Review Panel \when{2024}
    \item NSF Reviewer ($\times$3 occasions) \when{2024}
    \item NSF Review Panel \when{2024}
    \internalonly{
      \item \emph{Journal Referee:} PRD \when{2024}
    }

    \item NSF Reviewer \when{2023}
    \item DOE Reviewer \when{2023}
    \item LHC LLP Working Group Convener \when{2022--25}
    \item SESAPS Nominating Committee \when{2022--23}
    \item Physics Today Consultant \when{2022}
    \internalonly{
      \item \emph{Journal Referee:} Universe \when{2022}
      \item \emph{Journal Referee:} Universe \when{2022}
      \item \emph{Journal Referee:} Instruments \when{2021}
      \item \emph{Journal Referee:} PRD \when{2021}
      \item \emph{Journal Referee:} LHEP \when{2021}
      \item \emph{Journal Referee:} PRL \when{2020}
      \item \emph{Journal Referee:} PRL \when{2018}
    }

    \externalonly{
      \item Journal Referee for Physical Review Letters, Physical Review D, \when{2018--}
      \item \hspace{2em} EPJC, Instruments, Letters in High Energy Physics, Universe
    }
  \end{itemizetighter}

\section*{Selected Scientific Presentations}

  \emph{Invited Seminars and Colloquiua}


  % \begin{itemize}
  \begin{itemizetighter}


    % \item University of Florida -- \textit{Physics Colloquium, Forthcoming} \when{2026}

    \item UC San Diego -- \textit{Physics Colloquium, Cottrell Lecture} \when{2026} % May 7, 2026

    \item San Diego State University -- \textit{Physics Colloquium, Cottrell Lecture} \when{2026} % May 8, 2026

    \item Cornell University -- \textit{HEP Seminar} \when{2026} % April 24, 2026

    \item Fermilab -- \textit{Wine \& Cheese Seminar} \when{2026} % 16 Jan 2026

    \item Yale University -- \textit{NPA Seminar} \when{2024} % 9 May 2024

    \item New York University -- \textit{Physics Colloquium} \when{2024} % March 2024

    \item University of California, Irvine, Virtual -- \textit{AI Seminar} \when{2022}
    % July 12, 2022

    \item University of Kansas, Lawrence, KS -- \emph{Particle Physics on the Plains} \when{2022}
    % March 8, 2022

    \item University of Tennessee, Knoxville, TN -- \emph{HEP Seminar} \when{2021}%
    % April 21, 2021
    % \item \textbf{R-Parity Violating SUSY is Hard to Find}
    %   \begin{itemizetighter}
    %   \item University of Tennessee, \emph{Seminar}, Knoxville, TN \when{2021}%
    %   \end{itemizetighter}

    \item University of Tennessee, Knoxville, TN -- \emph{Physics Colloquium} \when{2020}%
    %  Nov 23, 2020
    % \item \textbf{What To Do When Nature Is Unnatural}
    %   \begin{itemizetighter}
    %   \item University of Tennessee, \emph{Physics Colloquium}, Knoxville, TN \when{2020} % 11 December 2018
    %   \end{itemizetighter}

    \item Fermilab, Batavia, IL -- \emph{LPC Physics Forum} \when{2020}
    % 19 Nov 2020
    % \item \textbf{Displaced Vertex Searches at ATLAS}
    %   \begin{itemizetighter}
    %   \item Fermilab, \emph{LPC Physics Forum}, Batavia, IL \when{2020} % 11 December 2018
    %   \end{itemizetighter}


    \item Lund University, Sweden -- \emph{Seminar} \when{2020} % Mar 5, 2020
    \item Michigan State University, East Lansing, MI -- \emph{Physics Colloquium} \when{2019} % 5 February 2019
    \item University of Illinois, Urbana-Champaign, IL -- \emph{Physics Colloquium} \when{2018} % 25 October 2018
    \item Fermilab, Batavia, IL --  \emph{Topic of the Week} \when{2018} % 22 October 2018
    \item SLAC, Palo Alto, CA -- \emph{Seminar} \when{2018} % 18 October 2018
    \item University of Oregon, Eugene, OR -- \emph{Seminar} \when{2018} % 15 October 2018
    \item University of Pennsylvania, Philadelphia, PA -- \emph{Seminar} \when{2018} % 7 August 2018
    \item Rutgers University, Piscataway, NJ -- \emph{Seminar} \when{2018} % 6 August 2018

    % \item \textbf{How SUSY Can Still Save The Day / The SUSY Swindle}
    %   \begin{itemizetighter}
    %   \item Lund University, Sweden \when{2020}
    %   \item Michigan State University, East Lansing, MI \when{2019} % 5 February 2019
    %   \item University of Illinois, Urbana-Champaign, IL \when{2018} % 25 October 2018
    %   \item Fermilab \emph{Topic of the Week}, Batavia, IL \when{2018} % 22 October 2018
    %   \item SLAC, Palo Alto, CA \when{2018} % 18 October 2018
    %   \item University of Oregon, Eugene, OR \when{2018} % 15 October 2018
    %   \item University of Pennsylvania, Philadelphia, PA \when{2018} % 7 August 2018
    %   \item Rutgers University, Piscataway, NJ \when{2018} % 6 August 2018
    %   \end{itemizetighter}


    \item Harvard University, Cambridge, MA -- \emph{LPPC Seminar} \when{2018} % 11 October 2018
    % \item \textbf{On the Higgs Branching Ratios}
    %   \begin{itemizetighter}
    %   \item Harvard University - LPPC Seminar, Cambridge, MA\when{2018} % 11 October 2018
    %   \end{itemizetighter}


    \item Oskar Klein Centre, Stockholm, Sweden -- \emph{HEP Seminar}\when{2018} % 22 May 2018
    \item Lawrence Berkeley National Laboratory, Berkeley, CA  -- \emph{RPM Seminar}\when{2017} % 17 August 2017
    % \item \textbf{Searches for Sneaky SUSY at the ATLAS Experiment \& A Toy Entropic Explanation for the Hierarchy Problem}
    %   \begin{itemizetighter}
    %   \item Oskar Klein Centre, Stockholm, Sweden \when{2018} % 22 May 2018
    %   \item Lawrence Berkeley National Laboratory, Berkeley, CA \when{2017} % 17 August 2017
    %   \end{itemizetighter}

    \item Harvard University, Cambridge, MA -- \emph{LPPC Seminar}\when{2015} % 30 September 2015
  % \item \textbf{Searches for R-Parity Violating SUSY at the ATLAS Experiment}
  %   \begin{itemizetighter}
  %   \item Harvard University - LPPC Seminar, Cambridge, MA\when{2015}% 30 September 2015
  %   \end{itemizetighter}

    \item University of Adelaide, Australia -- \emph{CSSM Seminar}\when{2014} % 13 August 2014
    \item University of Melbourne, Australia -- \emph{Seminar}\when{2014} % 10 July 2014
    \item Institut de F\'isica d'Altes Energies (IFAE), Barcelona, Spain -- \emph{Seminar}\when{2014} % 7 January 2014

  % \item \textbf{A Search for B-Violating Supersymmetry in Multijet Signatures and Light-quark vs. Gluon Jet Tagging}
  %   \begin{itemizetighter}
  %   \item University of Adelaide, Australia\when{2014}% 13 August 2014
  %   \item University of Melbourne, Australia\when{2014}% 10 July 2014
  %   \item Institut de F\'isica d'Altes Energies (IFAE), Barcelona, Spain\when{2014} % 7 January 2014
  %   \item Santa Cruz Institute of Particle Physics Seminar, Geneva, Switzerland\when{2013}% 12
  %   \end{itemizetighter}

    \item 2nd International Spring School on Particle Physics and Philosophy -- \emph{Seminar}\when{2014}
  % \item \textbf{...But What If I Like Naturalness?}
  %   \begin{itemizetighter}
  %   \item 2nd International Spring School on Particle Physics and Philosophy
  %   \item Wuppertal, Germany\when{2014}% 31 March 2014
  %   \end{itemizetighter}

    \item Santa Cruz Institute of Particle Physics Seminar, Geneva, Switzerland -- \emph{Seminar}\when{2013} % 12

  % \item \textbf{RPV Stops at the LHC}
  %   \begin{itemizetighter}
  %   \item LPC Workshop on Exotic Top Partners
  %   \item Fermi National Accelerator Laboratory, LHC Physics Center, Batavia, IL\when{2013}% 26 September 2013
  %   \end{itemizetighter}


% \end{itemize}
\end{itemizetighter}

  \emph{Conference Presentations}

  % \begin{itemize}
  \begin{itemizetighter}

  \item DPF, \textit{Energy Frontier Session}, Fermilab \when{2026} % July 23 2026
  \item DPF, \textit{Outreach \& Engagement Session}, Fermilab \when{2026} % July 21 2026

  \item AI4MuC -- \textit{Overview Talk} \when{2026} % June 29, 2026


  \item IDEC Conference, Columbia College of Chicago  \when{2026} % Mar 2026

  \item Good $\nu$s, University of Pittsburgh -- \emph{Invited}  \when{2025} % 29 Oct 2025

  \item Physics At The Highest Energies With Colliders, Galileo Galilei Institute, Florence, Italy -- \emph{Invited}  \when{2025} % 29 July 2025


  \item LHCP2025, NTU, Taipei, Taiwan -- \emph{Invited}  \when{2025} % 4 March 2025


  \item Aspen Center for Physics -- \emph{Invited Muon Collider Overview Plenary}  \when{2025} % 4 March 2025

  \item Inaugural US Muon Collider Meeting, Fermilab -- \emph{Invited Experimental Overview Plenary}  \when{2024} % 7 August 2024

  \item Workshop on Long-Lived Particles, University of Tokyo, Tokyo, Japan -- \emph{o.b.o. CMS}  \when{2024} % 4 July 2024

  \item DPF/Pheno, Pittsburgh, PA \when{2024} % May 2024

  \item SESAPS2023, EKU, Richmond, KY -- \emph{Invited `Highlight' Talk}\when{2023} % XX Nov 2023

  \item Fermilab Users Meeting -- \emph{Invited Plenary} \when{2023}
  % June 28, 2023

  \item UCSB, Kavli Institute of Theoretical Physics, Muon Collider Workshop -- \emph{Invited Plenary} \when{2023}
  % March 8, 2023

  \item CMS Exotica Workshop -- \emph{Invited Plenary} \when{2022} % 11 Dec 2022

  \item ML4Jets, Rutgers University, Piscataway, NJ \when{2022} % 1 Nov 2022

  \item Snowmass Community Summer Study Workshop, Seattle, WA \when{2022} % 21 July 2022

  \item APS April Meeting, New York City, NY \when{2022} % 10 April 2022

  \item CMS Exotica Workshop -- \emph{Invited Plenary} \when{2021} % 24 Nov 2021

  \item SESAPS2021, FSU, Tallahassee, Fl -- \emph{Invited `Highlight' Talk}\when{2021} % 19 Nov 2021


  % \item \textbf{ColliderScope: Reaching New Audiences with Electronic(s) Music}
  %   \begin{itemizetighter}
  %   \item 18th International Particle Physics Outreach Group meeting, CERN \when{2019} % 28 November 2019
  %   \item International Conference on New Frontiers in Physics, Kolymbari, Crete \when{2019} % 21 August 2019
  %   \item APS Division of Particles and Fields Meeting, Boston, MA \when{2019} % 30 July 2019
  %   \end{itemizetighter}

  \item LHCP2021, Virtual -- \emph{Invited Plenary, o.b.o. ATLAS and CMS}\when{2021}
  % \item 6/9/2021 Neutral FIPs at the LHC and Beyond

  \item APS April Meeting, Virtual -- \emph{Invited Plenary, o.b.o. ATLAS}\when{2021}
  % \item \textbf{New Long-Lived Particle Results from the LHC} -- \emph{Plenary Talk} \when{2021} % 4/17/2021
  %   \begin{itemizetighter}
  %   \item APS April Meeting, Virtual
  %   \end{itemizetighter}

  \item A Rainbow Of Dark Sectors, Aspen Center for Physics, Virtual -- \emph{Invited Plenary}\when{2021}
  % \item \textbf{The Lifetime of the Long-Lived Particle Era at the LHC} -- \emph{Plenary Talk} \when{2021} % 24 Mar 2021
  %   \begin{itemizetighter}
  %   \item A Rainbow Of Dark Sectors, Aspen Center for Physics, Virtual
  %   \end{itemizetighter}


  \item LHC Experiments Committee (LHCC) Meeting, CERN, Geneva -- \emph{Open Session} \when{2020}
    % Nov 18, 2020
    % \item \textbf{Updates from the ATLAS Experiment}
    %   \begin{itemizetighter}
    %   \item LHC Experiments Committee (LHCC) Meeting Open Session, CERN, Geneva \when{2020}
    %   \end{itemizetighter}


  \item ICNFP, Kolymbari, Crete, Greece -- \emph{o.b.o. ATLAS}\when{2019} % 21 August 2019
  % \item \textbf{Searches for Long-Lived Particles with the ATLAS Detector}
  %   \begin{itemizetighter}
  %   \item \emph{On behalf of the ATLAS Collaboration}
  %   \item International Conference on New Frontiers in Physics, Kolymbari, Crete \when{2019} % 21 August 2019
  %   \end{itemizetighter}


  \item Higgs Cross-Section Working Group Meeting, CERN, Geneva -- \emph{Invited Plenary} \when{2018} % 11 December 2018
  % \item \textbf{Long-Lived Particles and the Higgs}
  %   \begin{itemizetighter}
  %   \item Higgs Cross-Section Working Group Meeting, CERN, Geneva \when{2018} % 11 December 2018
  %   \end{itemizetighter}

    \item USATLAS Workshop, Pittsburgh, PA -- \emph{Invited Plenary} \when{2018} % 1 August 2018
    % \item \textbf{SUSY and Exotics Overview} -- ``\emph{Headliner}'' Talk
    %   \begin{itemizetighter}
    %   \item USATLAS Workshop, University of Pittsburgh, Pittsburgh, PA \when{2018} % 1 August 2018
    %   \end{itemizetighter}

    % \item \textbf{The SUSY Paradox \& A Toy Entropic Explanation for the Hierarchy Problem}
    %   \begin{itemizetighter}
    %   \item Oskar Klein Centre, Stockholm, Sweden \when{2018} % 22 May 2018
    %   \end{itemizetighter}

  \item SUSY2017, Mumbai, India -- \emph{o.b.o. ATLAS} \when{2017} % 11 December 2017
  % \item \textbf{Searches for squarks and gluinos in scenarios with R-parity violating sparticle decays, or long- lived sparticles with ATLAS}
  %   \begin{itemizetighter}
  %   \item \emph{On behalf of the ATLAS Collaboration}
  %    \item SUSY2017, Mumbai, India \when{2017} % 11 December 2017
  %   \end{itemizetighter}

  \item IHEP-T2E, Kuala Lumpur, Malaysia -- \emph{Plenary, o.b.o. ATLAS} \when{2017} % 27 February 2017
  % \item \textbf{SUSY Searches in the ATLAS Experiment} -- \emph{Plenary Talk}
  %   \begin{itemizetighter}
  %   \item \emph{On behalf of the ATLAS Collaboration}
  %   \item IHEP-T2E, Kuala Lumpur, Malaysia \when{2017} % 27 February 2017
  %   \end{itemizetighter}

  \item SUSY2016, Melbourne, Victoria, Australia\when{2016} % 3 July 2016
  % \item \textbf{The Recursive Jigsaw Reconstruction Technique}
  %   \begin{itemizetighter}
  %   \item SUSY2016, Melbourne, Victoria, Australia \when{2016} % 3 July 2016
  %   \end{itemizetighter}

  % \item \textbf{New Searches for Strongly-Produced SUSY at the ATLAS Experiment}
  %   \begin{itemizetighter}
  %   \item CoEPP Annual Workshop, Torquay, Victoria, Australia\when{2016} % 18 February 2016
  %   \end{itemizetighter}

  \item SUSY2015, Lake Tahoe, CA\when{2015} % 24 August 2015
  % \item \textbf{Applications of the Recursive Jigsaw Technique}
  %   \begin{itemizetighter}
  %   \item SUSY2015, Lake Tahoe, CA\when{2015}% 24 August 2015
  %   \end{itemizetighter}

  \item Kruger2014, Kruger Gate, South Africa -- \emph{o.b.o. ATLAS} \when{2014} % 6 December 2014
  % \item \textbf{SUSY Searches in the ATLAS Experiment} -- \emph{Plenary Talk}
  %   \begin{itemizetighter}
  %   \item \emph{On behalf of the ATLAS Collaboration}
  %   \item Kruger2014, Kruger Gate, South Africa\when{2014}% 6 December 2014
  %   \end{itemizetighter}


  \item Fermilab, LPC Workshop on Exotic Top Partners, Batavia, IL -- \emph{Invited Plenary}\when{2013} % 26 September 2013

  \item SUSY2012, Peking University, Beijing, China -- \emph{o.b.o. ATLAS} \when{2012} % 17 August 2012
  % \item \textbf{Search for supersymmetry in resonant production and R-parity violating signatures with the ATLAS detector}
  %   \begin{itemizetighter}
  %   \item \emph{On behalf of the ATLAS Collaboration}
  %   \item SUSY2012, Peking University, Beijing, China\when{2012}% 17 August 2012
  %   \end{itemizetighter}

% \end{itemize}
\end{itemizetighter}


  \emph{Conference and Workshop Organization}

\begin{itemizetighter}

  \item \textit{USMCC 2026}, Stanford University, \emph{Co-chair (w/ M. Peskin), Program Committee} \when{2026} % Dec 2026
  \item \textit{USMCC 2025}, University of Chicago, IL, \emph{Detector Design and Performance Session Organizer [of 3]} \when{2025} % Aug 2025
  \item \textit{LHCP 2025}, Taipei, Taiwan, \emph{Poster Committee Member} \when{2025} % May 2025
  \item \textit{CPAD 2024}, Knoxville, TN, \emph{Chair of Local Organizing Committee [of 5]} \when{2024} % Nov 2024
  \item \textit{LHCP 2024}, Northeastern University, Boston, MA, \emph{Poster Committee Member} \when{2024} % May 2024

  \item \textit{The Future of High Energy Physics}, Aspen Center for Physics, \emph{Conference Organizer [of 4]} \when{2024} % Mar 2024
  \item \textit{PITT PACC Muon Collider Physics Benchmark Workshop}, Pittsburgh, PA, \emph{Invited Closing Remarks} \when{2023} % Nov 2023
  \item \textit{ML4Jets}, Rutgers University, \emph{Session Chair} \when{2022} % Oct 2022
  \item \textit{APS April Meeting}, New York, NY, \emph{Session Chair} \when{2022} % April 10 2022
  \item \textit{ATLAS SUSY and Exotics Workshop}, Virtual \when{2020} % September 2020
  \item \textit{ATLAS Reaching New Heights Workshop}, CERN \when{2019} % December 2019
  \item \textit{ATLAS SUSY Workshop}, Lecce, Italy \when{2019} % September 2019
  \item \textit{APS Division of Particles and Fields Meeting}, Northeastern University \when{2019} % 29 July - 2 August 2019
  \item \textit{ATLAS SUSY Workshop}, Stockholm, Sweden \when{2018} % 23-25 May 2018
  \item \textit{ATLAS SUSY and Exotics Workshop}, Bucharest, Romania \when{2017} % 8-12 May 2017
  \item \textit{Recursive Jigsaw Workshop}, Harvard University \when{2015} % 21 September 2015
  \item \textit{Workshop on LHC Searches}, LBNL \when{2014} % 30 January 2014

  \end{itemizetighter}

% \clearpage


\section*{Teaching}

  \begin{itemizetighter}
    \item Physics of Cosmic Rays \emph{Seminar \& Design/Build Course w/ UT Architecture \& Design} (PHYS 494) \when{2025}
    \item Classical Mechanics I \textit{\&} II (PHYS 311/312) \when{2021--24}
    \item UTK HEP/Astro Seminar    \when{2022}
    \item Graduate Research Participation Seminar  \when{2021}
  \end{itemizetighter}


  \internalonly{

    \section*{Current Student \textit{\&} Postdoc Mentorship}

    \begin{itemizetighter}
      \item Olivia Clark [Undergrad] \when{2025--}
      \item Micah Hillman [Undergrad] \when{2025--}
      \item \rightnote{UTK EUReCA 2nd Place Presenter Award}

      \item Allie Dabney [Undergrad] \when{2025--}
      \item Shivam [Ph.D.] \when{2024--}
      \item Caroline Riggall [Ph.D.] \when{2024--}
      \item Cordney Nash [Undergrad] \when{2024--}
      \item \rightnote{UTK `Volunteer of Distinction'}

      \item Colby Thompson [Ph.D.] \when{2023--}
      \item Adam Vendrasco [Ph.D.] \when{2022--}
      \item John Lawless [Ph.D.] \when{2022--}
      \item \rightnote{Stelson Fellowship for Outstanding Beginning Research}

      \item Emery Nibigira [Postdoc] \when{2022--}
      \item \rightnote{Fermilab LPC Distinguished Researcher}

      \item \rightnote{Lindau Nobel Laureate Meeting Selected Participant}

    \end{itemizetighter}


    \section*{Student Thesis Committee Membership}

    \begin{itemizetighter}
      \item Colby Thompson [Ph.D.], \emph{chair} \when{2025--}
      \item Raghav Chari [Undergrad, College Honors] \when{2025}
      \item John Lawless [Ph.D.], \emph{chair} \when{2024--}
      \item Brandi Skipworth [Ph.D.] \when{2024--}
      \item Christal Martin [Ph.D.] \when{2024--}
      \item Ewa Glimos [Ph.D.] \when{2024--}
      \item Taylor Sussmane [Undergrad, Honors], \emph{chair} \when{2024}
      \item John Melhorn [Undergrad, Honors] \when{2024}
      \item Luke Carpenter [Undergrad, Honors] \when{2024}
      \item Rebecca Godri [Ph.D.] \when{2023--}
      \item Daniel Ally [Ph.D.] \when{2022--}
      \item Jesse Harris [Ph.D.] \when{2022--}
      \item Satyajit Roy [Ph.D.] \when{2021--24}
      \item Austin Schmier [Ph.D.] \when{2021--24}
      \item Himal Acharya [Ph.D.] \when{2021}
      \item External Ph.D. Dissertation Review for Yale University \when{2021}
      \item External M.S. Thesis Review for University of Adelaide \when{2021}
    \end{itemizetighter}

  }

  \section*{Past Student \textit{\&} Postdoc Mentorship}


  \textit{n.b. Square brackets denote initial position of mentorship relationship. Current position in italics.}

    \begin{itemizetighter}
      \item Dr. Alexander Tuna [Postdoc, UTK] \hfill \emph{Research Staff, UCSD}
      \item Dr. Ankush Kanuganti [Postdoc, UTK] \hfill \emph{Postdoc, BNL}
      \item Micah Hillman [Undergrad, UTK] \hfill \emph{Physics Postbacc, Princeton}
      \item Charles Bell [Undergrad, UTK] \hfill \emph{Physics PhD Student, UMichigan, NSF GRFP Awardee, Cover of Science,}
        \item \hfill \emph{UTK Talley Award for Outstanding Undergrad Research,}
        \item \hfill \emph{UTK Douglas V. Roseberry Distinguished Upper Class Major Award}
      \item Lacey Dishman [Undergrad, UTK] \hfill \emph{Physics PhD Student, Boston University}
      \item Taylor Sussmane [Undergrad, UTK] \hfill \emph{Physics PhD Student, University of Wisconsin, Madison,}
        \item \hfill \emph{Douglas V. Roseberry Distinguished Upper Class Major Award,}
        \item \hfill \emph{UTK Talley Award for Outstanding Undergrad Leadership}
      \item Everett Hagen [Undergrad, UTK]
      \item Jack Peltier [Undergrad, UTK]\hfill \emph{Goldwater Scholar}
      \item Puck Cravens [Undergrad, UTK]
      \item Tanner Wolfe [Undergrad, UTK]
      \item Alexander Sizemore [Undergrad, UTK]
      \item Matthew Durkee [Undergrad, UTK] \hfill \emph{Defense Contractor}
      \item Philipp Wagenknecht [Ph.D., UTK] \hfill \emph{HEP Lab Technician, Brown University}
      \item Dr. Tamas Vami [Ph.D., Johns Hopkins] \hfill \emph{Postdoc, UCSB}
      \item Prof. Karri Folan DiPetrillo [Ph.D., Harvard] \hfill \emph{Asst. Professor, University of Chicago}
      \item Prof. Jennifer Roloff [Ph.D., Harvard] \hfill \emph{Asst. Professor, Brown University}
      \item Dr. Anthony Badea [Ph.D., Harvard] \hfill \emph{Schmidt AI Fellow, University of Chicago}
      \item Dr. Anne Fortman [Ph.D., Harvard] \hfill \emph{Postdoc, LBNL}
      \item Dr. Aaron White [Postdoc, Harvard] \hfill \emph{}
      \item Xiaohe Jia [Ph.D., Harvard] \hfill \emph{Industry}
      \item Madeline Bernstein [Undergrad, Harvard] \hfill \emph{Physics PhD Student, UC Berkeley}
      \item Dr. Stefanie Morgenstern [Ph.D., Dresden/CERN] \hfill \emph{CERN Fellow}
      \item Dr. Russell Smith [Ph.D., Columbia] \hfill \emph{Squarepoint, Financial Sector}
      \item Prof. Emma Alexander [Undergrad, Yale] \hfill \emph{Asst. Professor, Northwestern University}
      \item Dr. Ariel Ekblaw [Undergrad, Yale] \hfill \emph{Founder of the MIT Space Exploration Initiative, CEO of Aurelia Institute}
    \end{itemizetighter}


\section*{Selected Outreach Activities}

  \begin{itemize}

    \item \textbf{\href{https://cosmicrayn.utk.edu}{Cosmic Rayn: Cosmic Ray Educational Experience}} w/ UTK Architecture \textit{\&} Design \when{2024--26}

    \item \textbf{Quantum Canvases: Physics, the Arts, and the Humanities} \when{2023}

      \begin{itemizetighter}
      \item Lead organizer of a series of events in downtown Knoxville on the intersection of physics, the arts, and the humanities in collaboration with UTK Physics \textit{\&} Astronomy, CERN, and the UT Humanities Center.
      \item Electronic music event (``Harmonic Motion: Physics x Electronic Music'')
      \item Events around a collection of science fiction short stories
      \item A panel discussion on the intersection of physics and science fiction
      \item A humanities-facing physics colloquium event
      \end{itemizetighter}


    \item \textbf{Higgs Tenth Anniversary Celebration} \when{2022}

      \begin{itemizetighter}
      \item Primary organizer for a documentary screening of \emph{Particle Fever} and open panel session for the tenth anniversary of the discovery of the Higgs Boson. Held at Central Cinema in Knoxville, TN.
      \end{itemizetighter}


    \item \textbf{STEAM Trailer Design Project} \when{2021--22}

      \begin{itemizetighter}
      \item Worked with the UTK College of Architecture \textit{\&} Design to design and construct an educational experience for local middle schoolers with a focus on the physics of sound and visualization.
      \end{itemizetighter}

    \item \textbf{Keynote presentation at event for Eureka Science Museum}, Grand Junction, CO \when{2021}


    \item \href{http://www.cern.ch/ColliderScope}{\textbf{ColliderScope -- Particle-Physics-Inspired Electronic Music Project}} \when{2019--}

      \begin{itemizetighter}
      \item Music with sound waves that show particle physics images in Lissajous figures
      \item Live performances:
      \item \hspace{2em} Luck Reunion (at Willie Nelson's ranch), Austin, TX (\emph{w/ Micah Nelson, Daniel Lanois}) \when{2026}
      \item \hspace{2em} Fly By Night, Knoxville, TN (\emph{w/ signal:noise; VJ Set}) \when{2026}
      \item \hspace{2em} Fly By Night, Knoxville, TN (\emph{w/ signal:noise; VJ Set}) \when{2025}
      \item \hspace{2em} Fly By Night, Knoxville, TN (\emph{w/ signal:noise; VJ Set}) \when{2025}
      \item \hspace{2em} Scruffy City Hall, Knoxville, TN (\emph{w/ Harmonic Motion}) \when{2025}
      \item \hspace{2em} Scruffy City Hall, Knoxville, TN (\emph{w/ Harmonic Motion}) \when{2024}
      \item \hspace{2em} Old City Performing Arts Center, Knoxville, TN (\emph{w/ Harmonic Motion}) \when{2023}
      \item \hspace{2em} Blue Moon Tavern, Seattle, WA (\emph{w/ Snowmass summer meeting}) \when{2022}
      \item \hspace{2em} Roskilde Festival, Denmark \when{2022}
      \item \hspace{2em} Cross Club, Prague, YouTube, Live from the CERN Control Centre (\emph{w/ ICHEP2020}) \when{2020}
      \item \hspace{2em} Lund, Sweden \when{2020}
      \item \hspace{2em} Aeronaut Brewing, Somerville, MA \when{2019}
      \item \hspace{2em} ATLAS Experiment, Geneva, Switzerland \when{2019}
      \item \hspace{2em} CERN Open Days, Geneva, Switzerland \when{2019}
      \item \hspace{2em} Pohoda Music Festival, Tren\v{c}\'{i}n, Slovakia \when{2019}
      \item Commissioned Compositions:

      \item \hspace{2em} Intro Videos for the CERN Science Gateway Educational Program \when{2023}

      \item \hspace{2em} Intro Videos for CERN Music Festival Program \when{2019}

      \item Presented at multiple international and national conferences
    %   \item 18th International Particle Physics Outreach Group meeting, CERN \when{2019 % 28 November 2019}
    %   \item International Conference on New Frontiers in Physics, Kolymbari, Crete \when{2019 % 21 August 2019}
    %   \item APS Division of Particles and Fields Meeting, Boston, MA \when{2019 % 30 July 2019}
      \end{itemizetighter}


    \item \textbf{Presentations to Grade School Students}
    \begin{itemizetighter}
    \item L\&N STEM Academy, Knoxville, TN \when{2023}
    \item Tel Aviv University -- Future Scientists CERN Tour \when{2017--19}
    \item Five presentations to NJ and NY schools \when{2011--16}
    \end{itemizetighter}

    \item \textbf{Harvard NSW Activities at CERN} - \emph{Video} (\href{https://youtu.be/VHrJAkNsMhM}{YouTube})
    \begin{itemizetighter}
    \item Fully produced video for recruitment \when{2018}
    \end{itemizetighter}

  \item \href{http://www.cern.ch/inparticular}{\textbf{In Particular}} - \emph{Podcast}\when{2015--16}

    \begin{itemizetighter}
    \item Video, audio, and original music production
    % \item Production of ``Video Shorts''
    % \item Web development, graphic design, and episode production
    \item Tens of thousands of downloads ($\sim11$~TB served)
    \item Featured by CERN, iTunes, The Guardian Science Blogs, Ars Technica, Vice, and others
    \end{itemizetighter}


  \item \textbf{Visiting Scholar, Dharma Drum Mountain, Taiwan}
    \begin{itemizetighter}
    \item Several weeks of conversations with monastery's monks and academics \when{2014}
    \end{itemizetighter}


  \end{itemize}


\section*{Press \textit{\&} Interviews}

      \begin{itemizetighter}

      \item \href{https://physics.utk.edu/tova-holmes-and-larry-lee-selected-as-fermilab-distinguished-researchers/}{UTK Physics -- Tova Holmes And Larry Lee Selected As Fermilab Distinguished Researchers}\when{2026}

      \item \href{https://news.utk.edu/2025/02/12/university-of-tennessee-physicist-named-cottrell-scholar/}{UTK News -- University of Tennessee Physicist Named Cottrell Scholar}\when{2025}

      \item \emph{Science} -- L.~Lee, C.~Bell, \textbf{\href{https://www.science.org/toc/science/383/6690}{Cover of Issue 383 6690}}  \when{2024}


      \begin{itemizetightrightpad}
        \item \emph{Muon collider simulation work and renderings made in Unreal Engine were featured on the cover and in the cover article in the issue from March 29, 2024.}
      \end{itemizetightrightpad}

      \item \href{https://engage.aps.org/dpb/resources/newsletters/24-25-newsletter/f10-article}{APS DPB Annual Newsletter -- The Muon Collider: A Modern Machine for a Modern World}\when{2024}
      \item \href{https://cerncourier.com/a/a-new-generation-a-new-vision/}{CERN Courier -- A New Generation, A New Vision} \when{2024}
      \item \href{https://physics.utk.edu/the-particle-detectorists/}{UTK Physics -- The (Particle) Detectorists} \when{2024}
      \item \href{https://physics.utk.edu/four-ut-physics-students-win-doe-scgsr-support/}{UTK Physics -- NSF GRFP Awardee} \when{2024}
      \item \href{https://physics.utk.edu/the-art-of-muon-collisions/}{UTK Physics -- The Art of Muon Collisions} \when{2024}
      \item Interview for article for MIT Tech Review \when{2023}
      \item Interview for article for Symmetry Magazine on CERN music festival program \when{2023}
      \item Interview with Sparks Magazine on Asian Americans in STEM\when{2023}
      \item \href{https://physics.utk.edu/in-my-own-words-undergraduate-taylor-sussmane/}{UTK Physics -- In My Own Words: Undergraduate Taylor Sussmane} \when{2023}
      \item \href{https://physics.utk.edu/challenging-the-standard-model-lawrence-lee-wins-nsf-career-award/}{UTK Physics -- Challenging The Standard Model: Lawrence Lee Wins NSF CAREER Award}\when{2023}
      \item \href{https://physics.utk.edu/the-art-of-science/}{UTK Physics -- The Art of Science} \when{2023}
      \item \href{https://physics.utk.edu/the-yeti-returns/}{UTK Physics -- The YETI Returns} \when{2023}
      \item Interview with Danish podcast TechTopia \when{2022}
      \item \textit{ColliderScope} featured in \href{https://www.symmetrymagazine.org/article/lhc-music-through-the-colliderscope}{Symmetry Magazine} \when{2019}
      \item Interview for article for Symmetry Magazine on naturalness \when{2019}
      \item CERN Courier -- ATLAS Pushes the Limits on Supersymmetry \when{2019}

      \item Interview for Swedish-language podcast \textit{Professor Magenta} about science and art \when{2016}
      \item Public screening of film Particle Fever w/ panel Q\textit{\&}A\when{2014}


      \end{itemizetighter}


\section*{University, College, and Departmental Committees and Other Service Roles}

  \begin{itemizetighter}

    \item UTK Academy for Global Engagement - Mentor \when{2026--27}

    \item CAS Climate, Culture, and Community Committee \when{2025--}

    \item Faculty Search Committee - Particle Theory \when{2025--26}
    \item Faculty Search Committee - Lecturer \when{2025--26}

    \item ORIED Launch Pad Panel on Prestigious Fellowships \& Awards \when{2025}
    \item University NSF CAREER Workshop Panel \when{2025}
    \item Faculty Search Committee - Nuclear Theory \when{2024--25}

    \item Coordinator of Department Web Services \when{2024--}

    \item 4 UTK Senior Thesis Committees \when{2024--}

    \item University NSF CAREER Workshop Panel \when{2024}

    \item Provost Junior Faculty Fellows (\emph{UTK}) \when{2023--25}

    \item Graduate Recruitment Committee  \when{2023--}

    \item SPS Faculty Advisor  \when{2023--}

    \item Assessment Committee  \when{2023--24}

    \item Graduate Advising Committee \when{2022--}

    \item Planning Committee \when{2022--24}

    \item Faculty Search Committee - Medium Energy Nuclear Experiment \when{2022--23}

    \item Qualifying Exam Committee \when{2022--23, 24--25}

    \externalonly{
      \item 9 UTK PhD Thesis Committees \when{2021--}

    }

    \item Community Connections Committee (\emph{Chair}) \when{2021--23}

    \item Department Head Search Committee  (\emph{UTK College of Arts \& Sciences}) \when{2021--22}
  \end{itemizetighter}


% \section*{Professional Activities}

%   \begin{itemize}
%   \item Member, American Association for the Advancement of Science, 2013--Present
%   \item Member, New York Academy of Sciences, 2012--Present
%   \item Member, American Physical Society, 2008--2009
%   \end{itemize}

% \section*{Advising}


%   \begin{itemizetight}

%   \item Harvard University:
%   \begin{itemizetighter}
%     \item Brendon Bullard (PhD Candidate) \hfill[\emph{w/ M. Morii}]
%     \item Jennifer Roloff (2019 PhD) \hfill[\emph{w/ J. Huth}]
%     \item Karri Folan DiPetrillo (2019 PhD) \hfill[\emph{w/ M. Franklin}]
%     \item Adriana Rotaru (2019 Summer Undergraduate)  \hfill[\emph{w/ J. Huth}]
%     \item Madeline Bernstein (2017 Summer Undergraduate)  \hfill[\emph{w/ M. Franklin}]
%   \end{itemizetighter}

%   % \item Uppsala University, Students:
%   % \begin{itemizetighter}
%   %   \item Andreas Gillgren (2018-2019 Visiting Engineering Physics Student)  \hfill[\emph{w/ J. Huth}]
%   % \end{itemizetighter}

%   \item Columbia University:
%   \begin{itemizetighter}
%     \item Russell Smith (2016 PhD)  \hfill[\emph{w/ E. Hughes}]
%   \end{itemizetighter}

%   \item University of Adelaide, Three PhD and Two Undergraduate Students
%   % \begin{itemizetighter}
%   %     \item Jason Oliver (2019 PhD) \hfill[\emph{w/ P. Jackson}]
%   %     \item Anum Qureshi (2019 PhD) \hfill[\emph{'' ''}]
%   %     \item Damir Duvnjak (2014 Honours, 2019 PhD) \hfill[\emph{'' ''}]
%   %     \item Jens Kr\"oger (2015 Visiting) \hfill[\emph{'' ''}]
%   %     \item Finley Borgas (2014 Honours) \hfill[\emph{'' ''}]
%   % \end{itemizetighter}

%   \item Yale University, Five Undergraduate Honors Students
%   % \begin{itemizetight}
%   %     \item Madeleine Barrow (2012)
%   %     \item Ari Brill (2012)
%   %     \item Emma Alexander (2011)
%   %     \item James Giammona (2011)
%   %     \item Alexandra Turrini (2011)
%   % \end{itemizetight}
%   \end{itemizetight}

% \section*{Teaching}

% % \begin{changemargin}{1.5em}{0.0in}

%   \begin{itemizetight}
%     \item Yale University
%     \begin{itemizetight}
%       \item \emph{Fundamentals of Physics}, lecturing assistant for Prof. P. Tipton \when{2011}
%       \item \emph{Physics for the Life Sciences}, lecturing assistant for Prof. S. Mochrie \when{2010}
%       \item \emph{Modern Physical Measurement}, assistant for Prof. J. Harris \when{2010}
%       \item \emph{General Physics Laboratory}, assistant for Prof. K. Baker \when{2009}
%     \end{itemizetight}
%   \end{itemizetight}
% \end{changemargin}

% \clearpage
% \section*{References}

%   \begin{itemize}
%   \item \textbf{Prof. Melissa Franklin} - {\it Mallinckrodt Professor of Physics} - Harvard University
%   \item \hspace{1.6em} \href{mailto:Franklin@Physics.Harvard.edu}{Franklin@Physics.Harvard.edu}
%   \item \textbf{Prof. Tobias Golling} - {\it Associate Professor} - University of Geneva
%   \item \hspace{1.6em} \href{mailto:Tobias.Golling@unige.ch}{Tobias.Golling@unige.ch}
%   \item \textbf{Dr. Andreas Hoecker} - {\it Senior Research Physicist, ATLAS Spokesperson-Elect} - CERN
%   \item \hspace{1.6em} \href{mailto:Andreas.Hoecker@cern.ch}{Andreas.Hoecker@cern.ch}
%   \item \textbf{Prof. John Huth} - {\it Donner Professor of Science} - Harvard University
%   \item \hspace{1.6em} \href{mailto:Huth@g.Harvard.edu}{Huth@g.Harvard.edu}

%     % \pagebreak

%     % \item {\textit{\hspace{-1.6em} [References continued...] }}

%   \item \textbf{Prof. Paul Jackson} - {\it Associate Professor} - University of Adelaide
%   \item \hspace{1.6em} \href{mailto:P.Jackson@adelaide.edu.au}{P.Jackson@adelaide.edu.au}
%   \item \textbf{Dr. Zachary Marshall} - {\it Senior Scientist} - Lawrence Berkeley National Lab
%   \item \hspace{1.6em} \href{mailto:Zach.Marshall@cern.ch}{Zach.Marshall@cern.ch}

%   % \vspace{1em}
%   % \item \textit{Additional References}
%   % \begin{itemize}
%   % \item \textbf{Prof. Paul Tipton} - {\it Chair of Physics Department} - Yale University
%   % \item \hspace{1.6em} \href{mailto:paul.tipton@yale.edu}{Paul.Tipton@Yale.edu}
%   % \end{itemize}

%   \end{itemize}

\cvfooter

\end{document}
"
%%
%% which is what build.sh and .github/workflows/build-cv.yml do.

\newif\ifCVinternal
\ifdefined\CVinternal \CVinternaltrue \else \CVinternalfalse \fi

% \internalonly{<content>} -- included only in the internal edition.
\newcommand{\internalonly}[1]{\ifCVinternal #1\fi}

% \externalonly{<content>} -- included only in the external edition.
\newcommand{\externalonly}[1]{\ifCVinternal\else #1\fi}

% The internal edition says so in its PDF title, so the two files stay
% distinguishable once they have been detached from their filenames.
\ifCVinternal
  \def\LLedition{\name: Curriculum Vitae (internal edition)}
\else
  \def\LLedition{\name: Curriculum Vitae}
\fi

\hypersetup{
  colorlinks  = false,
  pdfauthor   = {\name},
  pdftitle    = {\LLedition},
  pdfsubject  = {Curriculum Vitae},
  pdfkeywords = {},
  pdfpagemode = UseNone,
}


%% ===========================================================================
%% Page layout
%% ===========================================================================

\geometry{
  body       = {6.5in, 9.0in},
  left       = 1.0in,
  top        = 1.0in,
  headheight = 14pt,   % enough room for the running header; silences fancyhdr
}

\pagestyle{fancy}
\fancyhf{}
\renewcommand{\headrulewidth}{0pt}
\fancyhead[L]{\name}
\fancyhead[R]{\thepage}

\setlength{\parindent}{0em}

\titlespacing{\section}{0pt}{0pt}{0pt}
\sectionfont{\rmfamily\mdseries\Large}
\subsectionfont{\rmfamily\mdseries\itshape\large}


%% ===========================================================================
%% Lists
%% ===========================================================================
%% The CV is built almost entirely from bullet-free lists that differ only in
%% how tightly they are set.  itemize itself is redefined, rather than
%% configured through enumitem, so that a plain \begin{itemize} anywhere picks
%% up the CV style.
%%
%%   itemize               default spacing between entries
%%   itemizetight          no extra space between entries
%%   itemizetighter        entries set as close as the leading allows
%%   itemizetightrightpad  tight, and inset from the right margin, for the
%%                         commentary set underneath an entry
%%
%% \LLlist{<left>}{<right>}{<itemsep>}{<parskip>}{<parsep>}

\newcommand{\LLlist}[5]{%
  \list{}{%
    \setlength{\leftmargin}{#1}%
    \setlength{\rightmargin}{#2}%
    \setlength{\itemsep}{#3}%
    \setlength{\parskip}{#4}%
    \setlength{\parsep}{#5}%
  }%
}

\renewenvironment{itemize}%
  {\LLlist{1.5em}{0em}{0.25em}{0pt}{0.25em}}{\endlist}
\newenvironment{itemizetight}%
  {\LLlist{1.5em}{0em}{0em}{-1pt}{0em}}{\endlist}
\newenvironment{itemizetighter}%
  {\LLlist{1.5em}{0em}{0em}{-0.4em}{0em}}{\endlist}
\newenvironment{itemizetightrightpad}%
  {\LLlist{1.5em}{3em}{0em}{-1pt}{0em}}{\endlist}

% Hanging indent for a run-on entry: \parshape 2 3.2em \linewidth \listindent
% \dimexpr\linewidth-\listindent\relax sets the first line full width and
% indents every line after it to \listindent.
\newlength{\listindent}
\setlength{\listindent}{4.8em}

% \changemargin{<left>}{<right>} ... \endchangemargin -- indents a block of the
% document; used to nest the collaboration sub-lists under their headings.
\def\changemargin#1#2{\list{}{\rightmargin#2\leftmargin#1}\item[]}
\let\endchangemargin=\endlist


%% ===========================================================================
%% The right-hand date column
%% ===========================================================================
%% Nearly every entry ends with a date set against the right margin.  \when
%% supplies both the \hfill and the fixed-width box that lines those dates up
%% into one column: the box is left-aligned, so "2024" and "2022--24" begin at
%% the same horizontal position.
%%
%% A \makebox is only as wide as it is declared, so an over-long date used to
%% overflow it silently and print out in the right margin -- six of the dates in
%% this CV ("2021--2026", "2022--23, 24--25", ...) are wider than the column.
%% \when measures the date first and sets an over-wide one flush right instead,
%% which costs the shared left edge on those few entries but keeps every date
%% inside the text block.

\newlength{\datecolumn}
\setlength{\datecolumn}{3.5em}

%% The glue ahead of a date is \datesep plus infinite stretch.  The stretch is
%% what pushes the date to the right margin; the fixed \datesep is a floor, so
%% an entry that grows long enough to reach its date reports an overfull box
%% instead of quietly printing its text hard against the year.  Every entry
%% currently clears it, so the floor costs nothing today -- it is there to catch
%% the next entry that does not.  \whenflush is the fix when one turns up: it
%% gives back the empty right-hand part of the date box, about 1.6em for a
%% four-digit year.
\newlength{\datesep}
\setlength{\datesep}{0.4em}

%% \LLdateglue is the glue that carries the date to the right margin.  It reads
%% as one \hfill most of the time, but is written out so that an entry too long
%% to hold its date can break: TeX takes the \penalty50, leaves the text on the
%% line it filled, and the \hbox{} -- which survives the break where glue would
%% not -- carries the second \hfill onto the next line, so the date still lands
%% flush right instead of stranded against the left margin.  \unskip drops the
%% source space ahead of the date, so the gap is \datesep however the entry was
%% typed.
\newcommand{\LLdateglue}{%
  \unskip\nobreak\hspace{\datesep}\hfill\penalty50 \hbox{}\nobreak\hfill
}

\newlength{\LLdatewidth}
\newcommand{\when}[1]{%
  \LLdateglue
  \settowidth{\LLdatewidth}{#1}%
  \ifdim\LLdatewidth>\datecolumn #1\else\makebox[\datecolumn][l]{#1}\fi
}

% \whenflush{<date>} -- for an entry whose text runs so close to the right
% margin that the shared left edge of the date column will not fit.  Sets the
% date flush right, giving up the alignment to keep the gap.
\newcommand{\whenflush}[1]{\LLdateglue #1}

% \rightnote{<text>} -- a right-aligned italic annotation (an award, a current
% position) hung underneath an entry.  The trailing pad makes it end where a
% four-digit year in the date column ends, rather than at the margin itself.
\newcommand{\rightnote}[1]{%
  \settowidth{\LLdatewidth}{2024}%
  \hfill\emph{#1}\hspace{\dimexpr\datecolumn-\LLdatewidth\relax}%
}


%% ===========================================================================
%% Live publication metrics from INSPIRE-HEP
%% ===========================================================================
%% Citation counts go stale the moment they are typed, so they are not written
%% into the CV by hand.  tools/fetch_inspire.py queries the INSPIRE-HEP API and
%% writes LL_InspireData.tex, a flat list of declarations:
%%
%%     \inspiresetcites{2690093}{3}         % citations of one paper
%%     \inspiresetstat{citations}{207,000}  % an author-level summary figure
%%
%% which the CV body reads back through \inspirepub and \inspirestat.  The
%% values set just below are fallbacks for a build that has never run the
%% fetcher -- an offline checkout, or Overleaf -- so the document always
%% compiles with plausible numbers.  build.sh and CI refresh the data first.

% Declarations written by tools/fetch_inspire.py.  The keys live in \csname, so
% the LL@ namespace prefix needs no \makeatletter here.
\newcommand{\inspiresetcites}[2]{\expandafter\gdef\csname LL@cites@#1\endcsname{#2}}
\newcommand{\inspiresetstat}[2]{\expandafter\gdef\csname LL@stat@#1\endcsname{#2}}

% Fallbacks: the last hand-checked figures, overridden by LL_InspireData.tex.
\inspiresetstat{papers}{1400}
\inspiresetstat{citations}{202,000}
\inspiresetstat{hindex}{209}

% \inspirestat{<key>} -- an author-level figure: papers, citations or hindex.
% A missing key is reported loudly rather than left silently blank.
\newcommand{\inspirestat}[1]{%
  \ifcsname LL@stat@#1\endcsname\csname LL@stat@#1\endcsname\else\textbf{[?~#1]}\fi
}

% \inspirecites{<record>} -- "[n citations]" for a record whose count is known.
% Nothing is printed for an unknown count, or for one below \citesmin, so a
% brand-new paper is left unannotated rather than advertising that nobody has
% cited it yet.
%
% The annotation costs about a dozen characters on every entry, which is enough
% to push an entry that already filled its line onto a second one.  Raising
% \citesmin, or shortening the label, is the dial if that happens too often.
\newcommand{\citesmin}{1}
\newcommand{\LLprintcites}[1]{%
  \ifnum#1<\citesmin\relax\else
    \nobreakspace{\small[#1\nobreakspace\ifnum#1=1\relax citation\else citations\fi]}%
  \fi
}
\newcommand{\inspirecites}[1]{%
  \ifcsname LL@cites@#1\endcsname\LLprintcites{\csname LL@cites@#1\endcsname}\fi
}

% \inspirepub{<record>}{<title>} -- a publication title in bold, linked to its
% INSPIRE-HEP record and followed by that record's live citation count.
\newcommand{\inspirepub}[2]{%
  \href{https://inspirehep.net/literature/#1}{\textbf{#2}}\inspirecites{#1}%
}

% \pub{<url>}{<title>} -- the same, for an entry with no INSPIRE-HEP record:
% arXiv-only preprints, CDS and CERN notes, bare DOI links.
\newcommand{\pub}[2]{\href{#1}{\textbf{#2}}}

% The generated data, if it has been fetched.  \IfFileExists keeps a fresh
% checkout compiling on the fallbacks above.
\IfFileExists{LL_InspireData.tex}{%% Publication metrics fetched from the INSPIRE-HEP API.
%% GENERATED FILE -- do not edit; run tools/fetch_inspire.py instead.
%%
%% \inspiresetcites{<record>}{<count>}  per-paper citation count
%% \inspiresetstat{<key>}{<value>}      author-level summary figure
%% Both are defined in LL_Preamble.tex.

%% Author-level summary, as shown at the top of the Publications section.
%% Papers and citations are rounded down because the CV says "over".
\inspiresetstat{papers}{1400}    % exact: 1,467
\inspiresetstat{citations}{207,000}  % exact: 207,783
\inspiresetstat{hindex}{211}

%% Per-paper citation counts, keyed by INSPIRE record id.
\inspiresetcites{3159704}{0}
\inspiresetcites{3137624}{2}
\inspiresetcites{3119390}{0}
\inspiresetcites{3103093}{4}
\inspiresetcites{2962374}{9}
\inspiresetcites{2958460}{1}
\inspiresetcites{2874678}{26}
\inspiresetcites{2840007}{27}
\inspiresetcites{2752508}{22}
\inspiresetcites{2690093}{3}
\inspiresetcites{2644916}{13}
\inspiresetcites{2642414}{406}
\inspiresetcites{2628398}{83}
\inspiresetcites{2005570}{23}
\inspiresetcites{1856547}{33}
\inspiresetcites{1831504}{102}
\inspiresetcites{1788448}{119}
\inspiresetcites{1756729}{145}
\inspiresetcites{1710078}{319}
\inspiresetcites{1701002}{245}
\inspiresetcites{1630632}{260}
\inspiresetcites{1609773}{260}
\inspiresetcites{1498566}{50}
\inspiresetcites{1458270}{135}
\inspiresetcites{1345355}{100}
\inspiresetcites{1298030}{152}
\inspiresetcites{1191022}{108}
}{%
  \typeout{^^JLL_CV: LL_InspireData.tex not found; using fallback INSPIRE
figures. Run tools/fetch_inspire.py to refresh them.^^J}%
}


%% ===========================================================================
%% Title block
%% ===========================================================================

% \cvheader -- the name, contact details and ORCID that open every document.
\newcommand{\cvheader}{%
  \begin{minipage}[t]{0.4\textwidth}
    {\huge \name}
  \end{minipage}%
  \begin{minipage}[t]{0.6\textwidth}
    \begin{flushright}
      \href{mailto:\email}{\email}\\
      \phone\\
      \href{https://\homepage}{\homepage}\\
      \href{https://orcid.org/\orcid}{\orcid}
      % A postal address, if one is ever wanted here:
      %   CERN CH-1211\\
      %   B\^atiment 40-1-D11\\
      %   23 Gen\`{e}ve, Switzerland
    \end{flushright}
  \end{minipage}%
}

% \cvfooter -- the "last updated" line that closes every document.
\newcommand{\cvfooter}{%
  \medskip
  \begin{center}
    \begin{small}
      Last updated: \today
    \end{small}
  \end{center}%
}
, and a body, so page
%% geometry, fonts, list styles and the custom commands below are only ever
%% edited in one place.  A document that wants to differ from a default
%% overrides it after the \input.
%%
%% Template attribution: Jason R. Blevins - Curriculum Vitae
%% Copyright (C) 2004-2013 Jason R. Blevins


%% ===========================================================================
%% Packages
%% ===========================================================================
%% hyperref goes last: it patches internals of most other packages.

\usepackage[T1]{fontenc}
\usepackage{CormorantGaramond}    % body and heading face
\usepackage{microtype}            % subtler justification, fewer overfull lines
\usepackage{wasysym}              % $\RHD$, the "lead role" marker
\usepackage{graphicx}             % for the (currently disabled) portrait

\usepackage{geometry}             % page size and margins
\usepackage{fancyhdr}             % running headers
\usepackage{lastpage}             % "page n of m"
\usepackage{titlesec}             % section spacing
\usepackage{sectsty}              % section fonts
\usepackage{enumitem}             % list configuration

\usepackage[hidelinks]{hyperref}  % clickable links, printed without decoration


%% ===========================================================================
%% Identity
%% ===========================================================================

\def\name{Lawrence Lee, PhD}
\def\email{Lawrence.Lee.Jr@cern.ch}
\def\phone{+1 856 765 4622}
\def\homepage{cern.ch/larry}
\def\orcid{0000-0002-5590-335X}


%% ===========================================================================
%% Internal and external editions
%% ===========================================================================
%% The internal edition carries material that is not for general circulation:
%% the current student and postdoc roster, thesis committee membership,
%% individual refereeing assignments, and student travel awards.  The external
%% edition is the default and needs no flag; the internal one is selected by
%% defining \CVinternal before the document is read:
%%
%%     pdflatex -jobname=LL_CV_internal "\def\CVinternal{}%% Lawrence Lee -- Curriculum Vitae
%%
%% Two editions are produced from this one source:
%%
%%   LL_CV.pdf           external, for general circulation
%%   LL_CV_internal.pdf  internal; additionally carries the current student and
%%                       postdoc roster, thesis committee membership, individual
%%                       refereeing assignments, and student travel awards
%%
%% Build both with ./build.sh.  Everything shared with LL_Publications.tex --
%% packages, page layout, list styles, the \internalonly / \externalonly
%% switches, and the live INSPIRE-HEP citation commands -- lives in
%% LL_Preamble.tex.  See README.md for the rest.

\documentclass[11pt,letterpaper]{article}

%!TEX root = LL_CV.tex
%%
%% Shared preamble for every document in this repository.
%%
%%   LL_CV.tex            the full CV, in an external and an internal edition
%%   LL_Publications.tex  the publication list on its own
%%
%% Both are just \documentclass, %!TEX root = LL_CV.tex
%%
%% Shared preamble for every document in this repository.
%%
%%   LL_CV.tex            the full CV, in an external and an internal edition
%%   LL_Publications.tex  the publication list on its own
%%
%% Both are just \documentclass, \input{LL_Preamble.tex}, and a body, so page
%% geometry, fonts, list styles and the custom commands below are only ever
%% edited in one place.  A document that wants to differ from a default
%% overrides it after the \input.
%%
%% Template attribution: Jason R. Blevins - Curriculum Vitae
%% Copyright (C) 2004-2013 Jason R. Blevins


%% ===========================================================================
%% Packages
%% ===========================================================================
%% hyperref goes last: it patches internals of most other packages.

\usepackage[T1]{fontenc}
\usepackage{CormorantGaramond}    % body and heading face
\usepackage{microtype}            % subtler justification, fewer overfull lines
\usepackage{wasysym}              % $\RHD$, the "lead role" marker
\usepackage{graphicx}             % for the (currently disabled) portrait

\usepackage{geometry}             % page size and margins
\usepackage{fancyhdr}             % running headers
\usepackage{lastpage}             % "page n of m"
\usepackage{titlesec}             % section spacing
\usepackage{sectsty}              % section fonts
\usepackage{enumitem}             % list configuration

\usepackage[hidelinks]{hyperref}  % clickable links, printed without decoration


%% ===========================================================================
%% Identity
%% ===========================================================================

\def\name{Lawrence Lee, PhD}
\def\email{Lawrence.Lee.Jr@cern.ch}
\def\phone{+1 856 765 4622}
\def\homepage{cern.ch/larry}
\def\orcid{0000-0002-5590-335X}


%% ===========================================================================
%% Internal and external editions
%% ===========================================================================
%% The internal edition carries material that is not for general circulation:
%% the current student and postdoc roster, thesis committee membership,
%% individual refereeing assignments, and student travel awards.  The external
%% edition is the default and needs no flag; the internal one is selected by
%% defining \CVinternal before the document is read:
%%
%%     pdflatex -jobname=LL_CV_internal "\def\CVinternal{}\input{LL_CV.tex}"
%%
%% which is what build.sh and .github/workflows/build-cv.yml do.

\newif\ifCVinternal
\ifdefined\CVinternal \CVinternaltrue \else \CVinternalfalse \fi

% \internalonly{<content>} -- included only in the internal edition.
\newcommand{\internalonly}[1]{\ifCVinternal #1\fi}

% \externalonly{<content>} -- included only in the external edition.
\newcommand{\externalonly}[1]{\ifCVinternal\else #1\fi}

% The internal edition says so in its PDF title, so the two files stay
% distinguishable once they have been detached from their filenames.
\ifCVinternal
  \def\LLedition{\name: Curriculum Vitae (internal edition)}
\else
  \def\LLedition{\name: Curriculum Vitae}
\fi

\hypersetup{
  colorlinks  = false,
  pdfauthor   = {\name},
  pdftitle    = {\LLedition},
  pdfsubject  = {Curriculum Vitae},
  pdfkeywords = {},
  pdfpagemode = UseNone,
}


%% ===========================================================================
%% Page layout
%% ===========================================================================

\geometry{
  body       = {6.5in, 9.0in},
  left       = 1.0in,
  top        = 1.0in,
  headheight = 14pt,   % enough room for the running header; silences fancyhdr
}

\pagestyle{fancy}
\fancyhf{}
\renewcommand{\headrulewidth}{0pt}
\fancyhead[L]{\name}
\fancyhead[R]{\thepage}

\setlength{\parindent}{0em}

\titlespacing{\section}{0pt}{0pt}{0pt}
\sectionfont{\rmfamily\mdseries\Large}
\subsectionfont{\rmfamily\mdseries\itshape\large}


%% ===========================================================================
%% Lists
%% ===========================================================================
%% The CV is built almost entirely from bullet-free lists that differ only in
%% how tightly they are set.  itemize itself is redefined, rather than
%% configured through enumitem, so that a plain \begin{itemize} anywhere picks
%% up the CV style.
%%
%%   itemize               default spacing between entries
%%   itemizetight          no extra space between entries
%%   itemizetighter        entries set as close as the leading allows
%%   itemizetightrightpad  tight, and inset from the right margin, for the
%%                         commentary set underneath an entry
%%
%% \LLlist{<left>}{<right>}{<itemsep>}{<parskip>}{<parsep>}

\newcommand{\LLlist}[5]{%
  \list{}{%
    \setlength{\leftmargin}{#1}%
    \setlength{\rightmargin}{#2}%
    \setlength{\itemsep}{#3}%
    \setlength{\parskip}{#4}%
    \setlength{\parsep}{#5}%
  }%
}

\renewenvironment{itemize}%
  {\LLlist{1.5em}{0em}{0.25em}{0pt}{0.25em}}{\endlist}
\newenvironment{itemizetight}%
  {\LLlist{1.5em}{0em}{0em}{-1pt}{0em}}{\endlist}
\newenvironment{itemizetighter}%
  {\LLlist{1.5em}{0em}{0em}{-0.4em}{0em}}{\endlist}
\newenvironment{itemizetightrightpad}%
  {\LLlist{1.5em}{3em}{0em}{-1pt}{0em}}{\endlist}

% Hanging indent for a run-on entry: \parshape 2 3.2em \linewidth \listindent
% \dimexpr\linewidth-\listindent\relax sets the first line full width and
% indents every line after it to \listindent.
\newlength{\listindent}
\setlength{\listindent}{4.8em}

% \changemargin{<left>}{<right>} ... \endchangemargin -- indents a block of the
% document; used to nest the collaboration sub-lists under their headings.
\def\changemargin#1#2{\list{}{\rightmargin#2\leftmargin#1}\item[]}
\let\endchangemargin=\endlist


%% ===========================================================================
%% The right-hand date column
%% ===========================================================================
%% Nearly every entry ends with a date set against the right margin.  \when
%% supplies both the \hfill and the fixed-width box that lines those dates up
%% into one column: the box is left-aligned, so "2024" and "2022--24" begin at
%% the same horizontal position.
%%
%% A \makebox is only as wide as it is declared, so an over-long date used to
%% overflow it silently and print out in the right margin -- six of the dates in
%% this CV ("2021--2026", "2022--23, 24--25", ...) are wider than the column.
%% \when measures the date first and sets an over-wide one flush right instead,
%% which costs the shared left edge on those few entries but keeps every date
%% inside the text block.

\newlength{\datecolumn}
\setlength{\datecolumn}{3.5em}

%% The glue ahead of a date is \datesep plus infinite stretch.  The stretch is
%% what pushes the date to the right margin; the fixed \datesep is a floor, so
%% an entry that grows long enough to reach its date reports an overfull box
%% instead of quietly printing its text hard against the year.  Every entry
%% currently clears it, so the floor costs nothing today -- it is there to catch
%% the next entry that does not.  \whenflush is the fix when one turns up: it
%% gives back the empty right-hand part of the date box, about 1.6em for a
%% four-digit year.
\newlength{\datesep}
\setlength{\datesep}{0.4em}

%% \LLdateglue is the glue that carries the date to the right margin.  It reads
%% as one \hfill most of the time, but is written out so that an entry too long
%% to hold its date can break: TeX takes the \penalty50, leaves the text on the
%% line it filled, and the \hbox{} -- which survives the break where glue would
%% not -- carries the second \hfill onto the next line, so the date still lands
%% flush right instead of stranded against the left margin.  \unskip drops the
%% source space ahead of the date, so the gap is \datesep however the entry was
%% typed.
\newcommand{\LLdateglue}{%
  \unskip\nobreak\hspace{\datesep}\hfill\penalty50 \hbox{}\nobreak\hfill
}

\newlength{\LLdatewidth}
\newcommand{\when}[1]{%
  \LLdateglue
  \settowidth{\LLdatewidth}{#1}%
  \ifdim\LLdatewidth>\datecolumn #1\else\makebox[\datecolumn][l]{#1}\fi
}

% \whenflush{<date>} -- for an entry whose text runs so close to the right
% margin that the shared left edge of the date column will not fit.  Sets the
% date flush right, giving up the alignment to keep the gap.
\newcommand{\whenflush}[1]{\LLdateglue #1}

% \rightnote{<text>} -- a right-aligned italic annotation (an award, a current
% position) hung underneath an entry.  The trailing pad makes it end where a
% four-digit year in the date column ends, rather than at the margin itself.
\newcommand{\rightnote}[1]{%
  \settowidth{\LLdatewidth}{2024}%
  \hfill\emph{#1}\hspace{\dimexpr\datecolumn-\LLdatewidth\relax}%
}


%% ===========================================================================
%% Live publication metrics from INSPIRE-HEP
%% ===========================================================================
%% Citation counts go stale the moment they are typed, so they are not written
%% into the CV by hand.  tools/fetch_inspire.py queries the INSPIRE-HEP API and
%% writes LL_InspireData.tex, a flat list of declarations:
%%
%%     \inspiresetcites{2690093}{3}         % citations of one paper
%%     \inspiresetstat{citations}{207,000}  % an author-level summary figure
%%
%% which the CV body reads back through \inspirepub and \inspirestat.  The
%% values set just below are fallbacks for a build that has never run the
%% fetcher -- an offline checkout, or Overleaf -- so the document always
%% compiles with plausible numbers.  build.sh and CI refresh the data first.

% Declarations written by tools/fetch_inspire.py.  The keys live in \csname, so
% the LL@ namespace prefix needs no \makeatletter here.
\newcommand{\inspiresetcites}[2]{\expandafter\gdef\csname LL@cites@#1\endcsname{#2}}
\newcommand{\inspiresetstat}[2]{\expandafter\gdef\csname LL@stat@#1\endcsname{#2}}

% Fallbacks: the last hand-checked figures, overridden by LL_InspireData.tex.
\inspiresetstat{papers}{1400}
\inspiresetstat{citations}{202,000}
\inspiresetstat{hindex}{209}

% \inspirestat{<key>} -- an author-level figure: papers, citations or hindex.
% A missing key is reported loudly rather than left silently blank.
\newcommand{\inspirestat}[1]{%
  \ifcsname LL@stat@#1\endcsname\csname LL@stat@#1\endcsname\else\textbf{[?~#1]}\fi
}

% \inspirecites{<record>} -- "[n citations]" for a record whose count is known.
% Nothing is printed for an unknown count, or for one below \citesmin, so a
% brand-new paper is left unannotated rather than advertising that nobody has
% cited it yet.
%
% The annotation costs about a dozen characters on every entry, which is enough
% to push an entry that already filled its line onto a second one.  Raising
% \citesmin, or shortening the label, is the dial if that happens too often.
\newcommand{\citesmin}{1}
\newcommand{\LLprintcites}[1]{%
  \ifnum#1<\citesmin\relax\else
    \nobreakspace{\small[#1\nobreakspace\ifnum#1=1\relax citation\else citations\fi]}%
  \fi
}
\newcommand{\inspirecites}[1]{%
  \ifcsname LL@cites@#1\endcsname\LLprintcites{\csname LL@cites@#1\endcsname}\fi
}

% \inspirepub{<record>}{<title>} -- a publication title in bold, linked to its
% INSPIRE-HEP record and followed by that record's live citation count.
\newcommand{\inspirepub}[2]{%
  \href{https://inspirehep.net/literature/#1}{\textbf{#2}}\inspirecites{#1}%
}

% \pub{<url>}{<title>} -- the same, for an entry with no INSPIRE-HEP record:
% arXiv-only preprints, CDS and CERN notes, bare DOI links.
\newcommand{\pub}[2]{\href{#1}{\textbf{#2}}}

% The generated data, if it has been fetched.  \IfFileExists keeps a fresh
% checkout compiling on the fallbacks above.
\IfFileExists{LL_InspireData.tex}{\input{LL_InspireData.tex}}{%
  \typeout{^^JLL_CV: LL_InspireData.tex not found; using fallback INSPIRE
figures. Run tools/fetch_inspire.py to refresh them.^^J}%
}


%% ===========================================================================
%% Title block
%% ===========================================================================

% \cvheader -- the name, contact details and ORCID that open every document.
\newcommand{\cvheader}{%
  \begin{minipage}[t]{0.4\textwidth}
    {\huge \name}
  \end{minipage}%
  \begin{minipage}[t]{0.6\textwidth}
    \begin{flushright}
      \href{mailto:\email}{\email}\\
      \phone\\
      \href{https://\homepage}{\homepage}\\
      \href{https://orcid.org/\orcid}{\orcid}
      % A postal address, if one is ever wanted here:
      %   CERN CH-1211\\
      %   B\^atiment 40-1-D11\\
      %   23 Gen\`{e}ve, Switzerland
    \end{flushright}
  \end{minipage}%
}

% \cvfooter -- the "last updated" line that closes every document.
\newcommand{\cvfooter}{%
  \medskip
  \begin{center}
    \begin{small}
      Last updated: \today
    \end{small}
  \end{center}%
}
, and a body, so page
%% geometry, fonts, list styles and the custom commands below are only ever
%% edited in one place.  A document that wants to differ from a default
%% overrides it after the \input.
%%
%% Template attribution: Jason R. Blevins - Curriculum Vitae
%% Copyright (C) 2004-2013 Jason R. Blevins


%% ===========================================================================
%% Packages
%% ===========================================================================
%% hyperref goes last: it patches internals of most other packages.

\usepackage[T1]{fontenc}
\usepackage{CormorantGaramond}    % body and heading face
\usepackage{microtype}            % subtler justification, fewer overfull lines
\usepackage{wasysym}              % $\RHD$, the "lead role" marker
\usepackage{graphicx}             % for the (currently disabled) portrait

\usepackage{geometry}             % page size and margins
\usepackage{fancyhdr}             % running headers
\usepackage{lastpage}             % "page n of m"
\usepackage{titlesec}             % section spacing
\usepackage{sectsty}              % section fonts
\usepackage{enumitem}             % list configuration

\usepackage[hidelinks]{hyperref}  % clickable links, printed without decoration


%% ===========================================================================
%% Identity
%% ===========================================================================

\def\name{Lawrence Lee, PhD}
\def\email{Lawrence.Lee.Jr@cern.ch}
\def\phone{+1 856 765 4622}
\def\homepage{cern.ch/larry}
\def\orcid{0000-0002-5590-335X}


%% ===========================================================================
%% Internal and external editions
%% ===========================================================================
%% The internal edition carries material that is not for general circulation:
%% the current student and postdoc roster, thesis committee membership,
%% individual refereeing assignments, and student travel awards.  The external
%% edition is the default and needs no flag; the internal one is selected by
%% defining \CVinternal before the document is read:
%%
%%     pdflatex -jobname=LL_CV_internal "\def\CVinternal{}%% Lawrence Lee -- Curriculum Vitae
%%
%% Two editions are produced from this one source:
%%
%%   LL_CV.pdf           external, for general circulation
%%   LL_CV_internal.pdf  internal; additionally carries the current student and
%%                       postdoc roster, thesis committee membership, individual
%%                       refereeing assignments, and student travel awards
%%
%% Build both with ./build.sh.  Everything shared with LL_Publications.tex --
%% packages, page layout, list styles, the \internalonly / \externalonly
%% switches, and the live INSPIRE-HEP citation commands -- lives in
%% LL_Preamble.tex.  See README.md for the rest.

\documentclass[11pt,letterpaper]{article}

\input{LL_Preamble.tex}


\begin{document}
\thispagestyle{empty} % no running header on the first page

\cvheader

% Alternative title blocks, kept for reference:
%   \centerline{\huge \bf \name}          % name centred and bold
%   \begin{picture}(0,0)                  % portrait in the top right corner
%     \put(380,-30){\hbox{\includegraphics[height=12em]{Portrait.jpg}}}
%   \end{picture}

\bigskip

% Spacer that holds the vertical gap ahead of the first section.  It pairs with
% the summary paragraph below, which is currently disabled.
\begin{minipage}[t]{1.3em}
  \hspace{1em}
\end{minipage}
% \begin{minipage}[t]{0.75\textwidth}
% \emph{Broadly-trained particle physicist with expertise in searching for physics beyond the standard model at particle colliders. Particular expertise in supersymmetry, long-lived particles, and unconventional searches. Technical skills include expertise in detector data acquisition system, analysis software, and detector simulation.}
% \end{minipage}


\section*{Professional Experience}

  \begin{itemizetight}


    \item Associate Professor, University of Tennessee, Knoxville \when{2026--}
    \item Assistant Professor, University of Tennessee, Knoxville \when{2021--26}

    % \begin{itemizetighter}
    %   \item {University of Tennessee}, Knoxville, TN
    % \end{itemizetighter}

    \item Research Associate, Harvard University - \emph{Supervisor: J. Huth} \when{2016--21}

    % \begin{itemizetighter}
    %   \item {Harvard University}, Cambridge, MA -- \emph{Supervisor: J. Huth}
    % \end{itemizetighter}

    \item Research Associate, University of Adelaide - \emph{Supervisor: P. Jackson} \when{2014--16}

    % \begin{itemizetighter}
    %   \item {The University of Adelaide}, Adelaide, Australia -- \emph{Supervisor: P. Jackson}
    % \end{itemizetighter}

  \end{itemizetight}


\section*{Education}

  \begin{itemizetight}
    \item Ph.D., M.Phil., M.S., Physics, {Yale University} - \emph{Supervisor: T. Golling} \when{2014}
    % \begin{itemizetighter}
    %       \item \emph{Supervisor: T. Golling}
    % \end{itemizetighter}
    % \item M.S., M.Phil. Physics, {Yale University} \when{2012}
    \item B.S. Physics, {Rutgers University} - \emph{Supervisors: R. Ransome, R. Gilman, R. Tumulka} \when{2009}

  %   \begin{itemizetighter}
  %         \item \emph{Supervisors: R. Ransome, R. Gilman, R. Tumulka}
  %   \end{itemizetighter}
  \end{itemizetight}


% \section*{Fields of Research Interest}

%     \begin{itemizetightrightpad}
%     \item High Energy Collider Physics, Experimental searches for physics Beyond the Standard Model, Supersymmetry, Long-Lived Particles, Foundations of Quantum Mechanics
%     % \item Quantum Information, Philosophical Foundations of Quantum Mechanics
%     \end{itemizetightrightpad}


\section*{Awards \textit{\&} Honors}

\emph{External}

  \begin{itemizetighter}
    \item Cottrell Lectureship \when{2026}
    \item Fermilab LHC Physics Center Senior Distinguished Researcher \when{2026}
    \item Breakthrough Prize in Fundamental Physics Laureate  \when{2025}
    \item \hspace{2em}\emph{as part of the ATLAS and CMS Collaborations}
    \item Cottrell Scholar - Research Corporation for Science Advancement \when{2025}
    \item NSF CAREER Award \when{2023}
    \item European Physical Society's High Energy and Particle Physics Prize \when{2013}
    \item \hspace{2em}\emph{as part of the ATLAS Collaboration}
  \end{itemizetighter}


\emph{University of Tennessee, Knoxville}

\begin{itemizetighter}
  \item University of Tennessee Alumni Association -  Alumni Outstanding Teacher Award \when{2025}
  \item UTK Division of Student Success - Faculty Research Mentor Award \when{2025}
  \item UTK Provost - Junior Faculty Fellow \when{2023--25}
  \item UTK Physics - Research Advisor of the Year \when{2023--24}
  \item UTK College of Arts and Sciences - Faculty Academic Outreach Teaching Award \when{2022}
  \item UTK Physics - Teacher of the Year \when{2021--22}
\end{itemizetighter}


\section*{Funding}

\emph{Multi-Year Grants}

  \begin{itemizetighter}
    \item BNL Subcontract on Accelerator Physics - \$190k \when{2025--29}
    \item Cottrell Scholar - Research Corporation for Science Advancement - \$120k \when{2025--28}
    \item DOE Award for UTK HEP (4 PIs) - \$1.425M \when{2024--27}
    \item NSF CAREER Award - \$1M \when{2023--28}
    \item USCMS Upgrade Project Funding - \$73k \when{2022--24}
    \item Single-PI DOE Award - \$250k \when{2022--24}
  \end{itemizetighter}


  \emph{Other Funding}

  \begin{itemizetighter}
    \item Fermilab LHC Physics Center - \$56k\when{2026}
    \item Universities Research Association - \$7k\when{2026}
    \item Fermilab LHC Physics Center (Postdoc Emery Nibigira Led) - \$43k\when{2024}
    \item Universities Research Association (Postdoc Emery Nibigira Led) - \$20k\when{2024}
    \item UT-Battelle (ORNL Management Contractor) - \$74k \when{2024}
    \internalonly{
      \item UTK Graduate Student Senate Travel Award - John Lawless - \$1k\when{2024}
      \item SESAPS Travel Grant - Caroline Riggall - \$325 \when{2024}
      \item SESAPS Travel Grant - Charles Bell - \$500 \when{2023}
      \item SESAPS Travel Grant - Lacey Dishman - \$500 \when{2023}
      \item UTK Faculty Research Assistants Funding - Tanner Wolfe - \$2.3k \when{2023}
      \item FNAL LPC Guest \textit{\&} Visitor Funding - Emery Nibigira - \$5.2k \when{2023}
      \item FNAL LPC Guest \textit{\&} Visitor Funding - John Lawless - \$3.9k \when{2023}
      \item UTK Faculty Research Assistants Funding - Taylor Sussmane - \$2.3k \when{2022}
      \item SESAPS Travel Grant - Taylor Sussmane - \$500 \when{2022}
      \item SESAPS Travel Grant - Alexander Sizemore - \$500 \when{2022}
      \item FNAL LPC Guest \textit{\&} Visitor Funding - Ramon Ogaz - \$2.8k \when{2022}
      \item FNAL LPC Guest \textit{\&} Visitor Funding - Philipp Wagenknecht - \$5.2k \when{2022}
      \item UTK Physics Summer Research Fellowship - Matthew Durkee - \$6k \when{2022}
    }


    \item George Southgate Fellowship, The University of Adelaide - 5k AUD \when{2015}
  \end{itemizetighter}


\section*{Research Experience}

  \begin{itemize}


  \item \textsc{Member of the CMS Collaboration} \when{2021--}


  \begin{changemargin}{0in}{0.0in}
    \begin{itemize}

    \item \textsc{Physics Analysis}

    \vspace{0.1em}
    \begin{itemizetighter}

        \item \textit{Co-convener} LHCC Long-Lived Particle Working Group \when{2022--25}

        \item Heavy Stable Charged Particle Search \when{2022--}

        % \begin{itemizetightrightpad}
        % \item \it Searching for detector-stable particles carrying electric charge and leaving anomalous ionization and timing signatures in the CMS detector.
        % \end{itemizetightrightpad}


        \item Multijet Search \when{2023--}

        % \begin{itemizetightrightpad}
        % \item \it Searching for new physics in multijet final states using conventional techniques and machine learning
        % \end{itemizetightrightpad}

        \item Recursive Jigsaw Reconstruction Search \when{2022}

        % \begin{itemizetightrightpad}
        % \item \it Contributed to calculation of exclusion limits for a search for compressed supersymmetry spectra using the Recursive Jigsaw Reconstruction method.
        % \end{itemizetightrightpad}

    \end{itemizetighter}


    \item \textsc{Phase-II Outer Tracker Upgrade}

    \vspace{0.1em}
    \begin{itemizetighter}

      \item \textit{US CMS Level-3 Manager} for OT Electronics \when{2024--}
      \item \textit{US CMS Control Account Manager} for OT DAQ \when{2024--}

      \item Phase-II Tracker DAQ Software \when{2021--}

      \item Outer Tracker Module Assembly and Grading \when{2022--}

    \end{itemizetighter}


    \end{itemize}
  \end{changemargin}


  \item \textsc{Member of the ATLAS Collaboration} \when{2009--21}

  \begin{changemargin}{0in}{0.0in}
    \begin{itemize}
      % \parshape 2 3.2em \linewidth \listindent \dimexpr\linewidth-\listindent\relax

    % \item { \small [{\it Nota Bene:} Labels in \textsc{small caps} denote official ATLAS designations.] }


    % SUSY Software Review, Contours
    % Maintainer of big framework

    % NSW DAQ

    % \vspace{.5em}

    \item \textsc{Physics Analysis}

    \vspace{0.1em}
    \begin{itemizetighter}

        \item Common BSM Result Interpretation Framework \when{2018--21}

        % \begin{itemizetightrightpad}
        % \item \it Responsible for standards of exclusion interpretation and limit interpolation within the ATLAS SUSY working group, implementing new methods for presentation of results
        % \end{itemizetightrightpad}

        \item \textit{Co-convener} of SUSY RPV/LL sub-group \when{2017--19}

        % \begin{itemizetightrightpad}
        % \item \it Defined standards and direction of searches for long-lived particles and $R$-parity-violating supersymmetry. Served two terms.
        % \end{itemizetightrightpad}

        \item Searches for Long-Lived Particles \when{2016--21}

        % \begin{itemizetightrightpad}
        % \item \it Led many analysis teams in searches for long-lived particles in a variety of signatures. See Publications.
        % \end{itemizetightrightpad}
        % \pagebreak

        \item Common Analysis Software Development \when{2015--21}

        % \begin{itemizetightrightpad}
        % \item \it Primary developer of common analysis software used by roughly ten analysis teams
        % \end{itemizetightrightpad}


        \item Inclusive searches for Supersymmetry \when{2014--16}

        % \begin{itemizetightrightpad}
        % \item \it Led flagship search for supersymmetry in inclusive zero-lepton final states in 2016 including leading the first searches for new physics using ``Recursive Jigsaw'' variables. Created and commissioned novel triggers for use during Run-2 of the LHC geared toward finding new physics in low-mass-splitting scenarios.
        % \end{itemizetightrightpad}


        \item All-Hadronic BSM Signatures \when{2011--14}

        % \begin{itemizetightrightpad}
        % \item \it Led two searches for new phenomena in the context of $R$-Parity violating supersymmetry in fully hadronic final states
        % \end{itemizetightrightpad}


        \item Quark-vs-Gluon Jet Tagging \when{2010--12}

        % \begin{itemizetightrightpad}
        % \item \it Developed the first quark- vs. gluon-jet discriminator at a hadron collider
        % \end{itemizetightrightpad}

    \end{itemizetighter}

    \item \textsc{New Small Wheel Muon Spectrometer Upgrade}

    \vspace{0.1em}
    \begin{itemizetighter}

      \item \textit{Coordinator} Micromegas Trigger \when{2020}

      % \begin{itemizetightrightpad}
      % \item \it Integration of trigger electronics for the micromegas detector of the New Small Wheel (NSW) upgrade for the ATLAS muon spectrometer
      % \end{itemizetightrightpad}

      \item Online Software Coordination \when{2019--20} % \hspace{0.9em}}

      % \begin{itemizetightrightpad}
      % \item \it Designed, implemented, and/or maintain the readout, configuration, and calibration software systems for the NSW
      % \end{itemizetightrightpad}

      % \item Micromegas Digitization \when{2017--18}

      % \begin{itemizetightrightpad}
      % \item \it Responsible for digitization and simulation of the readout and trigger paths for the micromegas detector for the NSW
      % \end{itemizetightrightpad}

    \end{itemizetighter}


    \end{itemize}
  \end{changemargin}

  \item \textsc{Muon Collider Projects} \when{2020--}

  \begin{changemargin}{0in}{0.0in}
    \begin{itemizetight}

      \item \emph{USMCC Deputy Communication Coordinator} \when{2025--}
      \item 10 TeV MAIA Detector Concept Characterization \when{2023--}
      \item \emph{Taskforce Member} IMCC Software Taskforce \when{2024}

      \item Automating Professional 3D Rendering for Collider Event Displays \when{2023--}

      \item Simulation Software Workflow Development \when{2020--22}

    \end{itemizetight}
  \end{changemargin}


  \item \textsc{Accelerator Physics Projects} \when{2024--}

  \begin{changemargin}{0in}{0.0in}
    \begin{itemizetight}

    \item Helical FOFO Design for Muon Cooling for a Future Muon Collider \when{2024--}

    \item Using Self-Consistent Beam Distributions to Reduce Intra-Beam Stripping \when{2024--}

    \end{itemizetight}
  \end{changemargin}


  \item \textsc{External Collider Physics}


    \begin{itemizetight}

    \item Experimental Impact of Jet Fragmentation Reference Frames At Particle Colliders \when{2022--}


    \item Future collider projections for photon-induced production of charged BSM particles \when{2018--21}

    \item Bell Inequalities at Future $ee$ Colliders \when{2013--}

    \end{itemizetight}


  \item \textsc{Nucleon Physics Research} \when{2007--09}

  % \item Developed the prototype trigger and data acquisition systems for Fermilab E-906/SeaQuest \when{2008--2009}

  \begin{changemargin}{0in}{0.0in}
    \begin{itemizetightrightpad}
            \parshape 2 3.2em \linewidth \listindent \dimexpr\linewidth-\listindent\relax

    \item Developed the prototype trigger and data acquisition systems for Fermilab E-906/SeaQuest\vspace{0.2em}
    \item Assembled the photomultiplier tube units used in MINERvA\vspace{0.2em}
    \item Upgraded a Fabry-Perot cavity for use in a JLAB Compton polarimeter\vspace{0.2em}
    \end{itemizetightrightpad}
  \end{changemargin}

\end{itemize}


% \pagebreak

\section*{Publications}

  \begin{itemizetight}

  \input{LL_PubInclude.tex}

  \subsection*{One CMS Analysis Review Committee (ARC)}

  \subsection*{Nine ATLAS Editorial Boards (EBs)}

  % \subsection*{Journal Referee for Physical Review Letters, Physical Review D, Letters in High Energy Physics}

  % \begin{itemize}

  % \item Served on 8 ATLAS Editorial Boards

  % % \item {\it\href{https://atlas.web.cern.ch/Atlas/GROUPS/PHYSICS/CONFNOTES/ATLAS-CONF-2019-027/}{Soft $b$-hadron tagging for compressed SUSY scenarios}},~ATLAS-CONF-2019-027 (2019)

  % % \item {\it Search for new phenomena in a lepton plus high jet multiplicity final state with the ATLAS experiment using $\sqrt{s}$ = 13~TeV proton-proton collision data},~ATLAS-CONF-2016-094, ATLAS-CONF-2017-013, JHEP 09 (2017) 88

  % % \item {\it A search for pair-produced resonances in four jets final states at $\sqrt{s}=13$~TeV with the ATLAS detector},~ATLAS-CONF-2016-022 (2016)

  % % \item {\it\href{https://cds.cern.ch/record/2152392}{A search for R-parity violating decays of the top squark in four jet final states with the ATLAS detector at $\sqrt{s}=13$~TeV}},~ATLAS-CONF-2016-022 (2016)

  % % \item {\it\href{http://link.springer.com/article/10.1140/epjc/s10052-016-4065-1}{A new method to distinguish hadronically decaying boosted Z bosons from W bosons using the ATLAS detector}},~Eur. Phys. J. C \textbf{76} 238 (2016)

  % % \item {\it\href{https://cds.cern.ch/record/2017303}{Constraints on promptly decaying supersymmetric particles with lepton-number- and R-parity-violating interactions using Run-1 ATLAS data}},~ATLAS-CONF-2015-018 (2015)

  % \end{itemize}


  % \subsection*{Other Unpublished Work}

  % \begin{itemize}

  % % FQM Stuff
  % \item $\RHD$ GRW Theory and the Thermodynamic Arrow of Time (2010)
  % \item $\RHD$ Relativistic Considerations of the GRW Theory \emph{w/ R. Tumulka} (2008)
  % \item $\RHD$ Trouble with Locality: Adding Nonlocal Correlations \emph{w/ R. Tumulka} (2008)
  % \item $\RHD$ GRW Theory \emph{w/ R. Tumulka} (2008)

  % \end{itemize}


\end{itemizetight}

% \newpage

% \section*{Grants, Fellowships, \& Awards}

%   \begin{itemize}
%   \item George Southgate Fellowship, The University of Adelaide, 2015
%   \item 2nd International Spring School on Particle Physics and Philosophy Travel Award, 2014
%   \item AAAS/Science Program for Excellence in Science, 2013
%   % \item New York Academy of Science Sponsored Membership, 2012
%   \item Yale University Graduate Student Assembly Conference Travel Fund, 2012
%   % \item NSF Graduate Research Fellowship, Honorable Mention, 2011, 2010
%   \item Division of Nuclear Physics, American Physical Society Fall Meeting, Conference Travel Fund, 2008
%   % \item Summer Research Fellowship, Jefferson Laboratory, Rutgers University, 2008
%   \end{itemize}


\section*{Professional Activities}


  \begin{itemizetighter}
    \internalonly{
      \item \emph{Journal Referee:} PRD \when{2025}
    }
    \item RCSA Reviewer \when{2026} % For Cottrell
    \item Lead organizer, $\mu$CATALYST summer accelerator program \when{2026}
    \item Chair of Fermilab \emph{Colliders of Tomorrow} Committee \when{2025--}
    \item US LHC Users Association (USLUA) Executive Committee Member \when{2022--25}
    \item NSF Reviewer \when{2025} % For Kathy
    \item NSF Review Panel \when{2025}
    \item DOE Reviewer \when{2025}
    \item USCMS Phase-II Upgrade Level-3 Manager for Outer Tracker DAQ \when{2024--}
    \internalonly{
      \item \emph{Journal Referee:} PRD \when{2024}
      \item \emph{Journal Referee:} EPJC \when{2024}
    }

    \item Swiss National Science Foundation Reviewer \when{2024}
    \item International Muon Collider Collaboration Software Task Force \when{2024}
    \item US Muon Collider Collaboration Membership Coordinator \when{2024}
    \item US CMS (DOE+NSF) Project Review Panel \when{2024}
    \item NSF Reviewer ($\times$3 occasions) \when{2024}
    \item NSF Review Panel \when{2024}
    \internalonly{
      \item \emph{Journal Referee:} PRD \when{2024}
    }

    \item NSF Reviewer \when{2023}
    \item DOE Reviewer \when{2023}
    \item LHC LLP Working Group Convener \when{2022--25}
    \item SESAPS Nominating Committee \when{2022--23}
    \item Physics Today Consultant \when{2022}
    \internalonly{
      \item \emph{Journal Referee:} Universe \when{2022}
      \item \emph{Journal Referee:} Universe \when{2022}
      \item \emph{Journal Referee:} Instruments \when{2021}
      \item \emph{Journal Referee:} PRD \when{2021}
      \item \emph{Journal Referee:} LHEP \when{2021}
      \item \emph{Journal Referee:} PRL \when{2020}
      \item \emph{Journal Referee:} PRL \when{2018}
    }

    \externalonly{
      \item Journal Referee for Physical Review Letters, Physical Review D, \when{2018--}
      \item \hspace{2em} EPJC, Instruments, Letters in High Energy Physics, Universe
    }
  \end{itemizetighter}

\section*{Selected Scientific Presentations}

  \emph{Invited Seminars and Colloquiua}


  % \begin{itemize}
  \begin{itemizetighter}


    % \item University of Florida -- \textit{Physics Colloquium, Forthcoming} \when{2026}

    \item UC San Diego -- \textit{Physics Colloquium, Cottrell Lecture} \when{2026} % May 7, 2026

    \item San Diego State University -- \textit{Physics Colloquium, Cottrell Lecture} \when{2026} % May 8, 2026

    \item Cornell University -- \textit{HEP Seminar} \when{2026} % April 24, 2026

    \item Fermilab -- \textit{Wine \& Cheese Seminar} \when{2026} % 16 Jan 2026

    \item Yale University -- \textit{NPA Seminar} \when{2024} % 9 May 2024

    \item New York University -- \textit{Physics Colloquium} \when{2024} % March 2024

    \item University of California, Irvine, Virtual -- \textit{AI Seminar} \when{2022}
    % July 12, 2022

    \item University of Kansas, Lawrence, KS -- \emph{Particle Physics on the Plains} \when{2022}
    % March 8, 2022

    \item University of Tennessee, Knoxville, TN -- \emph{HEP Seminar} \when{2021}%
    % April 21, 2021
    % \item \textbf{R-Parity Violating SUSY is Hard to Find}
    %   \begin{itemizetighter}
    %   \item University of Tennessee, \emph{Seminar}, Knoxville, TN \when{2021}%
    %   \end{itemizetighter}

    \item University of Tennessee, Knoxville, TN -- \emph{Physics Colloquium} \when{2020}%
    %  Nov 23, 2020
    % \item \textbf{What To Do When Nature Is Unnatural}
    %   \begin{itemizetighter}
    %   \item University of Tennessee, \emph{Physics Colloquium}, Knoxville, TN \when{2020} % 11 December 2018
    %   \end{itemizetighter}

    \item Fermilab, Batavia, IL -- \emph{LPC Physics Forum} \when{2020}
    % 19 Nov 2020
    % \item \textbf{Displaced Vertex Searches at ATLAS}
    %   \begin{itemizetighter}
    %   \item Fermilab, \emph{LPC Physics Forum}, Batavia, IL \when{2020} % 11 December 2018
    %   \end{itemizetighter}


    \item Lund University, Sweden -- \emph{Seminar} \when{2020} % Mar 5, 2020
    \item Michigan State University, East Lansing, MI -- \emph{Physics Colloquium} \when{2019} % 5 February 2019
    \item University of Illinois, Urbana-Champaign, IL -- \emph{Physics Colloquium} \when{2018} % 25 October 2018
    \item Fermilab, Batavia, IL --  \emph{Topic of the Week} \when{2018} % 22 October 2018
    \item SLAC, Palo Alto, CA -- \emph{Seminar} \when{2018} % 18 October 2018
    \item University of Oregon, Eugene, OR -- \emph{Seminar} \when{2018} % 15 October 2018
    \item University of Pennsylvania, Philadelphia, PA -- \emph{Seminar} \when{2018} % 7 August 2018
    \item Rutgers University, Piscataway, NJ -- \emph{Seminar} \when{2018} % 6 August 2018

    % \item \textbf{How SUSY Can Still Save The Day / The SUSY Swindle}
    %   \begin{itemizetighter}
    %   \item Lund University, Sweden \when{2020}
    %   \item Michigan State University, East Lansing, MI \when{2019} % 5 February 2019
    %   \item University of Illinois, Urbana-Champaign, IL \when{2018} % 25 October 2018
    %   \item Fermilab \emph{Topic of the Week}, Batavia, IL \when{2018} % 22 October 2018
    %   \item SLAC, Palo Alto, CA \when{2018} % 18 October 2018
    %   \item University of Oregon, Eugene, OR \when{2018} % 15 October 2018
    %   \item University of Pennsylvania, Philadelphia, PA \when{2018} % 7 August 2018
    %   \item Rutgers University, Piscataway, NJ \when{2018} % 6 August 2018
    %   \end{itemizetighter}


    \item Harvard University, Cambridge, MA -- \emph{LPPC Seminar} \when{2018} % 11 October 2018
    % \item \textbf{On the Higgs Branching Ratios}
    %   \begin{itemizetighter}
    %   \item Harvard University - LPPC Seminar, Cambridge, MA\when{2018} % 11 October 2018
    %   \end{itemizetighter}


    \item Oskar Klein Centre, Stockholm, Sweden -- \emph{HEP Seminar}\when{2018} % 22 May 2018
    \item Lawrence Berkeley National Laboratory, Berkeley, CA  -- \emph{RPM Seminar}\when{2017} % 17 August 2017
    % \item \textbf{Searches for Sneaky SUSY at the ATLAS Experiment \& A Toy Entropic Explanation for the Hierarchy Problem}
    %   \begin{itemizetighter}
    %   \item Oskar Klein Centre, Stockholm, Sweden \when{2018} % 22 May 2018
    %   \item Lawrence Berkeley National Laboratory, Berkeley, CA \when{2017} % 17 August 2017
    %   \end{itemizetighter}

    \item Harvard University, Cambridge, MA -- \emph{LPPC Seminar}\when{2015} % 30 September 2015
  % \item \textbf{Searches for R-Parity Violating SUSY at the ATLAS Experiment}
  %   \begin{itemizetighter}
  %   \item Harvard University - LPPC Seminar, Cambridge, MA\when{2015}% 30 September 2015
  %   \end{itemizetighter}

    \item University of Adelaide, Australia -- \emph{CSSM Seminar}\when{2014} % 13 August 2014
    \item University of Melbourne, Australia -- \emph{Seminar}\when{2014} % 10 July 2014
    \item Institut de F\'isica d'Altes Energies (IFAE), Barcelona, Spain -- \emph{Seminar}\when{2014} % 7 January 2014

  % \item \textbf{A Search for B-Violating Supersymmetry in Multijet Signatures and Light-quark vs. Gluon Jet Tagging}
  %   \begin{itemizetighter}
  %   \item University of Adelaide, Australia\when{2014}% 13 August 2014
  %   \item University of Melbourne, Australia\when{2014}% 10 July 2014
  %   \item Institut de F\'isica d'Altes Energies (IFAE), Barcelona, Spain\when{2014} % 7 January 2014
  %   \item Santa Cruz Institute of Particle Physics Seminar, Geneva, Switzerland\when{2013}% 12
  %   \end{itemizetighter}

    \item 2nd International Spring School on Particle Physics and Philosophy -- \emph{Seminar}\when{2014}
  % \item \textbf{...But What If I Like Naturalness?}
  %   \begin{itemizetighter}
  %   \item 2nd International Spring School on Particle Physics and Philosophy
  %   \item Wuppertal, Germany\when{2014}% 31 March 2014
  %   \end{itemizetighter}

    \item Santa Cruz Institute of Particle Physics Seminar, Geneva, Switzerland -- \emph{Seminar}\when{2013} % 12

  % \item \textbf{RPV Stops at the LHC}
  %   \begin{itemizetighter}
  %   \item LPC Workshop on Exotic Top Partners
  %   \item Fermi National Accelerator Laboratory, LHC Physics Center, Batavia, IL\when{2013}% 26 September 2013
  %   \end{itemizetighter}


% \end{itemize}
\end{itemizetighter}

  \emph{Conference Presentations}

  % \begin{itemize}
  \begin{itemizetighter}

  \item DPF, \textit{Energy Frontier Session}, Fermilab \when{2026} % July 23 2026
  \item DPF, \textit{Outreach \& Engagement Session}, Fermilab \when{2026} % July 21 2026

  \item AI4MuC -- \textit{Overview Talk} \when{2026} % June 29, 2026


  \item IDEC Conference, Columbia College of Chicago  \when{2026} % Mar 2026

  \item Good $\nu$s, University of Pittsburgh -- \emph{Invited}  \when{2025} % 29 Oct 2025

  \item Physics At The Highest Energies With Colliders, Galileo Galilei Institute, Florence, Italy -- \emph{Invited}  \when{2025} % 29 July 2025


  \item LHCP2025, NTU, Taipei, Taiwan -- \emph{Invited}  \when{2025} % 4 March 2025


  \item Aspen Center for Physics -- \emph{Invited Muon Collider Overview Plenary}  \when{2025} % 4 March 2025

  \item Inaugural US Muon Collider Meeting, Fermilab -- \emph{Invited Experimental Overview Plenary}  \when{2024} % 7 August 2024

  \item Workshop on Long-Lived Particles, University of Tokyo, Tokyo, Japan -- \emph{o.b.o. CMS}  \when{2024} % 4 July 2024

  \item DPF/Pheno, Pittsburgh, PA \when{2024} % May 2024

  \item SESAPS2023, EKU, Richmond, KY -- \emph{Invited `Highlight' Talk}\when{2023} % XX Nov 2023

  \item Fermilab Users Meeting -- \emph{Invited Plenary} \when{2023}
  % June 28, 2023

  \item UCSB, Kavli Institute of Theoretical Physics, Muon Collider Workshop -- \emph{Invited Plenary} \when{2023}
  % March 8, 2023

  \item CMS Exotica Workshop -- \emph{Invited Plenary} \when{2022} % 11 Dec 2022

  \item ML4Jets, Rutgers University, Piscataway, NJ \when{2022} % 1 Nov 2022

  \item Snowmass Community Summer Study Workshop, Seattle, WA \when{2022} % 21 July 2022

  \item APS April Meeting, New York City, NY \when{2022} % 10 April 2022

  \item CMS Exotica Workshop -- \emph{Invited Plenary} \when{2021} % 24 Nov 2021

  \item SESAPS2021, FSU, Tallahassee, Fl -- \emph{Invited `Highlight' Talk}\when{2021} % 19 Nov 2021


  % \item \textbf{ColliderScope: Reaching New Audiences with Electronic(s) Music}
  %   \begin{itemizetighter}
  %   \item 18th International Particle Physics Outreach Group meeting, CERN \when{2019} % 28 November 2019
  %   \item International Conference on New Frontiers in Physics, Kolymbari, Crete \when{2019} % 21 August 2019
  %   \item APS Division of Particles and Fields Meeting, Boston, MA \when{2019} % 30 July 2019
  %   \end{itemizetighter}

  \item LHCP2021, Virtual -- \emph{Invited Plenary, o.b.o. ATLAS and CMS}\when{2021}
  % \item 6/9/2021 Neutral FIPs at the LHC and Beyond

  \item APS April Meeting, Virtual -- \emph{Invited Plenary, o.b.o. ATLAS}\when{2021}
  % \item \textbf{New Long-Lived Particle Results from the LHC} -- \emph{Plenary Talk} \when{2021} % 4/17/2021
  %   \begin{itemizetighter}
  %   \item APS April Meeting, Virtual
  %   \end{itemizetighter}

  \item A Rainbow Of Dark Sectors, Aspen Center for Physics, Virtual -- \emph{Invited Plenary}\when{2021}
  % \item \textbf{The Lifetime of the Long-Lived Particle Era at the LHC} -- \emph{Plenary Talk} \when{2021} % 24 Mar 2021
  %   \begin{itemizetighter}
  %   \item A Rainbow Of Dark Sectors, Aspen Center for Physics, Virtual
  %   \end{itemizetighter}


  \item LHC Experiments Committee (LHCC) Meeting, CERN, Geneva -- \emph{Open Session} \when{2020}
    % Nov 18, 2020
    % \item \textbf{Updates from the ATLAS Experiment}
    %   \begin{itemizetighter}
    %   \item LHC Experiments Committee (LHCC) Meeting Open Session, CERN, Geneva \when{2020}
    %   \end{itemizetighter}


  \item ICNFP, Kolymbari, Crete, Greece -- \emph{o.b.o. ATLAS}\when{2019} % 21 August 2019
  % \item \textbf{Searches for Long-Lived Particles with the ATLAS Detector}
  %   \begin{itemizetighter}
  %   \item \emph{On behalf of the ATLAS Collaboration}
  %   \item International Conference on New Frontiers in Physics, Kolymbari, Crete \when{2019} % 21 August 2019
  %   \end{itemizetighter}


  \item Higgs Cross-Section Working Group Meeting, CERN, Geneva -- \emph{Invited Plenary} \when{2018} % 11 December 2018
  % \item \textbf{Long-Lived Particles and the Higgs}
  %   \begin{itemizetighter}
  %   \item Higgs Cross-Section Working Group Meeting, CERN, Geneva \when{2018} % 11 December 2018
  %   \end{itemizetighter}

    \item USATLAS Workshop, Pittsburgh, PA -- \emph{Invited Plenary} \when{2018} % 1 August 2018
    % \item \textbf{SUSY and Exotics Overview} -- ``\emph{Headliner}'' Talk
    %   \begin{itemizetighter}
    %   \item USATLAS Workshop, University of Pittsburgh, Pittsburgh, PA \when{2018} % 1 August 2018
    %   \end{itemizetighter}

    % \item \textbf{The SUSY Paradox \& A Toy Entropic Explanation for the Hierarchy Problem}
    %   \begin{itemizetighter}
    %   \item Oskar Klein Centre, Stockholm, Sweden \when{2018} % 22 May 2018
    %   \end{itemizetighter}

  \item SUSY2017, Mumbai, India -- \emph{o.b.o. ATLAS} \when{2017} % 11 December 2017
  % \item \textbf{Searches for squarks and gluinos in scenarios with R-parity violating sparticle decays, or long- lived sparticles with ATLAS}
  %   \begin{itemizetighter}
  %   \item \emph{On behalf of the ATLAS Collaboration}
  %    \item SUSY2017, Mumbai, India \when{2017} % 11 December 2017
  %   \end{itemizetighter}

  \item IHEP-T2E, Kuala Lumpur, Malaysia -- \emph{Plenary, o.b.o. ATLAS} \when{2017} % 27 February 2017
  % \item \textbf{SUSY Searches in the ATLAS Experiment} -- \emph{Plenary Talk}
  %   \begin{itemizetighter}
  %   \item \emph{On behalf of the ATLAS Collaboration}
  %   \item IHEP-T2E, Kuala Lumpur, Malaysia \when{2017} % 27 February 2017
  %   \end{itemizetighter}

  \item SUSY2016, Melbourne, Victoria, Australia\when{2016} % 3 July 2016
  % \item \textbf{The Recursive Jigsaw Reconstruction Technique}
  %   \begin{itemizetighter}
  %   \item SUSY2016, Melbourne, Victoria, Australia \when{2016} % 3 July 2016
  %   \end{itemizetighter}

  % \item \textbf{New Searches for Strongly-Produced SUSY at the ATLAS Experiment}
  %   \begin{itemizetighter}
  %   \item CoEPP Annual Workshop, Torquay, Victoria, Australia\when{2016} % 18 February 2016
  %   \end{itemizetighter}

  \item SUSY2015, Lake Tahoe, CA\when{2015} % 24 August 2015
  % \item \textbf{Applications of the Recursive Jigsaw Technique}
  %   \begin{itemizetighter}
  %   \item SUSY2015, Lake Tahoe, CA\when{2015}% 24 August 2015
  %   \end{itemizetighter}

  \item Kruger2014, Kruger Gate, South Africa -- \emph{o.b.o. ATLAS} \when{2014} % 6 December 2014
  % \item \textbf{SUSY Searches in the ATLAS Experiment} -- \emph{Plenary Talk}
  %   \begin{itemizetighter}
  %   \item \emph{On behalf of the ATLAS Collaboration}
  %   \item Kruger2014, Kruger Gate, South Africa\when{2014}% 6 December 2014
  %   \end{itemizetighter}


  \item Fermilab, LPC Workshop on Exotic Top Partners, Batavia, IL -- \emph{Invited Plenary}\when{2013} % 26 September 2013

  \item SUSY2012, Peking University, Beijing, China -- \emph{o.b.o. ATLAS} \when{2012} % 17 August 2012
  % \item \textbf{Search for supersymmetry in resonant production and R-parity violating signatures with the ATLAS detector}
  %   \begin{itemizetighter}
  %   \item \emph{On behalf of the ATLAS Collaboration}
  %   \item SUSY2012, Peking University, Beijing, China\when{2012}% 17 August 2012
  %   \end{itemizetighter}

% \end{itemize}
\end{itemizetighter}


  \emph{Conference and Workshop Organization}

\begin{itemizetighter}

  \item \textit{USMCC 2026}, Stanford University, \emph{Co-chair (w/ M. Peskin), Program Committee} \when{2026} % Dec 2026
  \item \textit{USMCC 2025}, University of Chicago, IL, \emph{Detector Design and Performance Session Organizer [of 3]} \when{2025} % Aug 2025
  \item \textit{LHCP 2025}, Taipei, Taiwan, \emph{Poster Committee Member} \when{2025} % May 2025
  \item \textit{CPAD 2024}, Knoxville, TN, \emph{Chair of Local Organizing Committee [of 5]} \when{2024} % Nov 2024
  \item \textit{LHCP 2024}, Northeastern University, Boston, MA, \emph{Poster Committee Member} \when{2024} % May 2024

  \item \textit{The Future of High Energy Physics}, Aspen Center for Physics, \emph{Conference Organizer [of 4]} \when{2024} % Mar 2024
  \item \textit{PITT PACC Muon Collider Physics Benchmark Workshop}, Pittsburgh, PA, \emph{Invited Closing Remarks} \when{2023} % Nov 2023
  \item \textit{ML4Jets}, Rutgers University, \emph{Session Chair} \when{2022} % Oct 2022
  \item \textit{APS April Meeting}, New York, NY, \emph{Session Chair} \when{2022} % April 10 2022
  \item \textit{ATLAS SUSY and Exotics Workshop}, Virtual \when{2020} % September 2020
  \item \textit{ATLAS Reaching New Heights Workshop}, CERN \when{2019} % December 2019
  \item \textit{ATLAS SUSY Workshop}, Lecce, Italy \when{2019} % September 2019
  \item \textit{APS Division of Particles and Fields Meeting}, Northeastern University \when{2019} % 29 July - 2 August 2019
  \item \textit{ATLAS SUSY Workshop}, Stockholm, Sweden \when{2018} % 23-25 May 2018
  \item \textit{ATLAS SUSY and Exotics Workshop}, Bucharest, Romania \when{2017} % 8-12 May 2017
  \item \textit{Recursive Jigsaw Workshop}, Harvard University \when{2015} % 21 September 2015
  \item \textit{Workshop on LHC Searches}, LBNL \when{2014} % 30 January 2014

  \end{itemizetighter}

% \clearpage


\section*{Teaching}

  \begin{itemizetighter}
    \item Physics of Cosmic Rays \emph{Seminar \& Design/Build Course w/ UT Architecture \& Design} (PHYS 494) \when{2025}
    \item Classical Mechanics I \textit{\&} II (PHYS 311/312) \when{2021--24}
    \item UTK HEP/Astro Seminar    \when{2022}
    \item Graduate Research Participation Seminar  \when{2021}
  \end{itemizetighter}


  \internalonly{

    \section*{Current Student \textit{\&} Postdoc Mentorship}

    \begin{itemizetighter}
      \item Olivia Clark [Undergrad] \when{2025--}
      \item Micah Hillman [Undergrad] \when{2025--}
      \item \rightnote{UTK EUReCA 2nd Place Presenter Award}

      \item Allie Dabney [Undergrad] \when{2025--}
      \item Shivam [Ph.D.] \when{2024--}
      \item Caroline Riggall [Ph.D.] \when{2024--}
      \item Cordney Nash [Undergrad] \when{2024--}
      \item \rightnote{UTK `Volunteer of Distinction'}

      \item Colby Thompson [Ph.D.] \when{2023--}
      \item Adam Vendrasco [Ph.D.] \when{2022--}
      \item John Lawless [Ph.D.] \when{2022--}
      \item \rightnote{Stelson Fellowship for Outstanding Beginning Research}

      \item Emery Nibigira [Postdoc] \when{2022--}
      \item \rightnote{Fermilab LPC Distinguished Researcher}

      \item \rightnote{Lindau Nobel Laureate Meeting Selected Participant}

    \end{itemizetighter}


    \section*{Student Thesis Committee Membership}

    \begin{itemizetighter}
      \item Colby Thompson [Ph.D.], \emph{chair} \when{2025--}
      \item Raghav Chari [Undergrad, College Honors] \when{2025}
      \item John Lawless [Ph.D.], \emph{chair} \when{2024--}
      \item Brandi Skipworth [Ph.D.] \when{2024--}
      \item Christal Martin [Ph.D.] \when{2024--}
      \item Ewa Glimos [Ph.D.] \when{2024--}
      \item Taylor Sussmane [Undergrad, Honors], \emph{chair} \when{2024}
      \item John Melhorn [Undergrad, Honors] \when{2024}
      \item Luke Carpenter [Undergrad, Honors] \when{2024}
      \item Rebecca Godri [Ph.D.] \when{2023--}
      \item Daniel Ally [Ph.D.] \when{2022--}
      \item Jesse Harris [Ph.D.] \when{2022--}
      \item Satyajit Roy [Ph.D.] \when{2021--24}
      \item Austin Schmier [Ph.D.] \when{2021--24}
      \item Himal Acharya [Ph.D.] \when{2021}
      \item External Ph.D. Dissertation Review for Yale University \when{2021}
      \item External M.S. Thesis Review for University of Adelaide \when{2021}
    \end{itemizetighter}

  }

  \section*{Past Student \textit{\&} Postdoc Mentorship}


  \textit{n.b. Square brackets denote initial position of mentorship relationship. Current position in italics.}

    \begin{itemizetighter}
      \item Dr. Alexander Tuna [Postdoc, UTK] \hfill \emph{Research Staff, UCSD}
      \item Dr. Ankush Kanuganti [Postdoc, UTK] \hfill \emph{Postdoc, BNL}
      \item Micah Hillman [Undergrad, UTK] \hfill \emph{Physics Postbacc, Princeton}
      \item Charles Bell [Undergrad, UTK] \hfill \emph{Physics PhD Student, UMichigan, NSF GRFP Awardee, Cover of Science,}
        \item \hfill \emph{UTK Talley Award for Outstanding Undergrad Research,}
        \item \hfill \emph{UTK Douglas V. Roseberry Distinguished Upper Class Major Award}
      \item Lacey Dishman [Undergrad, UTK] \hfill \emph{Physics PhD Student, Boston University}
      \item Taylor Sussmane [Undergrad, UTK] \hfill \emph{Physics PhD Student, University of Wisconsin, Madison,}
        \item \hfill \emph{Douglas V. Roseberry Distinguished Upper Class Major Award,}
        \item \hfill \emph{UTK Talley Award for Outstanding Undergrad Leadership}
      \item Everett Hagen [Undergrad, UTK]
      \item Jack Peltier [Undergrad, UTK]\hfill \emph{Goldwater Scholar}
      \item Puck Cravens [Undergrad, UTK]
      \item Tanner Wolfe [Undergrad, UTK]
      \item Alexander Sizemore [Undergrad, UTK]
      \item Matthew Durkee [Undergrad, UTK] \hfill \emph{Defense Contractor}
      \item Philipp Wagenknecht [Ph.D., UTK] \hfill \emph{HEP Lab Technician, Brown University}
      \item Dr. Tamas Vami [Ph.D., Johns Hopkins] \hfill \emph{Postdoc, UCSB}
      \item Prof. Karri Folan DiPetrillo [Ph.D., Harvard] \hfill \emph{Asst. Professor, University of Chicago}
      \item Prof. Jennifer Roloff [Ph.D., Harvard] \hfill \emph{Asst. Professor, Brown University}
      \item Dr. Anthony Badea [Ph.D., Harvard] \hfill \emph{Schmidt AI Fellow, University of Chicago}
      \item Dr. Anne Fortman [Ph.D., Harvard] \hfill \emph{Postdoc, LBNL}
      \item Dr. Aaron White [Postdoc, Harvard] \hfill \emph{}
      \item Xiaohe Jia [Ph.D., Harvard] \hfill \emph{Industry}
      \item Madeline Bernstein [Undergrad, Harvard] \hfill \emph{Physics PhD Student, UC Berkeley}
      \item Dr. Stefanie Morgenstern [Ph.D., Dresden/CERN] \hfill \emph{CERN Fellow}
      \item Dr. Russell Smith [Ph.D., Columbia] \hfill \emph{Squarepoint, Financial Sector}
      \item Prof. Emma Alexander [Undergrad, Yale] \hfill \emph{Asst. Professor, Northwestern University}
      \item Dr. Ariel Ekblaw [Undergrad, Yale] \hfill \emph{Founder of the MIT Space Exploration Initiative, CEO of Aurelia Institute}
    \end{itemizetighter}


\section*{Selected Outreach Activities}

  \begin{itemize}

    \item \textbf{\href{https://cosmicrayn.utk.edu}{Cosmic Rayn: Cosmic Ray Educational Experience}} w/ UTK Architecture \textit{\&} Design \when{2024--26}

    \item \textbf{Quantum Canvases: Physics, the Arts, and the Humanities} \when{2023}

      \begin{itemizetighter}
      \item Lead organizer of a series of events in downtown Knoxville on the intersection of physics, the arts, and the humanities in collaboration with UTK Physics \textit{\&} Astronomy, CERN, and the UT Humanities Center.
      \item Electronic music event (``Harmonic Motion: Physics x Electronic Music'')
      \item Events around a collection of science fiction short stories
      \item A panel discussion on the intersection of physics and science fiction
      \item A humanities-facing physics colloquium event
      \end{itemizetighter}


    \item \textbf{Higgs Tenth Anniversary Celebration} \when{2022}

      \begin{itemizetighter}
      \item Primary organizer for a documentary screening of \emph{Particle Fever} and open panel session for the tenth anniversary of the discovery of the Higgs Boson. Held at Central Cinema in Knoxville, TN.
      \end{itemizetighter}


    \item \textbf{STEAM Trailer Design Project} \when{2021--22}

      \begin{itemizetighter}
      \item Worked with the UTK College of Architecture \textit{\&} Design to design and construct an educational experience for local middle schoolers with a focus on the physics of sound and visualization.
      \end{itemizetighter}

    \item \textbf{Keynote presentation at event for Eureka Science Museum}, Grand Junction, CO \when{2021}


    \item \href{http://www.cern.ch/ColliderScope}{\textbf{ColliderScope -- Particle-Physics-Inspired Electronic Music Project}} \when{2019--}

      \begin{itemizetighter}
      \item Music with sound waves that show particle physics images in Lissajous figures
      \item Live performances:
      \item \hspace{2em} Luck Reunion (at Willie Nelson's ranch), Austin, TX (\emph{w/ Micah Nelson, Daniel Lanois}) \when{2026}
      \item \hspace{2em} Fly By Night, Knoxville, TN (\emph{w/ signal:noise; VJ Set}) \when{2026}
      \item \hspace{2em} Fly By Night, Knoxville, TN (\emph{w/ signal:noise; VJ Set}) \when{2025}
      \item \hspace{2em} Fly By Night, Knoxville, TN (\emph{w/ signal:noise; VJ Set}) \when{2025}
      \item \hspace{2em} Scruffy City Hall, Knoxville, TN (\emph{w/ Harmonic Motion}) \when{2025}
      \item \hspace{2em} Scruffy City Hall, Knoxville, TN (\emph{w/ Harmonic Motion}) \when{2024}
      \item \hspace{2em} Old City Performing Arts Center, Knoxville, TN (\emph{w/ Harmonic Motion}) \when{2023}
      \item \hspace{2em} Blue Moon Tavern, Seattle, WA (\emph{w/ Snowmass summer meeting}) \when{2022}
      \item \hspace{2em} Roskilde Festival, Denmark \when{2022}
      \item \hspace{2em} Cross Club, Prague, YouTube, Live from the CERN Control Centre (\emph{w/ ICHEP2020}) \when{2020}
      \item \hspace{2em} Lund, Sweden \when{2020}
      \item \hspace{2em} Aeronaut Brewing, Somerville, MA \when{2019}
      \item \hspace{2em} ATLAS Experiment, Geneva, Switzerland \when{2019}
      \item \hspace{2em} CERN Open Days, Geneva, Switzerland \when{2019}
      \item \hspace{2em} Pohoda Music Festival, Tren\v{c}\'{i}n, Slovakia \when{2019}
      \item Commissioned Compositions:

      \item \hspace{2em} Intro Videos for the CERN Science Gateway Educational Program \when{2023}

      \item \hspace{2em} Intro Videos for CERN Music Festival Program \when{2019}

      \item Presented at multiple international and national conferences
    %   \item 18th International Particle Physics Outreach Group meeting, CERN \when{2019 % 28 November 2019}
    %   \item International Conference on New Frontiers in Physics, Kolymbari, Crete \when{2019 % 21 August 2019}
    %   \item APS Division of Particles and Fields Meeting, Boston, MA \when{2019 % 30 July 2019}
      \end{itemizetighter}


    \item \textbf{Presentations to Grade School Students}
    \begin{itemizetighter}
    \item L\&N STEM Academy, Knoxville, TN \when{2023}
    \item Tel Aviv University -- Future Scientists CERN Tour \when{2017--19}
    \item Five presentations to NJ and NY schools \when{2011--16}
    \end{itemizetighter}

    \item \textbf{Harvard NSW Activities at CERN} - \emph{Video} (\href{https://youtu.be/VHrJAkNsMhM}{YouTube})
    \begin{itemizetighter}
    \item Fully produced video for recruitment \when{2018}
    \end{itemizetighter}

  \item \href{http://www.cern.ch/inparticular}{\textbf{In Particular}} - \emph{Podcast}\when{2015--16}

    \begin{itemizetighter}
    \item Video, audio, and original music production
    % \item Production of ``Video Shorts''
    % \item Web development, graphic design, and episode production
    \item Tens of thousands of downloads ($\sim11$~TB served)
    \item Featured by CERN, iTunes, The Guardian Science Blogs, Ars Technica, Vice, and others
    \end{itemizetighter}


  \item \textbf{Visiting Scholar, Dharma Drum Mountain, Taiwan}
    \begin{itemizetighter}
    \item Several weeks of conversations with monastery's monks and academics \when{2014}
    \end{itemizetighter}


  \end{itemize}


\section*{Press \textit{\&} Interviews}

      \begin{itemizetighter}

      \item \href{https://physics.utk.edu/tova-holmes-and-larry-lee-selected-as-fermilab-distinguished-researchers/}{UTK Physics -- Tova Holmes And Larry Lee Selected As Fermilab Distinguished Researchers}\when{2026}

      \item \href{https://news.utk.edu/2025/02/12/university-of-tennessee-physicist-named-cottrell-scholar/}{UTK News -- University of Tennessee Physicist Named Cottrell Scholar}\when{2025}

      \item \emph{Science} -- L.~Lee, C.~Bell, \textbf{\href{https://www.science.org/toc/science/383/6690}{Cover of Issue 383 6690}}  \when{2024}


      \begin{itemizetightrightpad}
        \item \emph{Muon collider simulation work and renderings made in Unreal Engine were featured on the cover and in the cover article in the issue from March 29, 2024.}
      \end{itemizetightrightpad}

      \item \href{https://engage.aps.org/dpb/resources/newsletters/24-25-newsletter/f10-article}{APS DPB Annual Newsletter -- The Muon Collider: A Modern Machine for a Modern World}\when{2024}
      \item \href{https://cerncourier.com/a/a-new-generation-a-new-vision/}{CERN Courier -- A New Generation, A New Vision} \when{2024}
      \item \href{https://physics.utk.edu/the-particle-detectorists/}{UTK Physics -- The (Particle) Detectorists} \when{2024}
      \item \href{https://physics.utk.edu/four-ut-physics-students-win-doe-scgsr-support/}{UTK Physics -- NSF GRFP Awardee} \when{2024}
      \item \href{https://physics.utk.edu/the-art-of-muon-collisions/}{UTK Physics -- The Art of Muon Collisions} \when{2024}
      \item Interview for article for MIT Tech Review \when{2023}
      \item Interview for article for Symmetry Magazine on CERN music festival program \when{2023}
      \item Interview with Sparks Magazine on Asian Americans in STEM\when{2023}
      \item \href{https://physics.utk.edu/in-my-own-words-undergraduate-taylor-sussmane/}{UTK Physics -- In My Own Words: Undergraduate Taylor Sussmane} \when{2023}
      \item \href{https://physics.utk.edu/challenging-the-standard-model-lawrence-lee-wins-nsf-career-award/}{UTK Physics -- Challenging The Standard Model: Lawrence Lee Wins NSF CAREER Award}\when{2023}
      \item \href{https://physics.utk.edu/the-art-of-science/}{UTK Physics -- The Art of Science} \when{2023}
      \item \href{https://physics.utk.edu/the-yeti-returns/}{UTK Physics -- The YETI Returns} \when{2023}
      \item Interview with Danish podcast TechTopia \when{2022}
      \item \textit{ColliderScope} featured in \href{https://www.symmetrymagazine.org/article/lhc-music-through-the-colliderscope}{Symmetry Magazine} \when{2019}
      \item Interview for article for Symmetry Magazine on naturalness \when{2019}
      \item CERN Courier -- ATLAS Pushes the Limits on Supersymmetry \when{2019}

      \item Interview for Swedish-language podcast \textit{Professor Magenta} about science and art \when{2016}
      \item Public screening of film Particle Fever w/ panel Q\textit{\&}A\when{2014}


      \end{itemizetighter}


\section*{University, College, and Departmental Committees and Other Service Roles}

  \begin{itemizetighter}

    \item UTK Academy for Global Engagement - Mentor \when{2026--27}

    \item CAS Climate, Culture, and Community Committee \when{2025--}

    \item Faculty Search Committee - Particle Theory \when{2025--26}
    \item Faculty Search Committee - Lecturer \when{2025--26}

    \item ORIED Launch Pad Panel on Prestigious Fellowships \& Awards \when{2025}
    \item University NSF CAREER Workshop Panel \when{2025}
    \item Faculty Search Committee - Nuclear Theory \when{2024--25}

    \item Coordinator of Department Web Services \when{2024--}

    \item 4 UTK Senior Thesis Committees \when{2024--}

    \item University NSF CAREER Workshop Panel \when{2024}

    \item Provost Junior Faculty Fellows (\emph{UTK}) \when{2023--25}

    \item Graduate Recruitment Committee  \when{2023--}

    \item SPS Faculty Advisor  \when{2023--}

    \item Assessment Committee  \when{2023--24}

    \item Graduate Advising Committee \when{2022--}

    \item Planning Committee \when{2022--24}

    \item Faculty Search Committee - Medium Energy Nuclear Experiment \when{2022--23}

    \item Qualifying Exam Committee \when{2022--23, 24--25}

    \externalonly{
      \item 9 UTK PhD Thesis Committees \when{2021--}

    }

    \item Community Connections Committee (\emph{Chair}) \when{2021--23}

    \item Department Head Search Committee  (\emph{UTK College of Arts \& Sciences}) \when{2021--22}
  \end{itemizetighter}


% \section*{Professional Activities}

%   \begin{itemize}
%   \item Member, American Association for the Advancement of Science, 2013--Present
%   \item Member, New York Academy of Sciences, 2012--Present
%   \item Member, American Physical Society, 2008--2009
%   \end{itemize}

% \section*{Advising}


%   \begin{itemizetight}

%   \item Harvard University:
%   \begin{itemizetighter}
%     \item Brendon Bullard (PhD Candidate) \hfill[\emph{w/ M. Morii}]
%     \item Jennifer Roloff (2019 PhD) \hfill[\emph{w/ J. Huth}]
%     \item Karri Folan DiPetrillo (2019 PhD) \hfill[\emph{w/ M. Franklin}]
%     \item Adriana Rotaru (2019 Summer Undergraduate)  \hfill[\emph{w/ J. Huth}]
%     \item Madeline Bernstein (2017 Summer Undergraduate)  \hfill[\emph{w/ M. Franklin}]
%   \end{itemizetighter}

%   % \item Uppsala University, Students:
%   % \begin{itemizetighter}
%   %   \item Andreas Gillgren (2018-2019 Visiting Engineering Physics Student)  \hfill[\emph{w/ J. Huth}]
%   % \end{itemizetighter}

%   \item Columbia University:
%   \begin{itemizetighter}
%     \item Russell Smith (2016 PhD)  \hfill[\emph{w/ E. Hughes}]
%   \end{itemizetighter}

%   \item University of Adelaide, Three PhD and Two Undergraduate Students
%   % \begin{itemizetighter}
%   %     \item Jason Oliver (2019 PhD) \hfill[\emph{w/ P. Jackson}]
%   %     \item Anum Qureshi (2019 PhD) \hfill[\emph{'' ''}]
%   %     \item Damir Duvnjak (2014 Honours, 2019 PhD) \hfill[\emph{'' ''}]
%   %     \item Jens Kr\"oger (2015 Visiting) \hfill[\emph{'' ''}]
%   %     \item Finley Borgas (2014 Honours) \hfill[\emph{'' ''}]
%   % \end{itemizetighter}

%   \item Yale University, Five Undergraduate Honors Students
%   % \begin{itemizetight}
%   %     \item Madeleine Barrow (2012)
%   %     \item Ari Brill (2012)
%   %     \item Emma Alexander (2011)
%   %     \item James Giammona (2011)
%   %     \item Alexandra Turrini (2011)
%   % \end{itemizetight}
%   \end{itemizetight}

% \section*{Teaching}

% % \begin{changemargin}{1.5em}{0.0in}

%   \begin{itemizetight}
%     \item Yale University
%     \begin{itemizetight}
%       \item \emph{Fundamentals of Physics}, lecturing assistant for Prof. P. Tipton \when{2011}
%       \item \emph{Physics for the Life Sciences}, lecturing assistant for Prof. S. Mochrie \when{2010}
%       \item \emph{Modern Physical Measurement}, assistant for Prof. J. Harris \when{2010}
%       \item \emph{General Physics Laboratory}, assistant for Prof. K. Baker \when{2009}
%     \end{itemizetight}
%   \end{itemizetight}
% \end{changemargin}

% \clearpage
% \section*{References}

%   \begin{itemize}
%   \item \textbf{Prof. Melissa Franklin} - {\it Mallinckrodt Professor of Physics} - Harvard University
%   \item \hspace{1.6em} \href{mailto:Franklin@Physics.Harvard.edu}{Franklin@Physics.Harvard.edu}
%   \item \textbf{Prof. Tobias Golling} - {\it Associate Professor} - University of Geneva
%   \item \hspace{1.6em} \href{mailto:Tobias.Golling@unige.ch}{Tobias.Golling@unige.ch}
%   \item \textbf{Dr. Andreas Hoecker} - {\it Senior Research Physicist, ATLAS Spokesperson-Elect} - CERN
%   \item \hspace{1.6em} \href{mailto:Andreas.Hoecker@cern.ch}{Andreas.Hoecker@cern.ch}
%   \item \textbf{Prof. John Huth} - {\it Donner Professor of Science} - Harvard University
%   \item \hspace{1.6em} \href{mailto:Huth@g.Harvard.edu}{Huth@g.Harvard.edu}

%     % \pagebreak

%     % \item {\textit{\hspace{-1.6em} [References continued...] }}

%   \item \textbf{Prof. Paul Jackson} - {\it Associate Professor} - University of Adelaide
%   \item \hspace{1.6em} \href{mailto:P.Jackson@adelaide.edu.au}{P.Jackson@adelaide.edu.au}
%   \item \textbf{Dr. Zachary Marshall} - {\it Senior Scientist} - Lawrence Berkeley National Lab
%   \item \hspace{1.6em} \href{mailto:Zach.Marshall@cern.ch}{Zach.Marshall@cern.ch}

%   % \vspace{1em}
%   % \item \textit{Additional References}
%   % \begin{itemize}
%   % \item \textbf{Prof. Paul Tipton} - {\it Chair of Physics Department} - Yale University
%   % \item \hspace{1.6em} \href{mailto:paul.tipton@yale.edu}{Paul.Tipton@Yale.edu}
%   % \end{itemize}

%   \end{itemize}

\cvfooter

\end{document}
"
%%
%% which is what build.sh and .github/workflows/build-cv.yml do.

\newif\ifCVinternal
\ifdefined\CVinternal \CVinternaltrue \else \CVinternalfalse \fi

% \internalonly{<content>} -- included only in the internal edition.
\newcommand{\internalonly}[1]{\ifCVinternal #1\fi}

% \externalonly{<content>} -- included only in the external edition.
\newcommand{\externalonly}[1]{\ifCVinternal\else #1\fi}

% The internal edition says so in its PDF title, so the two files stay
% distinguishable once they have been detached from their filenames.
\ifCVinternal
  \def\LLedition{\name: Curriculum Vitae (internal edition)}
\else
  \def\LLedition{\name: Curriculum Vitae}
\fi

\hypersetup{
  colorlinks  = false,
  pdfauthor   = {\name},
  pdftitle    = {\LLedition},
  pdfsubject  = {Curriculum Vitae},
  pdfkeywords = {},
  pdfpagemode = UseNone,
}


%% ===========================================================================
%% Page layout
%% ===========================================================================

\geometry{
  body       = {6.5in, 9.0in},
  left       = 1.0in,
  top        = 1.0in,
  headheight = 14pt,   % enough room for the running header; silences fancyhdr
}

\pagestyle{fancy}
\fancyhf{}
\renewcommand{\headrulewidth}{0pt}
\fancyhead[L]{\name}
\fancyhead[R]{\thepage}

\setlength{\parindent}{0em}

\titlespacing{\section}{0pt}{0pt}{0pt}
\sectionfont{\rmfamily\mdseries\Large}
\subsectionfont{\rmfamily\mdseries\itshape\large}


%% ===========================================================================
%% Lists
%% ===========================================================================
%% The CV is built almost entirely from bullet-free lists that differ only in
%% how tightly they are set.  itemize itself is redefined, rather than
%% configured through enumitem, so that a plain \begin{itemize} anywhere picks
%% up the CV style.
%%
%%   itemize               default spacing between entries
%%   itemizetight          no extra space between entries
%%   itemizetighter        entries set as close as the leading allows
%%   itemizetightrightpad  tight, and inset from the right margin, for the
%%                         commentary set underneath an entry
%%
%% \LLlist{<left>}{<right>}{<itemsep>}{<parskip>}{<parsep>}

\newcommand{\LLlist}[5]{%
  \list{}{%
    \setlength{\leftmargin}{#1}%
    \setlength{\rightmargin}{#2}%
    \setlength{\itemsep}{#3}%
    \setlength{\parskip}{#4}%
    \setlength{\parsep}{#5}%
  }%
}

\renewenvironment{itemize}%
  {\LLlist{1.5em}{0em}{0.25em}{0pt}{0.25em}}{\endlist}
\newenvironment{itemizetight}%
  {\LLlist{1.5em}{0em}{0em}{-1pt}{0em}}{\endlist}
\newenvironment{itemizetighter}%
  {\LLlist{1.5em}{0em}{0em}{-0.4em}{0em}}{\endlist}
\newenvironment{itemizetightrightpad}%
  {\LLlist{1.5em}{3em}{0em}{-1pt}{0em}}{\endlist}

% Hanging indent for a run-on entry: \parshape 2 3.2em \linewidth \listindent
% \dimexpr\linewidth-\listindent\relax sets the first line full width and
% indents every line after it to \listindent.
\newlength{\listindent}
\setlength{\listindent}{4.8em}

% \changemargin{<left>}{<right>} ... \endchangemargin -- indents a block of the
% document; used to nest the collaboration sub-lists under their headings.
\def\changemargin#1#2{\list{}{\rightmargin#2\leftmargin#1}\item[]}
\let\endchangemargin=\endlist


%% ===========================================================================
%% The right-hand date column
%% ===========================================================================
%% Nearly every entry ends with a date set against the right margin.  \when
%% supplies both the \hfill and the fixed-width box that lines those dates up
%% into one column: the box is left-aligned, so "2024" and "2022--24" begin at
%% the same horizontal position.
%%
%% A \makebox is only as wide as it is declared, so an over-long date used to
%% overflow it silently and print out in the right margin -- six of the dates in
%% this CV ("2021--2026", "2022--23, 24--25", ...) are wider than the column.
%% \when measures the date first and sets an over-wide one flush right instead,
%% which costs the shared left edge on those few entries but keeps every date
%% inside the text block.

\newlength{\datecolumn}
\setlength{\datecolumn}{3.5em}

%% The glue ahead of a date is \datesep plus infinite stretch.  The stretch is
%% what pushes the date to the right margin; the fixed \datesep is a floor, so
%% an entry that grows long enough to reach its date reports an overfull box
%% instead of quietly printing its text hard against the year.  Every entry
%% currently clears it, so the floor costs nothing today -- it is there to catch
%% the next entry that does not.  \whenflush is the fix when one turns up: it
%% gives back the empty right-hand part of the date box, about 1.6em for a
%% four-digit year.
\newlength{\datesep}
\setlength{\datesep}{0.4em}

%% \LLdateglue is the glue that carries the date to the right margin.  It reads
%% as one \hfill most of the time, but is written out so that an entry too long
%% to hold its date can break: TeX takes the \penalty50, leaves the text on the
%% line it filled, and the \hbox{} -- which survives the break where glue would
%% not -- carries the second \hfill onto the next line, so the date still lands
%% flush right instead of stranded against the left margin.  \unskip drops the
%% source space ahead of the date, so the gap is \datesep however the entry was
%% typed.
\newcommand{\LLdateglue}{%
  \unskip\nobreak\hspace{\datesep}\hfill\penalty50 \hbox{}\nobreak\hfill
}

\newlength{\LLdatewidth}
\newcommand{\when}[1]{%
  \LLdateglue
  \settowidth{\LLdatewidth}{#1}%
  \ifdim\LLdatewidth>\datecolumn #1\else\makebox[\datecolumn][l]{#1}\fi
}

% \whenflush{<date>} -- for an entry whose text runs so close to the right
% margin that the shared left edge of the date column will not fit.  Sets the
% date flush right, giving up the alignment to keep the gap.
\newcommand{\whenflush}[1]{\LLdateglue #1}

% \rightnote{<text>} -- a right-aligned italic annotation (an award, a current
% position) hung underneath an entry.  The trailing pad makes it end where a
% four-digit year in the date column ends, rather than at the margin itself.
\newcommand{\rightnote}[1]{%
  \settowidth{\LLdatewidth}{2024}%
  \hfill\emph{#1}\hspace{\dimexpr\datecolumn-\LLdatewidth\relax}%
}


%% ===========================================================================
%% Live publication metrics from INSPIRE-HEP
%% ===========================================================================
%% Citation counts go stale the moment they are typed, so they are not written
%% into the CV by hand.  tools/fetch_inspire.py queries the INSPIRE-HEP API and
%% writes LL_InspireData.tex, a flat list of declarations:
%%
%%     \inspiresetcites{2690093}{3}         % citations of one paper
%%     \inspiresetstat{citations}{207,000}  % an author-level summary figure
%%
%% which the CV body reads back through \inspirepub and \inspirestat.  The
%% values set just below are fallbacks for a build that has never run the
%% fetcher -- an offline checkout, or Overleaf -- so the document always
%% compiles with plausible numbers.  build.sh and CI refresh the data first.

% Declarations written by tools/fetch_inspire.py.  The keys live in \csname, so
% the LL@ namespace prefix needs no \makeatletter here.
\newcommand{\inspiresetcites}[2]{\expandafter\gdef\csname LL@cites@#1\endcsname{#2}}
\newcommand{\inspiresetstat}[2]{\expandafter\gdef\csname LL@stat@#1\endcsname{#2}}

% Fallbacks: the last hand-checked figures, overridden by LL_InspireData.tex.
\inspiresetstat{papers}{1400}
\inspiresetstat{citations}{202,000}
\inspiresetstat{hindex}{209}

% \inspirestat{<key>} -- an author-level figure: papers, citations or hindex.
% A missing key is reported loudly rather than left silently blank.
\newcommand{\inspirestat}[1]{%
  \ifcsname LL@stat@#1\endcsname\csname LL@stat@#1\endcsname\else\textbf{[?~#1]}\fi
}

% \inspirecites{<record>} -- "[n citations]" for a record whose count is known.
% Nothing is printed for an unknown count, or for one below \citesmin, so a
% brand-new paper is left unannotated rather than advertising that nobody has
% cited it yet.
%
% The annotation costs about a dozen characters on every entry, which is enough
% to push an entry that already filled its line onto a second one.  Raising
% \citesmin, or shortening the label, is the dial if that happens too often.
\newcommand{\citesmin}{1}
\newcommand{\LLprintcites}[1]{%
  \ifnum#1<\citesmin\relax\else
    \nobreakspace{\small[#1\nobreakspace\ifnum#1=1\relax citation\else citations\fi]}%
  \fi
}
\newcommand{\inspirecites}[1]{%
  \ifcsname LL@cites@#1\endcsname\LLprintcites{\csname LL@cites@#1\endcsname}\fi
}

% \inspirepub{<record>}{<title>} -- a publication title in bold, linked to its
% INSPIRE-HEP record and followed by that record's live citation count.
\newcommand{\inspirepub}[2]{%
  \href{https://inspirehep.net/literature/#1}{\textbf{#2}}\inspirecites{#1}%
}

% \pub{<url>}{<title>} -- the same, for an entry with no INSPIRE-HEP record:
% arXiv-only preprints, CDS and CERN notes, bare DOI links.
\newcommand{\pub}[2]{\href{#1}{\textbf{#2}}}

% The generated data, if it has been fetched.  \IfFileExists keeps a fresh
% checkout compiling on the fallbacks above.
\IfFileExists{LL_InspireData.tex}{%% Publication metrics fetched from the INSPIRE-HEP API.
%% GENERATED FILE -- do not edit; run tools/fetch_inspire.py instead.
%%
%% \inspiresetcites{<record>}{<count>}  per-paper citation count
%% \inspiresetstat{<key>}{<value>}      author-level summary figure
%% Both are defined in LL_Preamble.tex.

%% Author-level summary, as shown at the top of the Publications section.
%% Papers and citations are rounded down because the CV says "over".
\inspiresetstat{papers}{1400}    % exact: 1,467
\inspiresetstat{citations}{207,000}  % exact: 207,783
\inspiresetstat{hindex}{211}

%% Per-paper citation counts, keyed by INSPIRE record id.
\inspiresetcites{3159704}{0}
\inspiresetcites{3137624}{2}
\inspiresetcites{3119390}{0}
\inspiresetcites{3103093}{4}
\inspiresetcites{2962374}{9}
\inspiresetcites{2958460}{1}
\inspiresetcites{2874678}{26}
\inspiresetcites{2840007}{27}
\inspiresetcites{2752508}{22}
\inspiresetcites{2690093}{3}
\inspiresetcites{2644916}{13}
\inspiresetcites{2642414}{406}
\inspiresetcites{2628398}{83}
\inspiresetcites{2005570}{23}
\inspiresetcites{1856547}{33}
\inspiresetcites{1831504}{102}
\inspiresetcites{1788448}{119}
\inspiresetcites{1756729}{145}
\inspiresetcites{1710078}{319}
\inspiresetcites{1701002}{245}
\inspiresetcites{1630632}{260}
\inspiresetcites{1609773}{260}
\inspiresetcites{1498566}{50}
\inspiresetcites{1458270}{135}
\inspiresetcites{1345355}{100}
\inspiresetcites{1298030}{152}
\inspiresetcites{1191022}{108}
}{%
  \typeout{^^JLL_CV: LL_InspireData.tex not found; using fallback INSPIRE
figures. Run tools/fetch_inspire.py to refresh them.^^J}%
}


%% ===========================================================================
%% Title block
%% ===========================================================================

% \cvheader -- the name, contact details and ORCID that open every document.
\newcommand{\cvheader}{%
  \begin{minipage}[t]{0.4\textwidth}
    {\huge \name}
  \end{minipage}%
  \begin{minipage}[t]{0.6\textwidth}
    \begin{flushright}
      \href{mailto:\email}{\email}\\
      \phone\\
      \href{https://\homepage}{\homepage}\\
      \href{https://orcid.org/\orcid}{\orcid}
      % A postal address, if one is ever wanted here:
      %   CERN CH-1211\\
      %   B\^atiment 40-1-D11\\
      %   23 Gen\`{e}ve, Switzerland
    \end{flushright}
  \end{minipage}%
}

% \cvfooter -- the "last updated" line that closes every document.
\newcommand{\cvfooter}{%
  \medskip
  \begin{center}
    \begin{small}
      Last updated: \today
    \end{small}
  \end{center}%
}



\begin{document}
\thispagestyle{empty} % no running header on the first page

\cvheader

% Alternative title blocks, kept for reference:
%   \centerline{\huge \bf \name}          % name centred and bold
%   \begin{picture}(0,0)                  % portrait in the top right corner
%     \put(380,-30){\hbox{\includegraphics[height=12em]{Portrait.jpg}}}
%   \end{picture}

\bigskip

% Spacer that holds the vertical gap ahead of the first section.  It pairs with
% the summary paragraph below, which is currently disabled.
\begin{minipage}[t]{1.3em}
  \hspace{1em}
\end{minipage}
% \begin{minipage}[t]{0.75\textwidth}
% \emph{Broadly-trained particle physicist with expertise in searching for physics beyond the standard model at particle colliders. Particular expertise in supersymmetry, long-lived particles, and unconventional searches. Technical skills include expertise in detector data acquisition system, analysis software, and detector simulation.}
% \end{minipage}


\section*{Professional Experience}

  \begin{itemizetight}


    \item Associate Professor, University of Tennessee, Knoxville \when{2026--}
    \item Assistant Professor, University of Tennessee, Knoxville \when{2021--26}

    % \begin{itemizetighter}
    %   \item {University of Tennessee}, Knoxville, TN
    % \end{itemizetighter}

    \item Research Associate, Harvard University - \emph{Supervisor: J. Huth} \when{2016--21}

    % \begin{itemizetighter}
    %   \item {Harvard University}, Cambridge, MA -- \emph{Supervisor: J. Huth}
    % \end{itemizetighter}

    \item Research Associate, University of Adelaide - \emph{Supervisor: P. Jackson} \when{2014--16}

    % \begin{itemizetighter}
    %   \item {The University of Adelaide}, Adelaide, Australia -- \emph{Supervisor: P. Jackson}
    % \end{itemizetighter}

  \end{itemizetight}


\section*{Education}

  \begin{itemizetight}
    \item Ph.D., M.Phil., M.S., Physics, {Yale University} - \emph{Supervisor: T. Golling} \when{2014}
    % \begin{itemizetighter}
    %       \item \emph{Supervisor: T. Golling}
    % \end{itemizetighter}
    % \item M.S., M.Phil. Physics, {Yale University} \when{2012}
    \item B.S. Physics, {Rutgers University} - \emph{Supervisors: R. Ransome, R. Gilman, R. Tumulka} \when{2009}

  %   \begin{itemizetighter}
  %         \item \emph{Supervisors: R. Ransome, R. Gilman, R. Tumulka}
  %   \end{itemizetighter}
  \end{itemizetight}


% \section*{Fields of Research Interest}

%     \begin{itemizetightrightpad}
%     \item High Energy Collider Physics, Experimental searches for physics Beyond the Standard Model, Supersymmetry, Long-Lived Particles, Foundations of Quantum Mechanics
%     % \item Quantum Information, Philosophical Foundations of Quantum Mechanics
%     \end{itemizetightrightpad}


\section*{Awards \textit{\&} Honors}

\emph{External}

  \begin{itemizetighter}
    \item Cottrell Lectureship \when{2026}
    \item Fermilab LHC Physics Center Senior Distinguished Researcher \when{2026}
    \item Breakthrough Prize in Fundamental Physics Laureate  \when{2025}
    \item \hspace{2em}\emph{as part of the ATLAS and CMS Collaborations}
    \item Cottrell Scholar - Research Corporation for Science Advancement \when{2025}
    \item NSF CAREER Award \when{2023}
    \item European Physical Society's High Energy and Particle Physics Prize \when{2013}
    \item \hspace{2em}\emph{as part of the ATLAS Collaboration}
  \end{itemizetighter}


\emph{University of Tennessee, Knoxville}

\begin{itemizetighter}
  \item University of Tennessee Alumni Association -  Alumni Outstanding Teacher Award \when{2025}
  \item UTK Division of Student Success - Faculty Research Mentor Award \when{2025}
  \item UTK Provost - Junior Faculty Fellow \when{2023--25}
  \item UTK Physics - Research Advisor of the Year \when{2023--24}
  \item UTK College of Arts and Sciences - Faculty Academic Outreach Teaching Award \when{2022}
  \item UTK Physics - Teacher of the Year \when{2021--22}
\end{itemizetighter}


\section*{Funding}

\emph{Multi-Year Grants}

  \begin{itemizetighter}
    \item BNL Subcontract on Accelerator Physics - \$190k \when{2025--29}
    \item Cottrell Scholar - Research Corporation for Science Advancement - \$120k \when{2025--28}
    \item DOE Award for UTK HEP (4 PIs) - \$1.425M \when{2024--27}
    \item NSF CAREER Award - \$1M \when{2023--28}
    \item USCMS Upgrade Project Funding - \$73k \when{2022--24}
    \item Single-PI DOE Award - \$250k \when{2022--24}
  \end{itemizetighter}


  \emph{Other Funding}

  \begin{itemizetighter}
    \item Fermilab LHC Physics Center - \$56k\when{2026}
    \item Universities Research Association - \$7k\when{2026}
    \item Fermilab LHC Physics Center (Postdoc Emery Nibigira Led) - \$43k\when{2024}
    \item Universities Research Association (Postdoc Emery Nibigira Led) - \$20k\when{2024}
    \item UT-Battelle (ORNL Management Contractor) - \$74k \when{2024}
    \internalonly{
      \item UTK Graduate Student Senate Travel Award - John Lawless - \$1k\when{2024}
      \item SESAPS Travel Grant - Caroline Riggall - \$325 \when{2024}
      \item SESAPS Travel Grant - Charles Bell - \$500 \when{2023}
      \item SESAPS Travel Grant - Lacey Dishman - \$500 \when{2023}
      \item UTK Faculty Research Assistants Funding - Tanner Wolfe - \$2.3k \when{2023}
      \item FNAL LPC Guest \textit{\&} Visitor Funding - Emery Nibigira - \$5.2k \when{2023}
      \item FNAL LPC Guest \textit{\&} Visitor Funding - John Lawless - \$3.9k \when{2023}
      \item UTK Faculty Research Assistants Funding - Taylor Sussmane - \$2.3k \when{2022}
      \item SESAPS Travel Grant - Taylor Sussmane - \$500 \when{2022}
      \item SESAPS Travel Grant - Alexander Sizemore - \$500 \when{2022}
      \item FNAL LPC Guest \textit{\&} Visitor Funding - Ramon Ogaz - \$2.8k \when{2022}
      \item FNAL LPC Guest \textit{\&} Visitor Funding - Philipp Wagenknecht - \$5.2k \when{2022}
      \item UTK Physics Summer Research Fellowship - Matthew Durkee - \$6k \when{2022}
    }


    \item George Southgate Fellowship, The University of Adelaide - 5k AUD \when{2015}
  \end{itemizetighter}


\section*{Research Experience}

  \begin{itemize}


  \item \textsc{Member of the CMS Collaboration} \when{2021--}


  \begin{changemargin}{0in}{0.0in}
    \begin{itemize}

    \item \textsc{Physics Analysis}

    \vspace{0.1em}
    \begin{itemizetighter}

        \item \textit{Co-convener} LHCC Long-Lived Particle Working Group \when{2022--25}

        \item Heavy Stable Charged Particle Search \when{2022--}

        % \begin{itemizetightrightpad}
        % \item \it Searching for detector-stable particles carrying electric charge and leaving anomalous ionization and timing signatures in the CMS detector.
        % \end{itemizetightrightpad}


        \item Multijet Search \when{2023--}

        % \begin{itemizetightrightpad}
        % \item \it Searching for new physics in multijet final states using conventional techniques and machine learning
        % \end{itemizetightrightpad}

        \item Recursive Jigsaw Reconstruction Search \when{2022}

        % \begin{itemizetightrightpad}
        % \item \it Contributed to calculation of exclusion limits for a search for compressed supersymmetry spectra using the Recursive Jigsaw Reconstruction method.
        % \end{itemizetightrightpad}

    \end{itemizetighter}


    \item \textsc{Phase-II Outer Tracker Upgrade}

    \vspace{0.1em}
    \begin{itemizetighter}

      \item \textit{US CMS Level-3 Manager} for OT Electronics \when{2024--}
      \item \textit{US CMS Control Account Manager} for OT DAQ \when{2024--}

      \item Phase-II Tracker DAQ Software \when{2021--}

      \item Outer Tracker Module Assembly and Grading \when{2022--}

    \end{itemizetighter}


    \end{itemize}
  \end{changemargin}


  \item \textsc{Member of the ATLAS Collaboration} \when{2009--21}

  \begin{changemargin}{0in}{0.0in}
    \begin{itemize}
      % \parshape 2 3.2em \linewidth \listindent \dimexpr\linewidth-\listindent\relax

    % \item { \small [{\it Nota Bene:} Labels in \textsc{small caps} denote official ATLAS designations.] }


    % SUSY Software Review, Contours
    % Maintainer of big framework

    % NSW DAQ

    % \vspace{.5em}

    \item \textsc{Physics Analysis}

    \vspace{0.1em}
    \begin{itemizetighter}

        \item Common BSM Result Interpretation Framework \when{2018--21}

        % \begin{itemizetightrightpad}
        % \item \it Responsible for standards of exclusion interpretation and limit interpolation within the ATLAS SUSY working group, implementing new methods for presentation of results
        % \end{itemizetightrightpad}

        \item \textit{Co-convener} of SUSY RPV/LL sub-group \when{2017--19}

        % \begin{itemizetightrightpad}
        % \item \it Defined standards and direction of searches for long-lived particles and $R$-parity-violating supersymmetry. Served two terms.
        % \end{itemizetightrightpad}

        \item Searches for Long-Lived Particles \when{2016--21}

        % \begin{itemizetightrightpad}
        % \item \it Led many analysis teams in searches for long-lived particles in a variety of signatures. See Publications.
        % \end{itemizetightrightpad}
        % \pagebreak

        \item Common Analysis Software Development \when{2015--21}

        % \begin{itemizetightrightpad}
        % \item \it Primary developer of common analysis software used by roughly ten analysis teams
        % \end{itemizetightrightpad}


        \item Inclusive searches for Supersymmetry \when{2014--16}

        % \begin{itemizetightrightpad}
        % \item \it Led flagship search for supersymmetry in inclusive zero-lepton final states in 2016 including leading the first searches for new physics using ``Recursive Jigsaw'' variables. Created and commissioned novel triggers for use during Run-2 of the LHC geared toward finding new physics in low-mass-splitting scenarios.
        % \end{itemizetightrightpad}


        \item All-Hadronic BSM Signatures \when{2011--14}

        % \begin{itemizetightrightpad}
        % \item \it Led two searches for new phenomena in the context of $R$-Parity violating supersymmetry in fully hadronic final states
        % \end{itemizetightrightpad}


        \item Quark-vs-Gluon Jet Tagging \when{2010--12}

        % \begin{itemizetightrightpad}
        % \item \it Developed the first quark- vs. gluon-jet discriminator at a hadron collider
        % \end{itemizetightrightpad}

    \end{itemizetighter}

    \item \textsc{New Small Wheel Muon Spectrometer Upgrade}

    \vspace{0.1em}
    \begin{itemizetighter}

      \item \textit{Coordinator} Micromegas Trigger \when{2020}

      % \begin{itemizetightrightpad}
      % \item \it Integration of trigger electronics for the micromegas detector of the New Small Wheel (NSW) upgrade for the ATLAS muon spectrometer
      % \end{itemizetightrightpad}

      \item Online Software Coordination \when{2019--20} % \hspace{0.9em}}

      % \begin{itemizetightrightpad}
      % \item \it Designed, implemented, and/or maintain the readout, configuration, and calibration software systems for the NSW
      % \end{itemizetightrightpad}

      % \item Micromegas Digitization \when{2017--18}

      % \begin{itemizetightrightpad}
      % \item \it Responsible for digitization and simulation of the readout and trigger paths for the micromegas detector for the NSW
      % \end{itemizetightrightpad}

    \end{itemizetighter}


    \end{itemize}
  \end{changemargin}

  \item \textsc{Muon Collider Projects} \when{2020--}

  \begin{changemargin}{0in}{0.0in}
    \begin{itemizetight}

      \item \emph{USMCC Deputy Communication Coordinator} \when{2025--}
      \item 10 TeV MAIA Detector Concept Characterization \when{2023--}
      \item \emph{Taskforce Member} IMCC Software Taskforce \when{2024}

      \item Automating Professional 3D Rendering for Collider Event Displays \when{2023--}

      \item Simulation Software Workflow Development \when{2020--22}

    \end{itemizetight}
  \end{changemargin}


  \item \textsc{Accelerator Physics Projects} \when{2024--}

  \begin{changemargin}{0in}{0.0in}
    \begin{itemizetight}

    \item Helical FOFO Design for Muon Cooling for a Future Muon Collider \when{2024--}

    \item Using Self-Consistent Beam Distributions to Reduce Intra-Beam Stripping \when{2024--}

    \end{itemizetight}
  \end{changemargin}


  \item \textsc{External Collider Physics}


    \begin{itemizetight}

    \item Experimental Impact of Jet Fragmentation Reference Frames At Particle Colliders \when{2022--}


    \item Future collider projections for photon-induced production of charged BSM particles \when{2018--21}

    \item Bell Inequalities at Future $ee$ Colliders \when{2013--}

    \end{itemizetight}


  \item \textsc{Nucleon Physics Research} \when{2007--09}

  % \item Developed the prototype trigger and data acquisition systems for Fermilab E-906/SeaQuest \when{2008--2009}

  \begin{changemargin}{0in}{0.0in}
    \begin{itemizetightrightpad}
            \parshape 2 3.2em \linewidth \listindent \dimexpr\linewidth-\listindent\relax

    \item Developed the prototype trigger and data acquisition systems for Fermilab E-906/SeaQuest\vspace{0.2em}
    \item Assembled the photomultiplier tube units used in MINERvA\vspace{0.2em}
    \item Upgraded a Fabry-Perot cavity for use in a JLAB Compton polarimeter\vspace{0.2em}
    \end{itemizetightrightpad}
  \end{changemargin}

\end{itemize}


% \pagebreak

\section*{Publications}

  \begin{itemizetight}

  %!TEX root = LL_CV.tex
%%
%% The publication list, shared by LL_CV.tex and LL_Publications.tex.
%%
%% Titles are written with the commands defined in LL_Preamble.tex:
%%
%%   \inspirepub{<record>}{<title>}  a title linked to its INSPIRE-HEP record,
%%                                   annotated with that record's live citation
%%                                   count (see tools/fetch_inspire.py)
%%   \pub{<url>}{<title>}            a title linked anywhere else: an arXiv-only
%%                                   preprint, a CDS or CERN note, a bare DOI
%%   \inspirestat{<key>}             a live author-level figure used in the
%%                                   summary paragraph below
%%   \whenflush{<date>}              the date, flush right
%%
%% $\RHD$ marks a primary editor, primary analyzer, or intellectual lead role.

  \item  \href{https://inspirehep.net/authors/1071846}{Over \inspirestat{papers} published papers, mostly as part of the ATLAS and CMS Collaborations, with over \inspirestat{citations} citations, and an $h$-index of \inspirestat{hindex} [\textit{INSPIRE}].} Valid authorship in the ATLAS and CMS Collaborations requires a record of and continued contributions to the operation, maintenance, and improvements of the experiments in service of all of the collaboration's activities. Selected publications with significant direct contributions are listed below. \textit{{$\RHD$} indicates a primary editor, primary analyzer, or intellectual lead role.}


  \subsection*{Small Author List Publications}
  \begin{enumerate}

    \item S. Kakkar, L. Lee, N. Evans, A. Hoover, V. Morozov, \inspirepub{3119390}{Using Normalizing Flows in Normal-Form Space to Model Intra-Beam Stripping}, \emph{JACoW IPAC26} WEP4348 \newline \emph{(Proceedings, norm for publication in Accelerator Physics)} \whenflush{2026}

    \item $\RHD$ L.~Lee, J.~Lawless, C.~Riggall, \inspirepub{3137624}{Higgs Boson Spookiness: Probing Quantum Nonlocality with Spacetime-Resolved $H\!\to\!\tau^+\tau^-$ Decays}, \emph{arXiv}, 2603.28868 \whenflush{2026}

    \item $\RHD$ L.~Jeanty, L.~Lee, \inspirepub{3103093}{Rare and Experimentally Challenging Supersymmetry Signatures}, \emph{arXiv}, 2601.06358, \emph{Invited, Accepted by Annual Reviews} \whenflush{2026}

    \item C.~Riggall, R.~Ganesan, L.~Lee, T.~Holmes, K.~Yonehara, D.~Stratakis, D.~Neuffer, M.~Palmer, \inspirepub{3159704}{A Charge-Agnostic Design For 6D Muon Ionization Cooling}, \emph{JACoW COOL2025} \newline \emph{(Proceedings, norm for publication in Accelerator Physics)} \whenflush{2025}

    \item S. Kakkar, V. Morozov, L. Lee, N. Evans, A. Zhukov, A. Hoover, \inspirepub{3119390}{Beam Loss Modeling And Mitigation Due To
Intra-Beam Stripping}, \emph{JACoW NAPAC2025} THP014 \newline \emph{(Proceedings, norm for publication in Accelerator Physics)} \whenflush{2025}

    \item $\RHD$ T. Holmes, L.~Lee, \inspirepub{2958460}{Event Rates at a 10 TeV Muon Collider and Implications for Detector Design: Trigger, Data Acquisition, and Luminosity}, \emph{arXiv}, 2508.06239  \whenflush{2025}

    \begin{itemizetightrightpad}
      \item \emph{Discusses the experimental implications for electroweak process cross sections, particularly with respect to the collision rate and background process rates.}
    \end{itemizetightrightpad}

    \item $\RHD$ L.~Lee, C.~Bell, J.~Lawless, C.~Nash, E.~Nibigira, \inspirepub{2690093}{Experimental Impact of Jet Fragmentation Reference Frames At Particle Colliders}, \emph{PLB}, Volume 866, 139561  \whenflush{2025}

    \begin{itemizetightrightpad}
      \item \emph{Intellectual lead in argument and analysis demonstrating how modern jet physics relies on a frame-specific assumption that can fail in interesting corners of phase space.}
    \end{itemizetightrightpad}

  \item  C.~Bell, et al (\emph{21 others}) \inspirepub{2874678}{MAIA: A new detector concept for a 10 TeV muon collider},
  \\ \emph{arXiv:2502.00181}  \whenflush{2025}

  \begin{itemizetightrightpad}
    \item \emph{Main editor of Section VI. Responsible for all 3D renderings of the detector concept. General input on the overall physics output.}
  \end{itemizetightrightpad}


  \item C.~Accettura, et al (\emph{296 others}), \inspirepub{2642414}{Towards a Muon Collider}, \emph{Eur. Phys. J. C} 83, 864  \whenflush{2023}

  \begin{itemizetightrightpad}
    \item \emph{Contributed studies of Beam-Induced Background mitigation at a future muon collider.}
  \end{itemizetightrightpad}

  \item G.~Iakovidis, et al (\emph{100 others}), \inspirepub{2644916}{The New Small Wheel electronics}, \emph{JINST} 18 P05012 \whenflush{2023}

  \begin{itemizetightrightpad}
    \item \emph{Contributed to readout electronics integration with readout software and FELIX, integration of trigger hardware readout and dataflow, development of control and calibration software, and noise studies.}
  \end{itemizetightrightpad}

  \item $\RHD$ A.~Badea, W.~Fawcett, J.~Huth, T.J.~Khoo, R.~Poggi, L.~Lee, \inspirepub{2005570}{Solving Combinatorial Problems at Particle Colliders Using Machine Learning}, \emph{PRD} 106 1, 016001 \whenflush{2022}

  \begin{itemizetightrightpad}
    \item \emph{Designed physics-informed neural network to solve combinatorial ambiguities in collider events with many objects, especially for decays from pair-produced heavy resonances.}
  \end{itemizetightrightpad}

  \item C.~Bakalis, I.~Grayzman, L.~Lee, L.~Levinson, V.~Polychronakos, \pub{https://doi.org/10.1088/1748-0221/16/06/P06041}{Accessing register spaces in FPGAs within the ATLAS DAQ scheme via the SCA eXtension}, \emph{JINST} 16 P06041\whenflush{2021}

  \begin{itemizetightrightpad}
    \item \emph{Designed the software infrastructure around an FPGA implementation of a CERN GBT-SCA-like register interface for use in detector configuration.}
  \end{itemizetightrightpad}

  \item M.~Bauer, O.~Brandt, L.~Lee, C.~Ohm, \inspirepub{1756729}{ANUBIS: AN Underground Belayed In-Shaft search experiment}, \href{https://arxiv.org/abs/1909.13022}{arXiv: 1909.13022}\whenflush{2020}

  \begin{itemizetightrightpad}
    \item \emph{Responsible for original idea to instrument service shafts above the ATLAS Experiment for a dedicated long-lived particle experiment. Helped shape technical proposal outlined in paper.}
  \end{itemizetightrightpad}


  \item $\RHD$ L.~Lee, C.~Ohm, A.~Soffer, and T.~Yu, \inspirepub{1701002}{Collider Searches for Long-Lived Particles Beyond the Standard Model}, \emph{Progress in Particle and Nuclear Physics} 3695\whenflush{2019}


  \begin{itemizetightrightpad}
    \item \emph{Co-responsible for all experimental aspects of this broad review of collider searches for long-lived particles.}
  \end{itemizetightrightpad}


  \item X. Cid Vidal, et al., \inspirepub{1710078}{Beyond the Standard Model Physics at the HL-LHC and HE-LHC}, CERN Yellow Rep. Monogr. 7 (2019) 585-865 \whenflush{2019}

  \begin{itemizetightrightpad}
    \item \emph{Contributed studies of HL-LHC projections for long-lived gluino searches. Responsible for sensitivity projections.}
  \end{itemizetightrightpad}


  \end{enumerate}


  \subsection*{Selected Peer-Reviewed Publications from the ATLAS and CMS Collaborations}
  \begin{enumerate}[resume]
  \item $\RHD$ CMS Collaboration, \inspirepub{2962374}{A general search for supersymmetric particles in scenarios with compressed mass spectra using proton-proton collisions at $\sqrt{s}$ = 13 TeV}, \emph{Phys. Rev. D} 112 (2025) 11, 112023 \whenflush{2025}

  \begin{itemizetightrightpad}
    \item \emph{Limit setting and model interpretation.}
  \end{itemizetightrightpad}

  \item $\RHD$ CMS Collaboration, \inspirepub{2840007}{Search for heavy long-lived charged particles with large ionization energy loss in proton-proton collisions at $\sqrt{s}$ = 13 TeV}, \emph{JHEP} 04 (2025) 109  \whenflush{2025}

  \begin{itemizetightrightpad}
    \item \emph{Led multiple aspects of this analysis, particularly in ionization analysis design, simulation, and signal MC generation.}
  \end{itemizetightrightpad}

  \item $\RHD$ ATLAS Collaboration, \inspirepub{2752508}{A search for R-parity-violating supersymmetry in final states containing many jets in $pp$ collisions at $\sqrt{s} = 13$~TeV with the ATLAS detector}, \emph{JHEP} 05 (2024) 003  \whenflush{2024}

  \begin{itemizetightrightpad}
    \item \emph{Early designer of multijet analysis, software developer, and started the machine learning effort portion of this search.}
  \end{itemizetightrightpad}


 \item $\RHD$ ATLAS Collaboration, \inspirepub{2628398}{Search for long-lived, massive particles in events with displaced vertices and multiple jets in pp collisions at $\sqrt{s}=13$~TeV with the ATLAS detector}, \emph{JHEP} 06, 200 \whenflush{2023}

  \begin{itemizetightrightpad}
    \item \emph{Led initial design and implementation of a search for decays of BSM particles with displaced vertices and jets, including reconstruction algorithms, analysis software, and background estimation design.}
  \end{itemizetightrightpad}


 \item $\RHD$ ATLAS Collaboration, \inspirepub{1856547}{A search for the decays of stopped long-lived particles at $\sqrt{s}=13$~TeV with the ATLAS detector}, \emph{JHEP} 173 \whenflush{2021}

  \begin{itemizetightrightpad}
    \item \emph{Led search for very long lived particle decays in empty LHC bunch crossings.}
  \end{itemizetightrightpad}

  \item ATLAS Collaboration, \inspirepub{1831504}{Search for displaced leptons in $\sqrt{s} = 13$ TeV $pp$ collisions with the ATLAS detector}, \emph{Phys. Rev. Lett.} 127 051802 \whenflush{2021}

  \begin{itemizetightrightpad}
    \item \emph{Developed analysis software and provided design guidance for a search that extends the limits obtained from LEP for long-lived sleptons.}
  \end{itemizetightrightpad}


  \item $\RHD$ ATLAS Collaboration, \inspirepub{1788448}{Search for long-lived, massive particles in events with a displaced vertex and a muon with large impact parameter in $pp$ collisions at $\sqrt{s}=13$~TeV with the ATLAS detector}, \emph{Phys. Rev. D} 102 032006\whenflush{2020}

  \begin{itemizetightrightpad}
    \item \emph{Co-coordinator of analysis. Responsible for all aspects of the analysis, including background estimation scheme, physics targets, and analysis software.}
  \end{itemizetightrightpad}

  \item $\RHD$ ATLAS Collaboration, \inspirepub{1630632}{Search for long-lived, massive particles in events with displaced vertices and missing transverse momentum in 13~TeV pp collisions with the ATLAS detector}, \emph{Phys. Rev. D} 97 052012\whenflush{2018}

  \begin{itemizetightrightpad}
    \item \emph{Co-coordinator of analysis. Particularly responsible for background estimation methods and their validation, limit-setting, and data-flow model.}
  \end{itemizetightrightpad}

  \item ATLAS Collaboration, \inspirepub{1609773}{Search for new phenomena in high-mass diphoton final states using $37$~fb$^{-1}$ of proton-proton collisions collected at $\sqrt{s} = 13$~TeV with the ATLAS detector}, \emph{Phys. Lett. B} 775 105\whenflush{2017}

  \begin{itemizetightrightpad}
    \item \emph{Responsible for creation and validation of BSM signals and Higgs EFT models.}
  \end{itemizetightrightpad}

  \item ATLAS Collaboration, \inspirepub{1498566}{Search for new phenomena in events containing a same-flavour opposite-sign dilepton pair, jets, and large missing transverse momentum in $\sqrt{s} = 13$~TeV pp collisions with the ATLAS detector}, \emph{Eur. Phys. J C} \textbf{77} 144\whenflush{2017}

  \begin{itemizetightrightpad}
    \item \emph{Responsible for studies of the jet system dynamics.}
  \end{itemizetightrightpad}


    % \pagebreak

    % \item {\textit{\hspace{-1.6em} [Publications continued...] }}\vspace{-0.5em}

  \item $\RHD$ ATLAS Collaboration, \inspirepub{1458270}{Search for squarks and gluinos in final states with jets and missing transverse momentum at $\sqrt{s} = 13$~TeV with the ATLAS detector}, \emph{Eur. Phys. J. C} \textbf{76}: 392\whenflush{2016}

  \begin{itemizetightrightpad}
    \item \emph{Co-coordinator of search. Responsible for all aspects of limit-setting. Responsible for background estimation design and implementation for RJR analysis.}
  \end{itemizetightrightpad}


  \item $\RHD$ ATLAS Collaboration,
    \inspirepub{1345355}{Search for massive supersymmetric particles decaying to many jets using the ATLAS detector in pp collisions at $\sqrt{s} = 8$~TeV},
    \emph{Phys. Rev. D} \textbf{91}, 112016\whenflush{2015}

  \begin{itemizetightrightpad}
    \item \emph{Co-coordinator of search. Responsible for all aspects of the search from background estimation, signal simulation, and limit-setting.}
  \end{itemizetightrightpad}


  \item $\RHD$ ATLAS Collaboration, \inspirepub{1298030}{Light-quark and Gluon Jet Discrimination in pp Collisions at $\sqrt{s} = 7$~TeV with the ATLAS Detector},
    \emph{Eur. Phys. J. C} \textbf{74}: 3023\whenflush{2014}

  \begin{itemizetightrightpad}
    \item \emph{Responsible for analysis design, data-driven template determination and validation.}
  \end{itemizetightrightpad}

  \item $\RHD$ ATLAS Collaboration, \inspirepub{1191022}{Search for pair production of massive particles decaying into three quarks with the ATLAS detector in $\sqrt{s} = 7$~TeV pp collisions at the LHC},~Journal of High Energy Physics 12, 1--42\whenflush{2012}


  \begin{itemizetightrightpad}
    \item \emph{Responsible for all signal simulation, limit-setting, and interpretation.}
  \end{itemizetightrightpad}

  \end{enumerate}


  \subsection*{Selected Public Documents}

  \begin{itemize}


  \item A. Affolder, et al., \pub{https://arxiv.org/abs/2209.03607}{Solid State Detectors and Tracking for Snowmass},
    \emph{Snowmass 2021}, \whenflush{2022}

  \item $\RHD$ D. Ally, et al., \pub{https://arxiv.org/abs/2203.06773}{Strategies for Beam-Induced Background Reduction at Muon Colliders},
    \emph{Snowmass 2021}, \whenflush{2022}

  \item $\RHD$ J. Cochran, et al., \pub{https://arxiv.org/abs/2203.09585}{Particle Physics Outreach at Non-traditional Venues},
    \emph{Snowmass 2021}, \whenflush{2022}

  \item K. Black, et al., \pub{https://arxiv.org/abs/2203.08874}{Prospects for the Measurement of the Standard Model Higgs Pair Production at the Muon Colliders}, \emph{Snowmass 2021}, \whenflush{2022}

  \item N. Bartosik, et al., \pub{https://arxiv.org/abs/2203.07964}{Simulated Detector Performance at the Muon Collider}, \emph{Snowmass 2021}, \whenflush{2022}

  \item J. De Blas, et al., \pub{https://arxiv.org/abs/2203.07261}{The physics case of a 3 TeV muon collider stage}, \emph{Snowmass 2021}, \whenflush{2022}

  \item C. Aime, et al., \pub{https://arxiv.org/abs/2203.07256}{Muon Collider Physics Summary},
    \emph{Snowmass 2021}, \whenflush{2022}

  \item S. Jindariani, et al., \pub{https://arxiv.org/abs/2203.07224}{Promising Technologies and R\&D Directions for the Future Muon Collider Detectors}, \emph{Snowmass 2021}, \whenflush{2022}


  \item ATLAS Collaboration, \pub{https://atlas.web.cern.ch/Atlas/GROUPS/PHYSICS/PUBNOTES/ATL-PHYS-PUB-2019-019/}{Generation and Simulation of $R$-Hadrons},~ATL-PHYS-PUB-2019-019\whenflush{2019}
  % \item $\RHD$ {\it\href{https://atlas.web.cern.ch/Atlas/GROUPS/PHYSICS/CONFNOTES/ATLAS-CONF-2019-006/}{Search for long-lived, massive particles in events with a displaced vertex and a displaced muon in $pp$ collisions at $\sqrt{s}=13$~TeV with the ATLAS detector}},~ATLAS-CONF-2019-006 (2019)
  \item $\RHD$ ATLAS Collaboration, \pub{https://atlas.web.cern.ch/Atlas/GROUPS/PHYSICS/PUBNOTES/ATL-PHYS-PUB-2019-012/}{Supersymmetry Summary Plots}, Main overview and RPV plots \whenflush{2018--20}
  \item ATLAS Collaboration, \pub{https://atlas.web.cern.ch/Atlas/GROUPS/PHYSICS/PUBNOTES/ATL-PHYS-PUB-2019-013/}{Performance of vertex reconstruction algorithms for detection of new long-lived particle decays within the ATLAS inner detector},~ATL-PHYS-PUB-2019-013\whenflush{2019}
  % \item {\it\href{https://atlas.web.cern.ch/Atlas/GROUPS/PHYSICS/PUBNOTES/ATL-PHYS-PUB-2018-033/}{Sensitivity of the ATLAS experiment to long-lived particles with a displaced vertex and $E_T^{\mbox{miss}}$ signature at the HL-LHC}},~ATL-PHYS-PUB-2018-033 (2018)
  \item ATLAS Collaboration, \pub{https://atlas.web.cern.ch/Atlas/GROUPS/PHYSICS/CONFNOTES/ATLAS-CONF-2018-003/}{Reinterpretation of searches for supersymmetry in models with variable R-parity-violating coupling strength and long-lived R-hadrons},~ATLAS-CONF-2018-003\whenflush{2018}
  % \item $\RHD$ {\it\href{https://atlas.web.cern.ch/Atlas/GROUPS/PHYSICS/CONFNOTES/ATLAS-CONF-2017-026/}{Search for long-lived, massive particles in events with displaced vertices and missing transverse momentum in 13~TeV pp collisions with the ATLAS detector}},~ATLAS-CONF-2017-026 (2017)
  % \item {\it\href{https://atlas.web.cern.ch/Atlas/GROUPS/PHYSICS/CONFNOTES/ATLAS-CONF-2017-022/}{Search for squarks and gluinos in final states with jets and missing transverse momentum using $36$~fb$^{-1}$ of $\sqrt{s} = 13$~TeV pp collision data with the ATLAS detector.}}~ATLAS-CONF-2017-022 (2017)
  % \item $\RHD$ {\it\href{https://atlas.web.cern.ch/Atlas/GROUPS/PHYSICS/CONFNOTES/ATLAS-CONF-2016-078/}{Further searches for squarks and gluinos in final states with jets and missing transverse momentum at $\sqrt{s} = 13$~TeV with the ATLAS detector}},~ATLAS-CONF-2016-078 (2016)
  % \item {\it\href{https://atlas.web.cern.ch/Atlas/GROUPS/PHYSICS/CONFNOTES/ATLAS-CONF-2016-059/}{Search for scalar diphoton resonances with $15.4$~fb$^{-1}$ of data collected at $\sqrt{s} = 13$~TeV in 2015 and 2016 with the ATLAS detector}},~ATLAS-CONF-2016-059 (2016)
  % \item {\it\href{https://cds.cern.ch/record/2114828}{Search for squarks and gluinos in final states with jets and missing transverse momentum at $\sqrt{s} = 13$~TeV with the ATLAS detector}},~ATLAS-CONF-2015-062 (2015)
  \item ATLAS Collaboration, \pub{https://cds.cern.ch/record/2037905}{First look at pp collisions data at $\sqrt{s} = 13$~TeV in preparation for a search for squarks and gluinos in final states with jets and missing transverse momentum with the ATLAS detector},~ATL-PHYS-PUB-2015-028\whenflush{2015} \vspace{-1em}
  % \item $\RHD$ {\it\href{https://cds.cern.ch/record/1558689}{Search for massive particles decaying into multiple quarks with the ATLAS detector in $\sqrt{s} = 8$~TeV pp collisions}},~ATLAS-CONF-2013-091 (2013)
  \vspace{1em}
  \item \textit{Additional public documents supporting the publications described above.}

  \end{itemize}


  \subsection*{One CMS Analysis Review Committee (ARC)}

  \subsection*{Nine ATLAS Editorial Boards (EBs)}

  % \subsection*{Journal Referee for Physical Review Letters, Physical Review D, Letters in High Energy Physics}

  % \begin{itemize}

  % \item Served on 8 ATLAS Editorial Boards

  % % \item {\it\href{https://atlas.web.cern.ch/Atlas/GROUPS/PHYSICS/CONFNOTES/ATLAS-CONF-2019-027/}{Soft $b$-hadron tagging for compressed SUSY scenarios}},~ATLAS-CONF-2019-027 (2019)

  % % \item {\it Search for new phenomena in a lepton plus high jet multiplicity final state with the ATLAS experiment using $\sqrt{s}$ = 13~TeV proton-proton collision data},~ATLAS-CONF-2016-094, ATLAS-CONF-2017-013, JHEP 09 (2017) 88

  % % \item {\it A search for pair-produced resonances in four jets final states at $\sqrt{s}=13$~TeV with the ATLAS detector},~ATLAS-CONF-2016-022 (2016)

  % % \item {\it\href{https://cds.cern.ch/record/2152392}{A search for R-parity violating decays of the top squark in four jet final states with the ATLAS detector at $\sqrt{s}=13$~TeV}},~ATLAS-CONF-2016-022 (2016)

  % % \item {\it\href{http://link.springer.com/article/10.1140/epjc/s10052-016-4065-1}{A new method to distinguish hadronically decaying boosted Z bosons from W bosons using the ATLAS detector}},~Eur. Phys. J. C \textbf{76} 238 (2016)

  % % \item {\it\href{https://cds.cern.ch/record/2017303}{Constraints on promptly decaying supersymmetric particles with lepton-number- and R-parity-violating interactions using Run-1 ATLAS data}},~ATLAS-CONF-2015-018 (2015)

  % \end{itemize}


  % \subsection*{Other Unpublished Work}

  % \begin{itemize}

  % % FQM Stuff
  % \item $\RHD$ GRW Theory and the Thermodynamic Arrow of Time (2010)
  % \item $\RHD$ Relativistic Considerations of the GRW Theory \emph{w/ R. Tumulka} (2008)
  % \item $\RHD$ Trouble with Locality: Adding Nonlocal Correlations \emph{w/ R. Tumulka} (2008)
  % \item $\RHD$ GRW Theory \emph{w/ R. Tumulka} (2008)

  % \end{itemize}


\end{itemizetight}

% \newpage

% \section*{Grants, Fellowships, \& Awards}

%   \begin{itemize}
%   \item George Southgate Fellowship, The University of Adelaide, 2015
%   \item 2nd International Spring School on Particle Physics and Philosophy Travel Award, 2014
%   \item AAAS/Science Program for Excellence in Science, 2013
%   % \item New York Academy of Science Sponsored Membership, 2012
%   \item Yale University Graduate Student Assembly Conference Travel Fund, 2012
%   % \item NSF Graduate Research Fellowship, Honorable Mention, 2011, 2010
%   \item Division of Nuclear Physics, American Physical Society Fall Meeting, Conference Travel Fund, 2008
%   % \item Summer Research Fellowship, Jefferson Laboratory, Rutgers University, 2008
%   \end{itemize}


\section*{Professional Activities}


  \begin{itemizetighter}
    \internalonly{
      \item \emph{Journal Referee:} PRD \when{2025}
    }
    \item RCSA Reviewer \when{2026} % For Cottrell
    \item Lead organizer, $\mu$CATALYST summer accelerator program \when{2026}
    \item Chair of Fermilab \emph{Colliders of Tomorrow} Committee \when{2025--}
    \item US LHC Users Association (USLUA) Executive Committee Member \when{2022--25}
    \item NSF Reviewer \when{2025} % For Kathy
    \item NSF Review Panel \when{2025}
    \item DOE Reviewer \when{2025}
    \item USCMS Phase-II Upgrade Level-3 Manager for Outer Tracker DAQ \when{2024--}
    \internalonly{
      \item \emph{Journal Referee:} PRD \when{2024}
      \item \emph{Journal Referee:} EPJC \when{2024}
    }

    \item Swiss National Science Foundation Reviewer \when{2024}
    \item International Muon Collider Collaboration Software Task Force \when{2024}
    \item US Muon Collider Collaboration Membership Coordinator \when{2024}
    \item US CMS (DOE+NSF) Project Review Panel \when{2024}
    \item NSF Reviewer ($\times$3 occasions) \when{2024}
    \item NSF Review Panel \when{2024}
    \internalonly{
      \item \emph{Journal Referee:} PRD \when{2024}
    }

    \item NSF Reviewer \when{2023}
    \item DOE Reviewer \when{2023}
    \item LHC LLP Working Group Convener \when{2022--25}
    \item SESAPS Nominating Committee \when{2022--23}
    \item Physics Today Consultant \when{2022}
    \internalonly{
      \item \emph{Journal Referee:} Universe \when{2022}
      \item \emph{Journal Referee:} Universe \when{2022}
      \item \emph{Journal Referee:} Instruments \when{2021}
      \item \emph{Journal Referee:} PRD \when{2021}
      \item \emph{Journal Referee:} LHEP \when{2021}
      \item \emph{Journal Referee:} PRL \when{2020}
      \item \emph{Journal Referee:} PRL \when{2018}
    }

    \externalonly{
      \item Journal Referee for Physical Review Letters, Physical Review D, \when{2018--}
      \item \hspace{2em} EPJC, Instruments, Letters in High Energy Physics, Universe
    }
  \end{itemizetighter}

\section*{Selected Scientific Presentations}

  \emph{Invited Seminars and Colloquiua}


  % \begin{itemize}
  \begin{itemizetighter}


    % \item University of Florida -- \textit{Physics Colloquium, Forthcoming} \when{2026}

    \item UC San Diego -- \textit{Physics Colloquium, Cottrell Lecture} \when{2026} % May 7, 2026

    \item San Diego State University -- \textit{Physics Colloquium, Cottrell Lecture} \when{2026} % May 8, 2026

    \item Cornell University -- \textit{HEP Seminar} \when{2026} % April 24, 2026

    \item Fermilab -- \textit{Wine \& Cheese Seminar} \when{2026} % 16 Jan 2026

    \item Yale University -- \textit{NPA Seminar} \when{2024} % 9 May 2024

    \item New York University -- \textit{Physics Colloquium} \when{2024} % March 2024

    \item University of California, Irvine, Virtual -- \textit{AI Seminar} \when{2022}
    % July 12, 2022

    \item University of Kansas, Lawrence, KS -- \emph{Particle Physics on the Plains} \when{2022}
    % March 8, 2022

    \item University of Tennessee, Knoxville, TN -- \emph{HEP Seminar} \when{2021}%
    % April 21, 2021
    % \item \textbf{R-Parity Violating SUSY is Hard to Find}
    %   \begin{itemizetighter}
    %   \item University of Tennessee, \emph{Seminar}, Knoxville, TN \when{2021}%
    %   \end{itemizetighter}

    \item University of Tennessee, Knoxville, TN -- \emph{Physics Colloquium} \when{2020}%
    %  Nov 23, 2020
    % \item \textbf{What To Do When Nature Is Unnatural}
    %   \begin{itemizetighter}
    %   \item University of Tennessee, \emph{Physics Colloquium}, Knoxville, TN \when{2020} % 11 December 2018
    %   \end{itemizetighter}

    \item Fermilab, Batavia, IL -- \emph{LPC Physics Forum} \when{2020}
    % 19 Nov 2020
    % \item \textbf{Displaced Vertex Searches at ATLAS}
    %   \begin{itemizetighter}
    %   \item Fermilab, \emph{LPC Physics Forum}, Batavia, IL \when{2020} % 11 December 2018
    %   \end{itemizetighter}


    \item Lund University, Sweden -- \emph{Seminar} \when{2020} % Mar 5, 2020
    \item Michigan State University, East Lansing, MI -- \emph{Physics Colloquium} \when{2019} % 5 February 2019
    \item University of Illinois, Urbana-Champaign, IL -- \emph{Physics Colloquium} \when{2018} % 25 October 2018
    \item Fermilab, Batavia, IL --  \emph{Topic of the Week} \when{2018} % 22 October 2018
    \item SLAC, Palo Alto, CA -- \emph{Seminar} \when{2018} % 18 October 2018
    \item University of Oregon, Eugene, OR -- \emph{Seminar} \when{2018} % 15 October 2018
    \item University of Pennsylvania, Philadelphia, PA -- \emph{Seminar} \when{2018} % 7 August 2018
    \item Rutgers University, Piscataway, NJ -- \emph{Seminar} \when{2018} % 6 August 2018

    % \item \textbf{How SUSY Can Still Save The Day / The SUSY Swindle}
    %   \begin{itemizetighter}
    %   \item Lund University, Sweden \when{2020}
    %   \item Michigan State University, East Lansing, MI \when{2019} % 5 February 2019
    %   \item University of Illinois, Urbana-Champaign, IL \when{2018} % 25 October 2018
    %   \item Fermilab \emph{Topic of the Week}, Batavia, IL \when{2018} % 22 October 2018
    %   \item SLAC, Palo Alto, CA \when{2018} % 18 October 2018
    %   \item University of Oregon, Eugene, OR \when{2018} % 15 October 2018
    %   \item University of Pennsylvania, Philadelphia, PA \when{2018} % 7 August 2018
    %   \item Rutgers University, Piscataway, NJ \when{2018} % 6 August 2018
    %   \end{itemizetighter}


    \item Harvard University, Cambridge, MA -- \emph{LPPC Seminar} \when{2018} % 11 October 2018
    % \item \textbf{On the Higgs Branching Ratios}
    %   \begin{itemizetighter}
    %   \item Harvard University - LPPC Seminar, Cambridge, MA\when{2018} % 11 October 2018
    %   \end{itemizetighter}


    \item Oskar Klein Centre, Stockholm, Sweden -- \emph{HEP Seminar}\when{2018} % 22 May 2018
    \item Lawrence Berkeley National Laboratory, Berkeley, CA  -- \emph{RPM Seminar}\when{2017} % 17 August 2017
    % \item \textbf{Searches for Sneaky SUSY at the ATLAS Experiment \& A Toy Entropic Explanation for the Hierarchy Problem}
    %   \begin{itemizetighter}
    %   \item Oskar Klein Centre, Stockholm, Sweden \when{2018} % 22 May 2018
    %   \item Lawrence Berkeley National Laboratory, Berkeley, CA \when{2017} % 17 August 2017
    %   \end{itemizetighter}

    \item Harvard University, Cambridge, MA -- \emph{LPPC Seminar}\when{2015} % 30 September 2015
  % \item \textbf{Searches for R-Parity Violating SUSY at the ATLAS Experiment}
  %   \begin{itemizetighter}
  %   \item Harvard University - LPPC Seminar, Cambridge, MA\when{2015}% 30 September 2015
  %   \end{itemizetighter}

    \item University of Adelaide, Australia -- \emph{CSSM Seminar}\when{2014} % 13 August 2014
    \item University of Melbourne, Australia -- \emph{Seminar}\when{2014} % 10 July 2014
    \item Institut de F\'isica d'Altes Energies (IFAE), Barcelona, Spain -- \emph{Seminar}\when{2014} % 7 January 2014

  % \item \textbf{A Search for B-Violating Supersymmetry in Multijet Signatures and Light-quark vs. Gluon Jet Tagging}
  %   \begin{itemizetighter}
  %   \item University of Adelaide, Australia\when{2014}% 13 August 2014
  %   \item University of Melbourne, Australia\when{2014}% 10 July 2014
  %   \item Institut de F\'isica d'Altes Energies (IFAE), Barcelona, Spain\when{2014} % 7 January 2014
  %   \item Santa Cruz Institute of Particle Physics Seminar, Geneva, Switzerland\when{2013}% 12
  %   \end{itemizetighter}

    \item 2nd International Spring School on Particle Physics and Philosophy -- \emph{Seminar}\when{2014}
  % \item \textbf{...But What If I Like Naturalness?}
  %   \begin{itemizetighter}
  %   \item 2nd International Spring School on Particle Physics and Philosophy
  %   \item Wuppertal, Germany\when{2014}% 31 March 2014
  %   \end{itemizetighter}

    \item Santa Cruz Institute of Particle Physics Seminar, Geneva, Switzerland -- \emph{Seminar}\when{2013} % 12

  % \item \textbf{RPV Stops at the LHC}
  %   \begin{itemizetighter}
  %   \item LPC Workshop on Exotic Top Partners
  %   \item Fermi National Accelerator Laboratory, LHC Physics Center, Batavia, IL\when{2013}% 26 September 2013
  %   \end{itemizetighter}


% \end{itemize}
\end{itemizetighter}

  \emph{Conference Presentations}

  % \begin{itemize}
  \begin{itemizetighter}

  \item DPF, \textit{Energy Frontier Session}, Fermilab \when{2026} % July 23 2026
  \item DPF, \textit{Outreach \& Engagement Session}, Fermilab \when{2026} % July 21 2026

  \item AI4MuC -- \textit{Overview Talk} \when{2026} % June 29, 2026


  \item IDEC Conference, Columbia College of Chicago  \when{2026} % Mar 2026

  \item Good $\nu$s, University of Pittsburgh -- \emph{Invited}  \when{2025} % 29 Oct 2025

  \item Physics At The Highest Energies With Colliders, Galileo Galilei Institute, Florence, Italy -- \emph{Invited}  \when{2025} % 29 July 2025


  \item LHCP2025, NTU, Taipei, Taiwan -- \emph{Invited}  \when{2025} % 4 March 2025


  \item Aspen Center for Physics -- \emph{Invited Muon Collider Overview Plenary}  \when{2025} % 4 March 2025

  \item Inaugural US Muon Collider Meeting, Fermilab -- \emph{Invited Experimental Overview Plenary}  \when{2024} % 7 August 2024

  \item Workshop on Long-Lived Particles, University of Tokyo, Tokyo, Japan -- \emph{o.b.o. CMS}  \when{2024} % 4 July 2024

  \item DPF/Pheno, Pittsburgh, PA \when{2024} % May 2024

  \item SESAPS2023, EKU, Richmond, KY -- \emph{Invited `Highlight' Talk}\when{2023} % XX Nov 2023

  \item Fermilab Users Meeting -- \emph{Invited Plenary} \when{2023}
  % June 28, 2023

  \item UCSB, Kavli Institute of Theoretical Physics, Muon Collider Workshop -- \emph{Invited Plenary} \when{2023}
  % March 8, 2023

  \item CMS Exotica Workshop -- \emph{Invited Plenary} \when{2022} % 11 Dec 2022

  \item ML4Jets, Rutgers University, Piscataway, NJ \when{2022} % 1 Nov 2022

  \item Snowmass Community Summer Study Workshop, Seattle, WA \when{2022} % 21 July 2022

  \item APS April Meeting, New York City, NY \when{2022} % 10 April 2022

  \item CMS Exotica Workshop -- \emph{Invited Plenary} \when{2021} % 24 Nov 2021

  \item SESAPS2021, FSU, Tallahassee, Fl -- \emph{Invited `Highlight' Talk}\when{2021} % 19 Nov 2021


  % \item \textbf{ColliderScope: Reaching New Audiences with Electronic(s) Music}
  %   \begin{itemizetighter}
  %   \item 18th International Particle Physics Outreach Group meeting, CERN \when{2019} % 28 November 2019
  %   \item International Conference on New Frontiers in Physics, Kolymbari, Crete \when{2019} % 21 August 2019
  %   \item APS Division of Particles and Fields Meeting, Boston, MA \when{2019} % 30 July 2019
  %   \end{itemizetighter}

  \item LHCP2021, Virtual -- \emph{Invited Plenary, o.b.o. ATLAS and CMS}\when{2021}
  % \item 6/9/2021 Neutral FIPs at the LHC and Beyond

  \item APS April Meeting, Virtual -- \emph{Invited Plenary, o.b.o. ATLAS}\when{2021}
  % \item \textbf{New Long-Lived Particle Results from the LHC} -- \emph{Plenary Talk} \when{2021} % 4/17/2021
  %   \begin{itemizetighter}
  %   \item APS April Meeting, Virtual
  %   \end{itemizetighter}

  \item A Rainbow Of Dark Sectors, Aspen Center for Physics, Virtual -- \emph{Invited Plenary}\when{2021}
  % \item \textbf{The Lifetime of the Long-Lived Particle Era at the LHC} -- \emph{Plenary Talk} \when{2021} % 24 Mar 2021
  %   \begin{itemizetighter}
  %   \item A Rainbow Of Dark Sectors, Aspen Center for Physics, Virtual
  %   \end{itemizetighter}


  \item LHC Experiments Committee (LHCC) Meeting, CERN, Geneva -- \emph{Open Session} \when{2020}
    % Nov 18, 2020
    % \item \textbf{Updates from the ATLAS Experiment}
    %   \begin{itemizetighter}
    %   \item LHC Experiments Committee (LHCC) Meeting Open Session, CERN, Geneva \when{2020}
    %   \end{itemizetighter}


  \item ICNFP, Kolymbari, Crete, Greece -- \emph{o.b.o. ATLAS}\when{2019} % 21 August 2019
  % \item \textbf{Searches for Long-Lived Particles with the ATLAS Detector}
  %   \begin{itemizetighter}
  %   \item \emph{On behalf of the ATLAS Collaboration}
  %   \item International Conference on New Frontiers in Physics, Kolymbari, Crete \when{2019} % 21 August 2019
  %   \end{itemizetighter}


  \item Higgs Cross-Section Working Group Meeting, CERN, Geneva -- \emph{Invited Plenary} \when{2018} % 11 December 2018
  % \item \textbf{Long-Lived Particles and the Higgs}
  %   \begin{itemizetighter}
  %   \item Higgs Cross-Section Working Group Meeting, CERN, Geneva \when{2018} % 11 December 2018
  %   \end{itemizetighter}

    \item USATLAS Workshop, Pittsburgh, PA -- \emph{Invited Plenary} \when{2018} % 1 August 2018
    % \item \textbf{SUSY and Exotics Overview} -- ``\emph{Headliner}'' Talk
    %   \begin{itemizetighter}
    %   \item USATLAS Workshop, University of Pittsburgh, Pittsburgh, PA \when{2018} % 1 August 2018
    %   \end{itemizetighter}

    % \item \textbf{The SUSY Paradox \& A Toy Entropic Explanation for the Hierarchy Problem}
    %   \begin{itemizetighter}
    %   \item Oskar Klein Centre, Stockholm, Sweden \when{2018} % 22 May 2018
    %   \end{itemizetighter}

  \item SUSY2017, Mumbai, India -- \emph{o.b.o. ATLAS} \when{2017} % 11 December 2017
  % \item \textbf{Searches for squarks and gluinos in scenarios with R-parity violating sparticle decays, or long- lived sparticles with ATLAS}
  %   \begin{itemizetighter}
  %   \item \emph{On behalf of the ATLAS Collaboration}
  %    \item SUSY2017, Mumbai, India \when{2017} % 11 December 2017
  %   \end{itemizetighter}

  \item IHEP-T2E, Kuala Lumpur, Malaysia -- \emph{Plenary, o.b.o. ATLAS} \when{2017} % 27 February 2017
  % \item \textbf{SUSY Searches in the ATLAS Experiment} -- \emph{Plenary Talk}
  %   \begin{itemizetighter}
  %   \item \emph{On behalf of the ATLAS Collaboration}
  %   \item IHEP-T2E, Kuala Lumpur, Malaysia \when{2017} % 27 February 2017
  %   \end{itemizetighter}

  \item SUSY2016, Melbourne, Victoria, Australia\when{2016} % 3 July 2016
  % \item \textbf{The Recursive Jigsaw Reconstruction Technique}
  %   \begin{itemizetighter}
  %   \item SUSY2016, Melbourne, Victoria, Australia \when{2016} % 3 July 2016
  %   \end{itemizetighter}

  % \item \textbf{New Searches for Strongly-Produced SUSY at the ATLAS Experiment}
  %   \begin{itemizetighter}
  %   \item CoEPP Annual Workshop, Torquay, Victoria, Australia\when{2016} % 18 February 2016
  %   \end{itemizetighter}

  \item SUSY2015, Lake Tahoe, CA\when{2015} % 24 August 2015
  % \item \textbf{Applications of the Recursive Jigsaw Technique}
  %   \begin{itemizetighter}
  %   \item SUSY2015, Lake Tahoe, CA\when{2015}% 24 August 2015
  %   \end{itemizetighter}

  \item Kruger2014, Kruger Gate, South Africa -- \emph{o.b.o. ATLAS} \when{2014} % 6 December 2014
  % \item \textbf{SUSY Searches in the ATLAS Experiment} -- \emph{Plenary Talk}
  %   \begin{itemizetighter}
  %   \item \emph{On behalf of the ATLAS Collaboration}
  %   \item Kruger2014, Kruger Gate, South Africa\when{2014}% 6 December 2014
  %   \end{itemizetighter}


  \item Fermilab, LPC Workshop on Exotic Top Partners, Batavia, IL -- \emph{Invited Plenary}\when{2013} % 26 September 2013

  \item SUSY2012, Peking University, Beijing, China -- \emph{o.b.o. ATLAS} \when{2012} % 17 August 2012
  % \item \textbf{Search for supersymmetry in resonant production and R-parity violating signatures with the ATLAS detector}
  %   \begin{itemizetighter}
  %   \item \emph{On behalf of the ATLAS Collaboration}
  %   \item SUSY2012, Peking University, Beijing, China\when{2012}% 17 August 2012
  %   \end{itemizetighter}

% \end{itemize}
\end{itemizetighter}


  \emph{Conference and Workshop Organization}

\begin{itemizetighter}

  \item \textit{USMCC 2026}, Stanford University, \emph{Co-chair (w/ M. Peskin), Program Committee} \when{2026} % Dec 2026
  \item \textit{USMCC 2025}, University of Chicago, IL, \emph{Detector Design and Performance Session Organizer [of 3]} \when{2025} % Aug 2025
  \item \textit{LHCP 2025}, Taipei, Taiwan, \emph{Poster Committee Member} \when{2025} % May 2025
  \item \textit{CPAD 2024}, Knoxville, TN, \emph{Chair of Local Organizing Committee [of 5]} \when{2024} % Nov 2024
  \item \textit{LHCP 2024}, Northeastern University, Boston, MA, \emph{Poster Committee Member} \when{2024} % May 2024

  \item \textit{The Future of High Energy Physics}, Aspen Center for Physics, \emph{Conference Organizer [of 4]} \when{2024} % Mar 2024
  \item \textit{PITT PACC Muon Collider Physics Benchmark Workshop}, Pittsburgh, PA, \emph{Invited Closing Remarks} \when{2023} % Nov 2023
  \item \textit{ML4Jets}, Rutgers University, \emph{Session Chair} \when{2022} % Oct 2022
  \item \textit{APS April Meeting}, New York, NY, \emph{Session Chair} \when{2022} % April 10 2022
  \item \textit{ATLAS SUSY and Exotics Workshop}, Virtual \when{2020} % September 2020
  \item \textit{ATLAS Reaching New Heights Workshop}, CERN \when{2019} % December 2019
  \item \textit{ATLAS SUSY Workshop}, Lecce, Italy \when{2019} % September 2019
  \item \textit{APS Division of Particles and Fields Meeting}, Northeastern University \when{2019} % 29 July - 2 August 2019
  \item \textit{ATLAS SUSY Workshop}, Stockholm, Sweden \when{2018} % 23-25 May 2018
  \item \textit{ATLAS SUSY and Exotics Workshop}, Bucharest, Romania \when{2017} % 8-12 May 2017
  \item \textit{Recursive Jigsaw Workshop}, Harvard University \when{2015} % 21 September 2015
  \item \textit{Workshop on LHC Searches}, LBNL \when{2014} % 30 January 2014

  \end{itemizetighter}

% \clearpage


\section*{Teaching}

  \begin{itemizetighter}
    \item Physics of Cosmic Rays \emph{Seminar \& Design/Build Course w/ UT Architecture \& Design} (PHYS 494) \when{2025}
    \item Classical Mechanics I \textit{\&} II (PHYS 311/312) \when{2021--24}
    \item UTK HEP/Astro Seminar    \when{2022}
    \item Graduate Research Participation Seminar  \when{2021}
  \end{itemizetighter}


  \internalonly{

    \section*{Current Student \textit{\&} Postdoc Mentorship}

    \begin{itemizetighter}
      \item Olivia Clark [Undergrad] \when{2025--}
      \item Micah Hillman [Undergrad] \when{2025--}
      \item \rightnote{UTK EUReCA 2nd Place Presenter Award}

      \item Allie Dabney [Undergrad] \when{2025--}
      \item Shivam [Ph.D.] \when{2024--}
      \item Caroline Riggall [Ph.D.] \when{2024--}
      \item Cordney Nash [Undergrad] \when{2024--}
      \item \rightnote{UTK `Volunteer of Distinction'}

      \item Colby Thompson [Ph.D.] \when{2023--}
      \item Adam Vendrasco [Ph.D.] \when{2022--}
      \item John Lawless [Ph.D.] \when{2022--}
      \item \rightnote{Stelson Fellowship for Outstanding Beginning Research}

      \item Emery Nibigira [Postdoc] \when{2022--}
      \item \rightnote{Fermilab LPC Distinguished Researcher}

      \item \rightnote{Lindau Nobel Laureate Meeting Selected Participant}

    \end{itemizetighter}


    \section*{Student Thesis Committee Membership}

    \begin{itemizetighter}
      \item Colby Thompson [Ph.D.], \emph{chair} \when{2025--}
      \item Raghav Chari [Undergrad, College Honors] \when{2025}
      \item John Lawless [Ph.D.], \emph{chair} \when{2024--}
      \item Brandi Skipworth [Ph.D.] \when{2024--}
      \item Christal Martin [Ph.D.] \when{2024--}
      \item Ewa Glimos [Ph.D.] \when{2024--}
      \item Taylor Sussmane [Undergrad, Honors], \emph{chair} \when{2024}
      \item John Melhorn [Undergrad, Honors] \when{2024}
      \item Luke Carpenter [Undergrad, Honors] \when{2024}
      \item Rebecca Godri [Ph.D.] \when{2023--}
      \item Daniel Ally [Ph.D.] \when{2022--}
      \item Jesse Harris [Ph.D.] \when{2022--}
      \item Satyajit Roy [Ph.D.] \when{2021--24}
      \item Austin Schmier [Ph.D.] \when{2021--24}
      \item Himal Acharya [Ph.D.] \when{2021}
      \item External Ph.D. Dissertation Review for Yale University \when{2021}
      \item External M.S. Thesis Review for University of Adelaide \when{2021}
    \end{itemizetighter}

  }

  \section*{Past Student \textit{\&} Postdoc Mentorship}


  \textit{n.b. Square brackets denote initial position of mentorship relationship. Current position in italics.}

    \begin{itemizetighter}
      \item Dr. Alexander Tuna [Postdoc, UTK] \hfill \emph{Research Staff, UCSD}
      \item Dr. Ankush Kanuganti [Postdoc, UTK] \hfill \emph{Postdoc, BNL}
      \item Micah Hillman [Undergrad, UTK] \hfill \emph{Physics Postbacc, Princeton}
      \item Charles Bell [Undergrad, UTK] \hfill \emph{Physics PhD Student, UMichigan, NSF GRFP Awardee, Cover of Science,}
        \item \hfill \emph{UTK Talley Award for Outstanding Undergrad Research,}
        \item \hfill \emph{UTK Douglas V. Roseberry Distinguished Upper Class Major Award}
      \item Lacey Dishman [Undergrad, UTK] \hfill \emph{Physics PhD Student, Boston University}
      \item Taylor Sussmane [Undergrad, UTK] \hfill \emph{Physics PhD Student, University of Wisconsin, Madison,}
        \item \hfill \emph{Douglas V. Roseberry Distinguished Upper Class Major Award,}
        \item \hfill \emph{UTK Talley Award for Outstanding Undergrad Leadership}
      \item Everett Hagen [Undergrad, UTK]
      \item Jack Peltier [Undergrad, UTK]\hfill \emph{Goldwater Scholar}
      \item Puck Cravens [Undergrad, UTK]
      \item Tanner Wolfe [Undergrad, UTK]
      \item Alexander Sizemore [Undergrad, UTK]
      \item Matthew Durkee [Undergrad, UTK] \hfill \emph{Defense Contractor}
      \item Philipp Wagenknecht [Ph.D., UTK] \hfill \emph{HEP Lab Technician, Brown University}
      \item Dr. Tamas Vami [Ph.D., Johns Hopkins] \hfill \emph{Postdoc, UCSB}
      \item Prof. Karri Folan DiPetrillo [Ph.D., Harvard] \hfill \emph{Asst. Professor, University of Chicago}
      \item Prof. Jennifer Roloff [Ph.D., Harvard] \hfill \emph{Asst. Professor, Brown University}
      \item Dr. Anthony Badea [Ph.D., Harvard] \hfill \emph{Schmidt AI Fellow, University of Chicago}
      \item Dr. Anne Fortman [Ph.D., Harvard] \hfill \emph{Postdoc, LBNL}
      \item Dr. Aaron White [Postdoc, Harvard] \hfill \emph{}
      \item Xiaohe Jia [Ph.D., Harvard] \hfill \emph{Industry}
      \item Madeline Bernstein [Undergrad, Harvard] \hfill \emph{Physics PhD Student, UC Berkeley}
      \item Dr. Stefanie Morgenstern [Ph.D., Dresden/CERN] \hfill \emph{CERN Fellow}
      \item Dr. Russell Smith [Ph.D., Columbia] \hfill \emph{Squarepoint, Financial Sector}
      \item Prof. Emma Alexander [Undergrad, Yale] \hfill \emph{Asst. Professor, Northwestern University}
      \item Dr. Ariel Ekblaw [Undergrad, Yale] \hfill \emph{Founder of the MIT Space Exploration Initiative, CEO of Aurelia Institute}
    \end{itemizetighter}


\section*{Selected Outreach Activities}

  \begin{itemize}

    \item \textbf{\href{https://cosmicrayn.utk.edu}{Cosmic Rayn: Cosmic Ray Educational Experience}} w/ UTK Architecture \textit{\&} Design \when{2024--26}

    \item \textbf{Quantum Canvases: Physics, the Arts, and the Humanities} \when{2023}

      \begin{itemizetighter}
      \item Lead organizer of a series of events in downtown Knoxville on the intersection of physics, the arts, and the humanities in collaboration with UTK Physics \textit{\&} Astronomy, CERN, and the UT Humanities Center.
      \item Electronic music event (``Harmonic Motion: Physics x Electronic Music'')
      \item Events around a collection of science fiction short stories
      \item A panel discussion on the intersection of physics and science fiction
      \item A humanities-facing physics colloquium event
      \end{itemizetighter}


    \item \textbf{Higgs Tenth Anniversary Celebration} \when{2022}

      \begin{itemizetighter}
      \item Primary organizer for a documentary screening of \emph{Particle Fever} and open panel session for the tenth anniversary of the discovery of the Higgs Boson. Held at Central Cinema in Knoxville, TN.
      \end{itemizetighter}


    \item \textbf{STEAM Trailer Design Project} \when{2021--22}

      \begin{itemizetighter}
      \item Worked with the UTK College of Architecture \textit{\&} Design to design and construct an educational experience for local middle schoolers with a focus on the physics of sound and visualization.
      \end{itemizetighter}

    \item \textbf{Keynote presentation at event for Eureka Science Museum}, Grand Junction, CO \when{2021}


    \item \href{http://www.cern.ch/ColliderScope}{\textbf{ColliderScope -- Particle-Physics-Inspired Electronic Music Project}} \when{2019--}

      \begin{itemizetighter}
      \item Music with sound waves that show particle physics images in Lissajous figures
      \item Live performances:
      \item \hspace{2em} Luck Reunion (at Willie Nelson's ranch), Austin, TX (\emph{w/ Micah Nelson, Daniel Lanois}) \when{2026}
      \item \hspace{2em} Fly By Night, Knoxville, TN (\emph{w/ signal:noise; VJ Set}) \when{2026}
      \item \hspace{2em} Fly By Night, Knoxville, TN (\emph{w/ signal:noise; VJ Set}) \when{2025}
      \item \hspace{2em} Fly By Night, Knoxville, TN (\emph{w/ signal:noise; VJ Set}) \when{2025}
      \item \hspace{2em} Scruffy City Hall, Knoxville, TN (\emph{w/ Harmonic Motion}) \when{2025}
      \item \hspace{2em} Scruffy City Hall, Knoxville, TN (\emph{w/ Harmonic Motion}) \when{2024}
      \item \hspace{2em} Old City Performing Arts Center, Knoxville, TN (\emph{w/ Harmonic Motion}) \when{2023}
      \item \hspace{2em} Blue Moon Tavern, Seattle, WA (\emph{w/ Snowmass summer meeting}) \when{2022}
      \item \hspace{2em} Roskilde Festival, Denmark \when{2022}
      \item \hspace{2em} Cross Club, Prague, YouTube, Live from the CERN Control Centre (\emph{w/ ICHEP2020}) \when{2020}
      \item \hspace{2em} Lund, Sweden \when{2020}
      \item \hspace{2em} Aeronaut Brewing, Somerville, MA \when{2019}
      \item \hspace{2em} ATLAS Experiment, Geneva, Switzerland \when{2019}
      \item \hspace{2em} CERN Open Days, Geneva, Switzerland \when{2019}
      \item \hspace{2em} Pohoda Music Festival, Tren\v{c}\'{i}n, Slovakia \when{2019}
      \item Commissioned Compositions:

      \item \hspace{2em} Intro Videos for the CERN Science Gateway Educational Program \when{2023}

      \item \hspace{2em} Intro Videos for CERN Music Festival Program \when{2019}

      \item Presented at multiple international and national conferences
    %   \item 18th International Particle Physics Outreach Group meeting, CERN \when{2019 % 28 November 2019}
    %   \item International Conference on New Frontiers in Physics, Kolymbari, Crete \when{2019 % 21 August 2019}
    %   \item APS Division of Particles and Fields Meeting, Boston, MA \when{2019 % 30 July 2019}
      \end{itemizetighter}


    \item \textbf{Presentations to Grade School Students}
    \begin{itemizetighter}
    \item L\&N STEM Academy, Knoxville, TN \when{2023}
    \item Tel Aviv University -- Future Scientists CERN Tour \when{2017--19}
    \item Five presentations to NJ and NY schools \when{2011--16}
    \end{itemizetighter}

    \item \textbf{Harvard NSW Activities at CERN} - \emph{Video} (\href{https://youtu.be/VHrJAkNsMhM}{YouTube})
    \begin{itemizetighter}
    \item Fully produced video for recruitment \when{2018}
    \end{itemizetighter}

  \item \href{http://www.cern.ch/inparticular}{\textbf{In Particular}} - \emph{Podcast}\when{2015--16}

    \begin{itemizetighter}
    \item Video, audio, and original music production
    % \item Production of ``Video Shorts''
    % \item Web development, graphic design, and episode production
    \item Tens of thousands of downloads ($\sim11$~TB served)
    \item Featured by CERN, iTunes, The Guardian Science Blogs, Ars Technica, Vice, and others
    \end{itemizetighter}


  \item \textbf{Visiting Scholar, Dharma Drum Mountain, Taiwan}
    \begin{itemizetighter}
    \item Several weeks of conversations with monastery's monks and academics \when{2014}
    \end{itemizetighter}


  \end{itemize}


\section*{Press \textit{\&} Interviews}

      \begin{itemizetighter}

      \item \href{https://physics.utk.edu/tova-holmes-and-larry-lee-selected-as-fermilab-distinguished-researchers/}{UTK Physics -- Tova Holmes And Larry Lee Selected As Fermilab Distinguished Researchers}\when{2026}

      \item \href{https://news.utk.edu/2025/02/12/university-of-tennessee-physicist-named-cottrell-scholar/}{UTK News -- University of Tennessee Physicist Named Cottrell Scholar}\when{2025}

      \item \emph{Science} -- L.~Lee, C.~Bell, \textbf{\href{https://www.science.org/toc/science/383/6690}{Cover of Issue 383 6690}}  \when{2024}


      \begin{itemizetightrightpad}
        \item \emph{Muon collider simulation work and renderings made in Unreal Engine were featured on the cover and in the cover article in the issue from March 29, 2024.}
      \end{itemizetightrightpad}

      \item \href{https://engage.aps.org/dpb/resources/newsletters/24-25-newsletter/f10-article}{APS DPB Annual Newsletter -- The Muon Collider: A Modern Machine for a Modern World}\when{2024}
      \item \href{https://cerncourier.com/a/a-new-generation-a-new-vision/}{CERN Courier -- A New Generation, A New Vision} \when{2024}
      \item \href{https://physics.utk.edu/the-particle-detectorists/}{UTK Physics -- The (Particle) Detectorists} \when{2024}
      \item \href{https://physics.utk.edu/four-ut-physics-students-win-doe-scgsr-support/}{UTK Physics -- NSF GRFP Awardee} \when{2024}
      \item \href{https://physics.utk.edu/the-art-of-muon-collisions/}{UTK Physics -- The Art of Muon Collisions} \when{2024}
      \item Interview for article for MIT Tech Review \when{2023}
      \item Interview for article for Symmetry Magazine on CERN music festival program \when{2023}
      \item Interview with Sparks Magazine on Asian Americans in STEM\when{2023}
      \item \href{https://physics.utk.edu/in-my-own-words-undergraduate-taylor-sussmane/}{UTK Physics -- In My Own Words: Undergraduate Taylor Sussmane} \when{2023}
      \item \href{https://physics.utk.edu/challenging-the-standard-model-lawrence-lee-wins-nsf-career-award/}{UTK Physics -- Challenging The Standard Model: Lawrence Lee Wins NSF CAREER Award}\when{2023}
      \item \href{https://physics.utk.edu/the-art-of-science/}{UTK Physics -- The Art of Science} \when{2023}
      \item \href{https://physics.utk.edu/the-yeti-returns/}{UTK Physics -- The YETI Returns} \when{2023}
      \item Interview with Danish podcast TechTopia \when{2022}
      \item \textit{ColliderScope} featured in \href{https://www.symmetrymagazine.org/article/lhc-music-through-the-colliderscope}{Symmetry Magazine} \when{2019}
      \item Interview for article for Symmetry Magazine on naturalness \when{2019}
      \item CERN Courier -- ATLAS Pushes the Limits on Supersymmetry \when{2019}

      \item Interview for Swedish-language podcast \textit{Professor Magenta} about science and art \when{2016}
      \item Public screening of film Particle Fever w/ panel Q\textit{\&}A\when{2014}


      \end{itemizetighter}


\section*{University, College, and Departmental Committees and Other Service Roles}

  \begin{itemizetighter}

    \item UTK Academy for Global Engagement - Mentor \when{2026--27}

    \item CAS Climate, Culture, and Community Committee \when{2025--}

    \item Faculty Search Committee - Particle Theory \when{2025--26}
    \item Faculty Search Committee - Lecturer \when{2025--26}

    \item ORIED Launch Pad Panel on Prestigious Fellowships \& Awards \when{2025}
    \item University NSF CAREER Workshop Panel \when{2025}
    \item Faculty Search Committee - Nuclear Theory \when{2024--25}

    \item Coordinator of Department Web Services \when{2024--}

    \item 4 UTK Senior Thesis Committees \when{2024--}

    \item University NSF CAREER Workshop Panel \when{2024}

    \item Provost Junior Faculty Fellows (\emph{UTK}) \when{2023--25}

    \item Graduate Recruitment Committee  \when{2023--}

    \item SPS Faculty Advisor  \when{2023--}

    \item Assessment Committee  \when{2023--24}

    \item Graduate Advising Committee \when{2022--}

    \item Planning Committee \when{2022--24}

    \item Faculty Search Committee - Medium Energy Nuclear Experiment \when{2022--23}

    \item Qualifying Exam Committee \when{2022--23, 24--25}

    \externalonly{
      \item 9 UTK PhD Thesis Committees \when{2021--}

    }

    \item Community Connections Committee (\emph{Chair}) \when{2021--23}

    \item Department Head Search Committee  (\emph{UTK College of Arts \& Sciences}) \when{2021--22}
  \end{itemizetighter}


% \section*{Professional Activities}

%   \begin{itemize}
%   \item Member, American Association for the Advancement of Science, 2013--Present
%   \item Member, New York Academy of Sciences, 2012--Present
%   \item Member, American Physical Society, 2008--2009
%   \end{itemize}

% \section*{Advising}


%   \begin{itemizetight}

%   \item Harvard University:
%   \begin{itemizetighter}
%     \item Brendon Bullard (PhD Candidate) \hfill[\emph{w/ M. Morii}]
%     \item Jennifer Roloff (2019 PhD) \hfill[\emph{w/ J. Huth}]
%     \item Karri Folan DiPetrillo (2019 PhD) \hfill[\emph{w/ M. Franklin}]
%     \item Adriana Rotaru (2019 Summer Undergraduate)  \hfill[\emph{w/ J. Huth}]
%     \item Madeline Bernstein (2017 Summer Undergraduate)  \hfill[\emph{w/ M. Franklin}]
%   \end{itemizetighter}

%   % \item Uppsala University, Students:
%   % \begin{itemizetighter}
%   %   \item Andreas Gillgren (2018-2019 Visiting Engineering Physics Student)  \hfill[\emph{w/ J. Huth}]
%   % \end{itemizetighter}

%   \item Columbia University:
%   \begin{itemizetighter}
%     \item Russell Smith (2016 PhD)  \hfill[\emph{w/ E. Hughes}]
%   \end{itemizetighter}

%   \item University of Adelaide, Three PhD and Two Undergraduate Students
%   % \begin{itemizetighter}
%   %     \item Jason Oliver (2019 PhD) \hfill[\emph{w/ P. Jackson}]
%   %     \item Anum Qureshi (2019 PhD) \hfill[\emph{'' ''}]
%   %     \item Damir Duvnjak (2014 Honours, 2019 PhD) \hfill[\emph{'' ''}]
%   %     \item Jens Kr\"oger (2015 Visiting) \hfill[\emph{'' ''}]
%   %     \item Finley Borgas (2014 Honours) \hfill[\emph{'' ''}]
%   % \end{itemizetighter}

%   \item Yale University, Five Undergraduate Honors Students
%   % \begin{itemizetight}
%   %     \item Madeleine Barrow (2012)
%   %     \item Ari Brill (2012)
%   %     \item Emma Alexander (2011)
%   %     \item James Giammona (2011)
%   %     \item Alexandra Turrini (2011)
%   % \end{itemizetight}
%   \end{itemizetight}

% \section*{Teaching}

% % \begin{changemargin}{1.5em}{0.0in}

%   \begin{itemizetight}
%     \item Yale University
%     \begin{itemizetight}
%       \item \emph{Fundamentals of Physics}, lecturing assistant for Prof. P. Tipton \when{2011}
%       \item \emph{Physics for the Life Sciences}, lecturing assistant for Prof. S. Mochrie \when{2010}
%       \item \emph{Modern Physical Measurement}, assistant for Prof. J. Harris \when{2010}
%       \item \emph{General Physics Laboratory}, assistant for Prof. K. Baker \when{2009}
%     \end{itemizetight}
%   \end{itemizetight}
% \end{changemargin}

% \clearpage
% \section*{References}

%   \begin{itemize}
%   \item \textbf{Prof. Melissa Franklin} - {\it Mallinckrodt Professor of Physics} - Harvard University
%   \item \hspace{1.6em} \href{mailto:Franklin@Physics.Harvard.edu}{Franklin@Physics.Harvard.edu}
%   \item \textbf{Prof. Tobias Golling} - {\it Associate Professor} - University of Geneva
%   \item \hspace{1.6em} \href{mailto:Tobias.Golling@unige.ch}{Tobias.Golling@unige.ch}
%   \item \textbf{Dr. Andreas Hoecker} - {\it Senior Research Physicist, ATLAS Spokesperson-Elect} - CERN
%   \item \hspace{1.6em} \href{mailto:Andreas.Hoecker@cern.ch}{Andreas.Hoecker@cern.ch}
%   \item \textbf{Prof. John Huth} - {\it Donner Professor of Science} - Harvard University
%   \item \hspace{1.6em} \href{mailto:Huth@g.Harvard.edu}{Huth@g.Harvard.edu}

%     % \pagebreak

%     % \item {\textit{\hspace{-1.6em} [References continued...] }}

%   \item \textbf{Prof. Paul Jackson} - {\it Associate Professor} - University of Adelaide
%   \item \hspace{1.6em} \href{mailto:P.Jackson@adelaide.edu.au}{P.Jackson@adelaide.edu.au}
%   \item \textbf{Dr. Zachary Marshall} - {\it Senior Scientist} - Lawrence Berkeley National Lab
%   \item \hspace{1.6em} \href{mailto:Zach.Marshall@cern.ch}{Zach.Marshall@cern.ch}

%   % \vspace{1em}
%   % \item \textit{Additional References}
%   % \begin{itemize}
%   % \item \textbf{Prof. Paul Tipton} - {\it Chair of Physics Department} - Yale University
%   % \item \hspace{1.6em} \href{mailto:paul.tipton@yale.edu}{Paul.Tipton@Yale.edu}
%   % \end{itemize}

%   \end{itemize}

\cvfooter

\end{document}
"
%%
%% which is what build.sh and .github/workflows/build-cv.yml do.

\newif\ifCVinternal
\ifdefined\CVinternal \CVinternaltrue \else \CVinternalfalse \fi

% \internalonly{<content>} -- included only in the internal edition.
\newcommand{\internalonly}[1]{\ifCVinternal #1\fi}

% \externalonly{<content>} -- included only in the external edition.
\newcommand{\externalonly}[1]{\ifCVinternal\else #1\fi}

% The internal edition says so in its PDF title, so the two files stay
% distinguishable once they have been detached from their filenames.
\ifCVinternal
  \def\LLedition{\name: Curriculum Vitae (internal edition)}
\else
  \def\LLedition{\name: Curriculum Vitae}
\fi

\hypersetup{
  colorlinks  = false,
  pdfauthor   = {\name},
  pdftitle    = {\LLedition},
  pdfsubject  = {Curriculum Vitae},
  pdfkeywords = {},
  pdfpagemode = UseNone,
}


%% ===========================================================================
%% Page layout
%% ===========================================================================

\geometry{
  body       = {6.5in, 9.0in},
  left       = 1.0in,
  top        = 1.0in,
  headheight = 14pt,   % enough room for the running header; silences fancyhdr
}

\pagestyle{fancy}
\fancyhf{}
\renewcommand{\headrulewidth}{0pt}
\fancyhead[L]{\name}
\fancyhead[R]{\thepage}

\setlength{\parindent}{0em}

\titlespacing{\section}{0pt}{0pt}{0pt}
\sectionfont{\rmfamily\mdseries\Large}
\subsectionfont{\rmfamily\mdseries\itshape\large}


%% ===========================================================================
%% Lists
%% ===========================================================================
%% The CV is built almost entirely from bullet-free lists that differ only in
%% how tightly they are set.  itemize itself is redefined, rather than
%% configured through enumitem, so that a plain \begin{itemize} anywhere picks
%% up the CV style.
%%
%%   itemize               default spacing between entries
%%   itemizetight          no extra space between entries
%%   itemizetighter        entries set as close as the leading allows
%%   itemizetightrightpad  tight, and inset from the right margin, for the
%%                         commentary set underneath an entry
%%
%% \LLlist{<left>}{<right>}{<itemsep>}{<parskip>}{<parsep>}

\newcommand{\LLlist}[5]{%
  \list{}{%
    \setlength{\leftmargin}{#1}%
    \setlength{\rightmargin}{#2}%
    \setlength{\itemsep}{#3}%
    \setlength{\parskip}{#4}%
    \setlength{\parsep}{#5}%
  }%
}

\renewenvironment{itemize}%
  {\LLlist{1.5em}{0em}{0.25em}{0pt}{0.25em}}{\endlist}
\newenvironment{itemizetight}%
  {\LLlist{1.5em}{0em}{0em}{-1pt}{0em}}{\endlist}
\newenvironment{itemizetighter}%
  {\LLlist{1.5em}{0em}{0em}{-0.4em}{0em}}{\endlist}
\newenvironment{itemizetightrightpad}%
  {\LLlist{1.5em}{3em}{0em}{-1pt}{0em}}{\endlist}

% Hanging indent for a run-on entry: \parshape 2 3.2em \linewidth \listindent
% \dimexpr\linewidth-\listindent\relax sets the first line full width and
% indents every line after it to \listindent.
\newlength{\listindent}
\setlength{\listindent}{4.8em}

% \changemargin{<left>}{<right>} ... \endchangemargin -- indents a block of the
% document; used to nest the collaboration sub-lists under their headings.
\def\changemargin#1#2{\list{}{\rightmargin#2\leftmargin#1}\item[]}
\let\endchangemargin=\endlist


%% ===========================================================================
%% The right-hand date column
%% ===========================================================================
%% Nearly every entry ends with a date set against the right margin.  \when
%% supplies both the \hfill and the fixed-width box that lines those dates up
%% into one column: the box is left-aligned, so "2024" and "2022--24" begin at
%% the same horizontal position.
%%
%% A \makebox is only as wide as it is declared, so an over-long date used to
%% overflow it silently and print out in the right margin -- six of the dates in
%% this CV ("2021--2026", "2022--23, 24--25", ...) are wider than the column.
%% \when measures the date first and sets an over-wide one flush right instead,
%% which costs the shared left edge on those few entries but keeps every date
%% inside the text block.

\newlength{\datecolumn}
\setlength{\datecolumn}{3.5em}

%% The glue ahead of a date is \datesep plus infinite stretch.  The stretch is
%% what pushes the date to the right margin; the fixed \datesep is a floor, so
%% an entry that grows long enough to reach its date reports an overfull box
%% instead of quietly printing its text hard against the year.  Every entry
%% currently clears it, so the floor costs nothing today -- it is there to catch
%% the next entry that does not.  \whenflush is the fix when one turns up: it
%% gives back the empty right-hand part of the date box, about 1.6em for a
%% four-digit year.
\newlength{\datesep}
\setlength{\datesep}{0.4em}

%% \LLdateglue is the glue that carries the date to the right margin.  It reads
%% as one \hfill most of the time, but is written out so that an entry too long
%% to hold its date can break: TeX takes the \penalty50, leaves the text on the
%% line it filled, and the \hbox{} -- which survives the break where glue would
%% not -- carries the second \hfill onto the next line, so the date still lands
%% flush right instead of stranded against the left margin.  \unskip drops the
%% source space ahead of the date, so the gap is \datesep however the entry was
%% typed.
\newcommand{\LLdateglue}{%
  \unskip\nobreak\hspace{\datesep}\hfill\penalty50 \hbox{}\nobreak\hfill
}

\newlength{\LLdatewidth}
\newcommand{\when}[1]{%
  \LLdateglue
  \settowidth{\LLdatewidth}{#1}%
  \ifdim\LLdatewidth>\datecolumn #1\else\makebox[\datecolumn][l]{#1}\fi
}

% \whenflush{<date>} -- for an entry whose text runs so close to the right
% margin that the shared left edge of the date column will not fit.  Sets the
% date flush right, giving up the alignment to keep the gap.
\newcommand{\whenflush}[1]{\LLdateglue #1}

% \rightnote{<text>} -- a right-aligned italic annotation (an award, a current
% position) hung underneath an entry.  The trailing pad makes it end where a
% four-digit year in the date column ends, rather than at the margin itself.
\newcommand{\rightnote}[1]{%
  \settowidth{\LLdatewidth}{2024}%
  \hfill\emph{#1}\hspace{\dimexpr\datecolumn-\LLdatewidth\relax}%
}


%% ===========================================================================
%% Live publication metrics from INSPIRE-HEP
%% ===========================================================================
%% Citation counts go stale the moment they are typed, so they are not written
%% into the CV by hand.  tools/fetch_inspire.py queries the INSPIRE-HEP API and
%% writes LL_InspireData.tex, a flat list of declarations:
%%
%%     \inspiresetcites{2690093}{3}         % citations of one paper
%%     \inspiresetstat{citations}{207,000}  % an author-level summary figure
%%
%% which the CV body reads back through \inspirepub and \inspirestat.  The
%% values set just below are fallbacks for a build that has never run the
%% fetcher -- an offline checkout, or Overleaf -- so the document always
%% compiles with plausible numbers.  build.sh and CI refresh the data first.

% Declarations written by tools/fetch_inspire.py.  The keys live in \csname, so
% the LL@ namespace prefix needs no \makeatletter here.
\newcommand{\inspiresetcites}[2]{\expandafter\gdef\csname LL@cites@#1\endcsname{#2}}
\newcommand{\inspiresetstat}[2]{\expandafter\gdef\csname LL@stat@#1\endcsname{#2}}

% Fallbacks: the last hand-checked figures, overridden by LL_InspireData.tex.
\inspiresetstat{papers}{1400}
\inspiresetstat{citations}{202,000}
\inspiresetstat{hindex}{209}

% \inspirestat{<key>} -- an author-level figure: papers, citations or hindex.
% A missing key is reported loudly rather than left silently blank.
\newcommand{\inspirestat}[1]{%
  \ifcsname LL@stat@#1\endcsname\csname LL@stat@#1\endcsname\else\textbf{[?~#1]}\fi
}

% \inspirecites{<record>} -- "[n citations]" for a record whose count is known.
% Nothing is printed for an unknown count, or for one below \citesmin, so a
% brand-new paper is left unannotated rather than advertising that nobody has
% cited it yet.
%
% The annotation costs about a dozen characters on every entry, which is enough
% to push an entry that already filled its line onto a second one.  Raising
% \citesmin, or shortening the label, is the dial if that happens too often.
\newcommand{\citesmin}{1}
\newcommand{\LLprintcites}[1]{%
  \ifnum#1<\citesmin\relax\else
    \nobreakspace{\small[#1\nobreakspace\ifnum#1=1\relax citation\else citations\fi]}%
  \fi
}
\newcommand{\inspirecites}[1]{%
  \ifcsname LL@cites@#1\endcsname\LLprintcites{\csname LL@cites@#1\endcsname}\fi
}

% \inspirepub{<record>}{<title>} -- a publication title in bold, linked to its
% INSPIRE-HEP record and followed by that record's live citation count.
\newcommand{\inspirepub}[2]{%
  \href{https://inspirehep.net/literature/#1}{\textbf{#2}}\inspirecites{#1}%
}

% \pub{<url>}{<title>} -- the same, for an entry with no INSPIRE-HEP record:
% arXiv-only preprints, CDS and CERN notes, bare DOI links.
\newcommand{\pub}[2]{\href{#1}{\textbf{#2}}}

% The generated data, if it has been fetched.  \IfFileExists keeps a fresh
% checkout compiling on the fallbacks above.
\IfFileExists{LL_InspireData.tex}{%% Publication metrics fetched from the INSPIRE-HEP API.
%% GENERATED FILE -- do not edit; run tools/fetch_inspire.py instead.
%%
%% \inspiresetcites{<record>}{<count>}  per-paper citation count
%% \inspiresetstat{<key>}{<value>}      author-level summary figure
%% Both are defined in LL_Preamble.tex.

%% Author-level summary, as shown at the top of the Publications section.
%% Papers and citations are rounded down because the CV says "over".
\inspiresetstat{papers}{1400}    % exact: 1,467
\inspiresetstat{citations}{207,000}  % exact: 207,783
\inspiresetstat{hindex}{211}

%% Per-paper citation counts, keyed by INSPIRE record id.
\inspiresetcites{3159704}{0}
\inspiresetcites{3137624}{2}
\inspiresetcites{3119390}{0}
\inspiresetcites{3103093}{4}
\inspiresetcites{2962374}{9}
\inspiresetcites{2958460}{1}
\inspiresetcites{2874678}{26}
\inspiresetcites{2840007}{27}
\inspiresetcites{2752508}{22}
\inspiresetcites{2690093}{3}
\inspiresetcites{2644916}{13}
\inspiresetcites{2642414}{406}
\inspiresetcites{2628398}{83}
\inspiresetcites{2005570}{23}
\inspiresetcites{1856547}{33}
\inspiresetcites{1831504}{102}
\inspiresetcites{1788448}{119}
\inspiresetcites{1756729}{145}
\inspiresetcites{1710078}{319}
\inspiresetcites{1701002}{245}
\inspiresetcites{1630632}{260}
\inspiresetcites{1609773}{260}
\inspiresetcites{1498566}{50}
\inspiresetcites{1458270}{135}
\inspiresetcites{1345355}{100}
\inspiresetcites{1298030}{152}
\inspiresetcites{1191022}{108}
}{%
  \typeout{^^JLL_CV: LL_InspireData.tex not found; using fallback INSPIRE
figures. Run tools/fetch_inspire.py to refresh them.^^J}%
}


%% ===========================================================================
%% Title block
%% ===========================================================================

% \cvheader -- the name, contact details and ORCID that open every document.
\newcommand{\cvheader}{%
  \begin{minipage}[t]{0.4\textwidth}
    {\huge \name}
  \end{minipage}%
  \begin{minipage}[t]{0.6\textwidth}
    \begin{flushright}
      \href{mailto:\email}{\email}\\
      \phone\\
      \href{https://\homepage}{\homepage}\\
      \href{https://orcid.org/\orcid}{\orcid}
      % A postal address, if one is ever wanted here:
      %   CERN CH-1211\\
      %   B\^atiment 40-1-D11\\
      %   23 Gen\`{e}ve, Switzerland
    \end{flushright}
  \end{minipage}%
}

% \cvfooter -- the "last updated" line that closes every document.
\newcommand{\cvfooter}{%
  \medskip
  \begin{center}
    \begin{small}
      Last updated: \today
    \end{small}
  \end{center}%
}
